\section{渔业、林业与宠物食品：三个互不相关的子行业}
\label{sec:fishery}

本章合并讨论 \num{13} 家公司：水产养殖 \num{4}、海洋捕捞 \num{2}、林业 \num{4}、宠物食品 \num{3}
（此外三级行业“农业综合Ⅱ”下的 \num{2} 家也并入本章画像，合计 \num{15} 家）。
它们同属申万农林牧渔一级行业，但\hl{周期性来源完全不同}：
按 2026 年半年报口径，这 \num{13} 家公司合计盈利 \num{3.81} 亿元、仍有 \num{5} 家亏损；
2025 年年报口径下则合计净亏损 \num{7.39} 亿元（趋势对照）。

\begingroup\scriptsize\setlength{\tabcolsep}{2pt}
\begin{longtable}{@{}llrrrrrrr@{}}
\caption{农业上市公司分行业明细（水产养殖、海洋捕捞、林业Ⅲ、宠物食品；财务为 2026 年半年报，ROE 未年化，公告计数窗口 2023-09---2026-09）}\label{tab:comp-fishery}\\
\toprule
代码 & 公司简称 & 市值（亿元） & 营收（亿元） & 同比（\%） & 净利（亿元） & ROE（\%） & 再融资公告 & 重组公告 \\
\midrule
\endfirsthead
\multicolumn{9}{l}{\small（续）}\\
\toprule
代码 & 公司简称 & 市值（亿元） & 营收（亿元） & 同比（\%） & 净利（亿元） & ROE（\%） & 再融资公告 & 重组公告 \\
\midrule
\endhead
\bottomrule
\endlastfoot
000592 & 平潭发展 & 168.1 & 4.8 & -35.1 & -0.69 & -4.0 & 1 & 3 \\
002891 & 中宠股份 & 89.6 & 32.8 & 34.9 & 1.27 & 4.0 & 7 & 0 \\
301498 & 乖宝宠物 & 75.2 & 35.4 & 10.0 & 1.91 & 4.1 & 7 & 0 \\
000798 & 中水渔业 & 36.5 & 20.8 & 18.9 & 1.02 & 20.9 & 0 & 4 \\
600467 & 好当家 & 34.8 & 8.0 & 13.0 & 0.30 & 0.9 & 0 & 0 \\
300094 & 国联水产 & 32.1 & 14.5 & -12.5 & 0.10 & 10.0 & 1 & 4 \\
600265 & *ST景谷 & 30.4 & 0.6 & -54.6 & -0.00 & -0.0 & 16 & 23 \\
002679 & 福建金森 & 27.3 & 0.5 & -9.5 & -0.14 & -1.8 & 1 & 0 \\
600257 & 大湖股份 & 26.4 & 4.3 & 2.0 & 0.17 & 2.1 & 2 & 1 \\
002069 & 獐子岛 & 26.1 & 7.3 & -5.3 & 0.13 & 29.5 & 8 & 5 \\
300673 & 佩蒂股份 & 23.8 & 7.3 & 1.0 & 0.24 & 1.4 & 18 & 0 \\
600097 & 开创国际 & 22.2 & 11.8 & -5.4 & -0.30 & -1.3 & 0 & 2 \\
000663 & 永安林业 & 18.6 & 1.5 & 14.7 & -0.19 & -2.1 & 2 & 10 \\
\end{longtable}
\endgroup


\subsection{水产养殖：一家公司的亏损抵消了整个子行业}
\label{sec:fishery-aqua}

\num{4} 家水产养殖公司在 2026 年上半年全部实现盈利、合计 \num{0.70} 亿元
（上年同期为 $-5.08$ 亿元）：\textbf{好当家}（\num{0.30} 亿元）、\textbf{大湖股份}（\num{0.17} 亿元）、
\textbf{獐子岛}（\num{0.13} 亿元）与 \textbf{国联水产}（\num{0.10} 亿元）。
其中 \textbf{国联水产}的反转最能说明问题：2025 年它一家亏损 \hl{\num{16.31} 亿元}、
ROE \hl{\num{-160.2}\%}——\emph{亏损额已远超其净资产}，一家公司的亏损抵消了整个子行业；
2026 年上半年虽然转正，但利润仅 \num{0.10} 亿元，
2026 年 6 月末资产负债率仍高达 \num{95.2}\%。
国联水产的问题不在于水产周期，而在于\hl{对虾加工出口业务的
“高库存 + 汇率 + 美国关税”三重冲击}：
其营收结构与第 \ref{sec:other-aqua} 节描述的水产品进口增长、出口承压格局一致
（2025 年水产品贸易逆差扩大至 \num{42.18} 亿美元）。

从周期角度看，水产养殖的正向信号同样明确：2025 年二季度起淡水鱼价格大幅上涨
（草鱼塘口均价同比 +$43\%$）、养殖盈利改善，
且\hl{这一景气已在 2026 年上半年体现为报表利润}——水产养殖 \num{4} 家上市公司
2026 年上半年合计归母净利润 \hl{$+0.70$ 亿元、无一家亏损}（2025 年同期为 $-5.08$ 亿元），
与投苗到上市 \num{12}---\num{18} 个月的时滞一致。
但 \textbf{獐子岛}（半年 ROE \num{29.5}\%，\emph{由极薄的净资产放大}）
与 \textbf{好当家}（\num{0.87}\%）的绝对利润仍很小，
反映出底播养殖类企业的生物资产减值风险尚未完全释放。

\subsection{海洋捕捞：小板块、高 ROE}
\label{sec:fishery-marine}

海洋捕捞是全部农业三级行业中\hl{半年 ROE 中位数最高}的子行业
（两家公司的中位数为 \num{9.79}\%，未年化）：
\textbf{中水渔业}半年 ROE \num{20.9}\%、盈利 \num{1.02} 亿元，
\textbf{开创国际}营收 \num{11.8} 亿元、小幅亏损 \num{0.30} 亿元，
两家合计盈利 \num{0.72} 亿元（上年同期 \num{1.14} 亿元）。
远洋渔业的收益来自配额、船队与\hl{海外渔场资源}，
与国内农产品价格几乎无关，是农业板块中少见的\emph{资源型、非周期}业务。

\subsection{林业：最典型的“壳资源”集群}
\label{sec:fishery-forestry}

\num{4} 家林业公司在 2026 年上半年\hl{全部亏损}（合计 \num{-1.03} 亿元）：
\textbf{平潭发展}（\num{-0.69} 亿元，流通市值却达 \num{168.1} 亿元、为本章最高）、
\textbf{永安林业}（\num{-0.19} 亿元，\num{10} 条重组并购公告）、
\textbf{福建金森}（\num{-0.14} 亿元）与
\textbf{*ST景谷}（\num{-0.002} 亿元，\num{23} 条重组并购公告）。
林业的资源属性（林地资产）与其盈利能力长期脱节，
\hl{较高的市值与持续亏损并存，说明市场定价的是林地资产的重估或转型预期，
而不是木材生产的现金流}。\emph{这使林业成为第二部分中融资/重组需求最特殊的板块}：
其潜在融资需求往往与“保壳 + 资产注入”而非产能扩张相关。

\subsection{宠物食品：农业板块中唯一的高成长赛道}
\label{sec:fishery-pet}

\textbf{乖宝宠物}（半年净利润 \num{1.91} 亿元、半年 ROE \num{4.06}\%、营收同比 \num{+10.0}\%）、
\textbf{中宠股份}（\num{1.27} 亿元、\num{4.00}\%、\num{+34.9}\%）、
\textbf{佩蒂股份}（\num{0.24} 亿元、\num{1.36}\%、\num{+1.0}\%）
三家公司在 2026 年上半年全部盈利，但合计利润 \num{3.42} 亿元、
仅为上年同期（\num{6.60} 亿元）的一半左右，
\hl{半年 ROE 中位数 \num{4.00}\%，明显低于 2025 年年报口径的 \num{13.5}\%}。

宠物食品的逻辑与农业周期几乎无关：其需求来自城镇宠物数量与单只消费支出，
上游原料（鸡肉、谷物）价格下跌反而扩大毛利率。
\hl{“农产品价格下行 $\to$ 原料成本下降 $\to$ 毛利率提升”}
构成一条与其他农业子行业完全相反的传导链，
这也是两家公司在 2023---2025 年农业整体下行中保持高增长的原因。
但 2026 年上半年的数据给出了相反的信号：\hl{收入仍在增长（同比中位数 \num{+10.0}\%）、
利润却接近腰斩}，\emph{成长红利正在被竞争与费用侵蚀}，
这是本章最值得关注的边际变化。
其中 \textbf{乖宝宠物}于 2023 年 8 月上市、\textbf{索宝蛋白}于 2023 年 12 月上市，
是全部 \num{104} 家公司中\hl{上市时间最晚的两家}——农业板块的新增融资窗口，
主要留给了这些与周期脱钩的消费型企业。

\subsection{农业综合：农资流通与节水灌溉}
\label{sec:fishery-misc}

申万三级分类中“农业综合Ⅱ”下另有 \num{2} 家公司，
其业务不属于上述任何一类，但恰好代表了农业产业链的两种“服务型”环节：
\textbf{辉隆股份}（半年营收 \num{85.5} 亿元、净利润 \num{1.13} 亿元、半年 ROE \num{3.08}\%）
主营化肥、农药的内外贸分销与自主品牌复合肥生产，
\hl{其收入规模与盈利水平直接反映农资（化肥、农药）价格的波动}——
在化肥价格下行、农民种植收益受挤压的年份，流通环节的毛利被两头压缩；
\textbf{大禹节水}（半年营收 \num{15.9} 亿元、净利润 \num{-0.01} 亿元、半年 ROE \num{-0.04}\%）
从事农业高效节水与农村污水治理，
\hl{其订单来自政府水利项目，具有“逆农业周期”的财政属性}——
2026 年上半年微亏，说明\emph{其业绩节奏更多取决于财政投资而非农业景气本身}。

\begingroup\scriptsize\setlength{\tabcolsep}{2pt}
\begin{longtable}{@{}llrrrrrrr@{}}
\caption{农业上市公司分行业明细（农业综合Ⅱ；财务为 2025 年年报，公告计数窗口 2023-09---2026-09）}\label{tab:comp-misc}\\
\toprule
代码 & 公司简称 & 市值（亿元） & 营收（亿元） & 净利（亿元） & CAGR（\%） & ROE（\%） & 再融资公告 & 重组公告 \\
\midrule
\endfirsthead
\multicolumn{9}{l}{\small（续）}\\
\toprule
代码 & 公司简称 & 市值（亿元） & 营收（亿元） & 净利（亿元） & CAGR（\%） & ROE（\%） & 再融资公告 & 重组公告 \\
\midrule
\endhead
\bottomrule
\endlastfoot
002556 & 辉隆股份 & 47.4 & 151.3 & 1.94 & -7.9 & 5.2 & 0 & 0 \\
300021 & 大禹节水 & 33.4 & 37.6 & 0.49 & 4.4 & 2.2 & 18 & 9 \\
\end{longtable}
\endgroup


\begin{keybox}[渔业、林业与宠物食品的画像结论]
\normalsize
\textbf{① 一个板块，四种周期}：水产养殖看投苗与出口、海洋捕捞看资源与配额、
林业看资产重估、宠物食品看消费升级——\hl{它们唯一的共性是都被归入“农林牧渔”}。\par\vspace{3pt}
\textbf{② 出口加工型的风险与养殖周期不同步}：国联水产 2025 年一家亏损 \num{16.31} 亿元、
ROE \num{-160.2}\%（已超过其净资产），2026 年上半年虽转正（\num{0.10} 亿元），
但 6 月末负债率仍高达 \num{95.2}\%——\hl{加工出口型企业的修复慢于养殖周期}。\par\vspace{3pt}
\textbf{③ 宠物食品的成长红利出现放缓}：三家全部盈利，但合计利润由上年同期的 \num{6.60} 亿元
降至 \num{3.42} 亿元、半年 ROE 中位数 \num{4.00}\%；
\hl{叠加两家新上市公司（乖宝宠物、索宝蛋白），}
说明资本市场的融资功能正从“周期产能”转向“消费升级”，
\emph{但消费型企业的盈利兑现也开始承压}。\par\vspace{3pt}
\textbf{④ 2026 年上半年整体转正。}本组 \num{15} 家公司 2026 年上半年合计归母净利润
\num{4.92} 亿元（上年同期 \num{2.44} 亿元），其中水产养殖 \num{4} 家合计 \num{0.70} 亿元、
\hl{无一家亏损}；改善主要来自水产养殖（上年同期为 \num{-5.08} 亿元），
部分被宠物食品的回落（\num{6.60} $\to$ \num{3.42} 亿元）抵消。
\end{keybox}

\subsection{上市公司画像（\num{15} 家）}
\label{sec:co-fishery}

本节给出本组 \num{15} 家上市公司的画像，按流通市值降序排列：渔业、林业与宠物食品 \num{13} 家，
以及三级行业“农业综合Ⅱ”下的 \num{2} 家（农资流通与节水灌溉，
周期来源与农业生产资料价格及财政投资相关，见第 \ref{sec:fishery-misc} 节）。
数据口径与第 \ref{sec:comp-overview} 节一致：\hl{财务数据统一采用 2026 年半年报}
（合并报表口径、未年化，资产负债率为 2026 年 6 月末），
同时列出 2025 年半年报作为同口径对照、2023---2025 年年报作为趋势对照；
股价为月度收盘价，最后一个交易日为 2026 年 9 月 24 日；周期分段使用与第一部分相同的
ZigZag 算法，阈值取 25\%（个股波动大于商品价格，故高于生猪现货的 20\%）；
公告计数窗口为 2023 年 9 月 1 日至 2026 年 9 月 27 日。

% 农业上市公司画像：水产养殖、海洋捕捞、林业Ⅲ、宠物食品（共 13 家，按流通市值降序）
% 由 scripts/make_company_profiles.py 自动生成，请勿手工修改

\subsubsection{平潭发展（000592）}
\label{co:000592}

\pfl{公司基本情况} 平潭发展成立于 1993 年，注册地为福建平潭综合实验区的林业与木材加工企业，1996 年 3 月在深圳证券交易所上市；2007 至 2008 年经股权分置改革与重大资产重组，福建山田实业发展有限公司取得控制权并注入林业资产，现由自然人刘平山实际控制。根据《2024 年下半年上市公司行业分类结果》，公司所处行业为木材加工和木、竹、藤、棕、草制品业，主营中纤板生产销售、林木资产的抚育与采伐销售及化肥贸易等业务。2025 年归母净利润 -1.13 亿元，2026 年上半年归母净利润 -0.69 亿元，连续亏损；控股子公司中福海峡（平潭）置业有限公司已进入破产清算程序。 公司于 1996-03-27 在深交所主板首发上市，注册资本 19.32 亿元，总股本 19.32 亿股。 截至 2026 年 9 月 24 日收盘价 8.78 元，流通市值 168.1 亿元，近一年+158.2\%，流通市值在三级行业“林业Ⅲ”的 4 家公司中按由大到小排第\zhnumber{1}位，近一年涨跌幅高于全样本中位数（-10.1\%）168.3 个百分点。 按 2026 年半年报披露的行业构成，收入占比前三位为纤维板 74.9\%；房地产 10.9\%；林业 5.1\%。 股权结构方面，控股股东为福建山田实业发展有限公司。

\pfl{历史股价} 自 2021 年 1 月至 2026 年 9 月，股价由 2.31 元变为 8.78 元，累计 +280.1\%，区间最高 15.40 元（2025 年 12 月）、最低 1.30 元（2024 年 6 月）；当前价相当于区间最高价的 57.0\%，为区间最低价的 6.75 倍。 月度收益的年化波动率 58.9\%，高于全样本中位数 37.4\%，在三级行业“林业Ⅲ”的 4 家公司中波动率由高到低排第\zhnumber{1}位。

\begin{center}
\begin{minipage}{\textwidth}\centering
\begin{tikzpicture}
\begin{axis}[
  width=12.6cm, height=3.9cm, scale only axis,
  xmin=2020.922, xmax=2026.888, ymin=1.21, ymax=16.48,
  xtick={2021,2022,2023,2024,2025,2026},
  yearticks,
  yticklabel style={font=\scriptsize},
  ylabel={元/股}, ylabel style={font=\scriptsize},
  grid=major, grid style={LightGray, line width=0.3pt},
  axis line style={InkGray, line width=0.5pt},
  every axis plot/.append style={line width=0.9pt},
]
\addplot[AgriGreen] table[x=x,y=y]{data/clean/profiles/px_000592.dat};
\addplot[SoftRed, dashed, line width=0.7pt] table[x=x,y=y]{data/clean/profiles/pxzz_000592.dat};
\addplot[only marks, mark=*, mark size=1.1pt, SoftRed] table[x=x,y=y]{data/clean/profiles/pxzz_000592.dat};
\end{axis}
\end{tikzpicture}

\captionof{figure}[平潭发展（000592）股价与周期骨架]{平潭发展（000592）月度收盘价（元/股）与 ZigZag 周期骨架（阈值 25\%，圆点为周期转折点）}
\end{minipage}
\end{center}

\pfl{过去周期分析} 以 25\% 的 ZigZag 阈值划分，样本区间内股价形成 3 个主升段与 2 个主跌段。 最大主升段出现在 2024 年 6 月 至 2025 年 12 月，18 个月+1084.6\%；最大主跌段为 2021 年 5 月 至 2024 年 6 月，37 个月-67.1\%。 最近一次转折出现在 2026 年 9 月（上涨段自 2026 年 7 月起，2 个月+50.6\%），当前处于该段之后的震荡之中。 公司股价较申万农林牧渔指数提前约 23 个月见底，是板块中较早定价周期反转的公司之一。 林业的现金流依赖林木采伐与政策补贴，周期长、变现慢，公司 2026 年 6 月末资产负债率 39.3\%、半年 ROE -4.0\%（未年化），在本轮下行中属于随行业同步调整的一类。 综合区间形态看，该股单段涨跌幅度偏大、以趋势段为主（最大主升 +1084.6\%、最大主跌 -67.1\%），年化波动率 58.9\% 与全样本中位数 37.4\% 相比波动明显大于样本中位数。

\pfl{盈利情况} 2026 年上半年营业收入 4.82 亿元（同比 -35.1\%），归母净利润 -0.69 亿元，由上年同期的盈利转为亏损。 以年报为趋势对照，2023---2025 年营业收入由 12.33 亿元变为 13.20 亿元（两年年均复合增速 +3.5\%），归母净利润由 -3.08 亿元变为 -1.13 亿元。 按半年报口径，2026 年上半年的 ROE 为 -4.0\%（未年化）、销售净利率 -33.7\%、毛利率 7.9\%，ROE 在“林业Ⅲ”4 家公司中按数值由高到低排第\zhnumber{4}位。 按同一口径，三级行业“林业Ⅲ”4 家公司的半年 ROE 中位数为 -1.96\%，该公司低于其中位数 2.04 个百分点。 公司披露的应对举措集中在：2026 年 3 月，公司通过董事会与临时股东会决议，以债权人身份推动控股子公司中福置业破产清算。 综合看，公司仍处亏损区间、由盈转亏，半年 ROE 在三级行业内按由高到低排第\zhnumber{4}位，单位盈利能力的修复依赖行业景气与自身成本端的改善，报表弹性主要体现为价格与销量的方向性变化。

\begin{center}\small\setlength{\tabcolsep}{3.4pt}
\captionof{table}[平潭发展（000592）关键财务指标]{平潭发展（000592）关键财务指标（合并报表口径；两个半年报期均未年化；现金短债比 $=$ 货币资金 $\div$ 短期债务）}
\begin{adjustbox}{max width=\textwidth}
\begin{tabular}{@{}lrrrrr@{}}
\toprule
指标 & 2023 年 & 2024 年 & 2025 年 & 2025 年半年报 & 2026 年半年报 \\
\midrule
营业收入（亿元） & 12.33 & 15.63 & 13.20 & 7.43 & 4.82 \\
归母净利润（亿元） & -3.08 & -1.17 & -1.13 & 0.15 & -0.69 \\
毛利率（\%） & 6.7 & 11.2 & 6.9 & 7.5 & 7.9 \\
ROE（\%） & -13.9 & -5.9 & -6.2 & 0.8 & -4.0 \\
资产负债率（\%） & 48.0 & 46.9 & 45.0 & 41.3 & 39.3 \\
经营活动现金流净额（亿元） & 1.25 & -1.03 & 2.23 & 0.84 & 0.36 \\
货币资金（亿元） & 3.57 & 4.29 & 3.20 & 2.90 & 4.17 \\
有息负债合计（亿元） & 4.17 & 5.22 & 5.95 & 4.96 & 5.30 \\
现金短债比（倍） & 0.92 & 0.83 & 0.54 & 0.59 & 0.79 \\
\bottomrule
\end{tabular}
\end{adjustbox}
\end{center}

\begin{center}
\begin{minipage}{\textwidth}\centering
\begin{tikzpicture}
\begin{groupplot}[
  group style={group size=2 by 1, horizontal sep=0.60cm},
  width=6.85cm, height=4.60cm,
  tick label style={font=\tiny},
  title style={font=\scriptsize\heiti},
  yticklabel style={font=\tiny},
  ylabel style={font=\tiny},
  grid=major, grid style={LightGray, line width=0.3pt},
  axis line style={InkGray, line width=0.5pt},
  enlarge x limits={abs=0.8},
  legend style={font=\scriptsize, at={(0.5,-0.20)}, anchor=north,
                legend columns=2, draw=none, fill=none, column sep=6pt},
]
\nextgroupplot[
  ybar, bar width=4pt,
  ymin=-5.87, ymax=18.46,
  xtick={0,1,2,3,4,5.7},
  xticklabels={2021,2022,2023,2024,2025,2026H1},
  ylabel={亿元},
]
\addplot[fill=AgriGreen!75, draw=AgriGreen, bar shift=-2.6pt] table[x=i, y=rev]{data/clean/profiles/fin_000592.dat};
\addplot[fill=SoftRed!60, draw=SoftRed, bar shift=2.6pt] table[x=i, y=ni]{data/clean/profiles/fin_000592.dat};
\addplot[fill=AgriGreen!30, draw=AgriGreen, pattern=north east lines, bar shift=-2.6pt] table[x=x, y=rev]{data/clean/profiles/finh_000592.dat};
\addplot[fill=SoftRed!20, draw=SoftRed, pattern=north east lines, bar shift=2.6pt] table[x=x, y=ni]{data/clean/profiles/finh_000592.dat};
\legend{营业收入（年/半年）, 归母净利润（年/半年）}
\nextgroupplot[
  ymin=-18.9, ymax=54.2,
  xtick={0,1,2,3,4,5.7},
  xticklabels={2021,2022,2023,2024,2025,2026H1},
  ylabel={\%},
]
\addplot[AgriGold, mark=*, mark size=1.3pt] table[x=i, y=roe]{data/clean/profiles/fin_000592.dat};
\addplot[OilBlue, mark=square*, mark size=1.3pt] table[x=i, y=debt]{data/clean/profiles/fin_000592.dat};
\addplot[only marks, AgriGold, mark=*, mark size=2.0pt] table[x=x, y=roe]{data/clean/profiles/finh_000592.dat};
\addplot[only marks, OilBlue, mark=square*, mark size=2.0pt] table[x=x, y=debt]{data/clean/profiles/finh_000592.dat};
\legend{ROE, 资产负债率}
\end{groupplot}
\end{tikzpicture}
\begin{tikzpicture}
\begin{axis}[
  name=finbot,
  width=13.75cm, height=3.10cm, scale only axis,
  ymin=-0.60, ymax=6.67,
  xtick={0,1,2,3,4,5.7},
  xticklabels={2021,2022,2023,2024,2025,2026H1},
  tick label style={font=\tiny},
  yticklabel style={font=\tiny},
  ylabel={亿元}, ylabel style={font=\tiny},
  ybar, bar width=4pt,
  grid=major, grid style={LightGray, line width=0.3pt},
  axis line style={InkGray, line width=0.5pt},
  enlarge x limits={abs=0.8},
  legend style={font=\scriptsize, at={(0.5,-0.26)}, anchor=north,
                legend columns=2, draw=none, fill=none, column sep=10pt},
]
\addplot[fill=AgriGreen!55, draw=AgriGreen, bar shift=-3.0pt] table[x=i, y=cash]{data/clean/profiles/fin_000592.dat};
\addplot[fill=SoftRed!45, draw=SoftRed, bar shift=3.0pt] table[x=i, y=ibd]{data/clean/profiles/fin_000592.dat};
\addplot[fill=AgriGreen!25, draw=AgriGreen, pattern=north east lines, bar shift=-3.0pt] table[x=x, y=cash]{data/clean/profiles/finh_000592.dat};
\addplot[fill=SoftRed!20, draw=SoftRed, pattern=north east lines, bar shift=3.0pt] table[x=x, y=ibd]{data/clean/profiles/finh_000592.dat};
\legend{货币资金（左轴）, 有息负债（左轴）}
\end{axis}
\begin{axis}[
  at={(finbot.south east)}, anchor=south east,
  width=13.75cm, height=3.10cm, scale only axis,
  axis y line*=right, axis x line*=none, xtick=\empty,
  ymin=-1.87, ymax=3.21,
  tick label style={font=\tiny},
  yticklabel style={font=\tiny}, scaled y ticks=false,
  legend style={font=\scriptsize, at={(0.5,-0.44)}, anchor=north,
                legend columns=1, draw=none, fill=none},
]
\addplot[OilBlue, mark=*, mark size=2.2pt, line width=1.3pt] table[x=i, y=ocf]{data/clean/profiles/fin_000592.dat};
\addplot[only marks, OilBlue, mark=*, mark size=3.2pt] table[x=x, y=ocf]{data/clean/profiles/finh_000592.dat};
\legend{经营活动现金流净额（右轴，亿元，蓝点）}
\end{axis}
\end{tikzpicture}

\captionof{figure}[平潭发展（000592）近年主要财务指标]{平潭发展（000592）近年主要财务指标（左上：营业收入与归母净利润；右上：ROE 与资产负债率；下：货币资金、有息负债为左轴柱状，经营活动现金流净额为其右轴的蓝点折线。
2021---2025 年为年报口径，2026H1 为 2026 年半年报/半年末，与其前一年报之间留出间隔、以斜纹或空心点区分；半年报数据未年化）}
\end{minipage}
\end{center}

\pfl{财务分析} 2026 年 6 月末资产负债率 39.3\%，较 2025 年末下降 5.6 个百分点（样本中位数 52.2\%，所处三级行业内按由低到高排第\zhnumber{3}位）。 2026 年 6 月末货币资金 4.17 亿元、短期债务（短期借款与一年内到期非流动负债）5.27 亿元、长期债务 0.04 亿元，有息负债合计 5.30 亿元，现金短债比 0.79 倍——覆盖偏紧。 净有息负债 1.14 亿元，相当于其 2026 年上半年营业收入的 24\%。 流动比率 2.04、速动比率 1.18，短期流动性无明显缺口。 2026 年上半年经营活动现金流净额 0.36 亿元，净利润为负而经营现金流为正，差额主要来自折旧摊销与营运资本变动，说明亏损尚未完全转化为现金流出。 综合看，现金短债比 0.79 倍，短期资金安排偏紧；资产负债率低于样本中位数 12.9 个百分点；有息负债中短期债务占比更高，期限结构仍以滚动续作为主。 公司披露的重整与债务违约风险为：2023 年 12 月，公司已于 2023 年 12 月、2024 年 1 月召开董事会和股东会审议通过作为债权人向法院申请对中福建材城进行破产预重整事项。

\pfl{重大战略转型} 近三年公司的战略动作集中在：其他动作（1 项）：2025 年 11 月，公司出资 1.50 亿元认购宜昌东数一号投资有限责任公司增资。 公告口径上，近三年（2023 年 9 月---2026 年 9 月）公司披露重大重组与并购类公告 3 条、对外投资与转型类公告 0 条、回购与股权激励类公告 17 条，合计公告 328 条。 第一大行业“纤维板”的收入占比由 2021 年年报的 46.3\% 变为 74.9\%（28.6 个百分点），收入结构出现再平衡。 综合判断，公司当前的战略重心不是扩张而是化解债务与完成重整，业务层面的调整让位于司法重整程序。

\begin{center}\small\setlength{\tabcolsep}{4pt}
\captionof{table}[平潭发展（000592）历史融资与资本运作]{平潭发展（000592）历史融资与资本运作（据公司公告与公开报道整理；金额为披露值，含首次公开发行、再融资、债券、银行借款与重整投资等）}
\begin{adjustbox}{max width=\textwidth}
\begin{tabular}{@{}p{1.45cm}p{2.55cm}p{4.05cm}p{4.95cm}@{}}
\toprule
时间 & 方式 & 规模 & 用途与说明 \\
\midrule
1996-03 & 首次公开发行A股股票并在深圳证券交易所上市（IPO） & 发行价 4.33 元/\allowbreak{}股；首发前总股本 11，\allowbreak{}540.00 万股，\allowbreak{}实际发行 3 & 招股公告日 1996 年 3 月 7 日 \\
2010 & 向特定对象非公开发行A股股票（定增） & 发行 7,\allowbreak{}280.00 万股，\allowbreak{}发行后公司股本变更为 6.52 亿元 & 2011 年度股东大会决议以资本公积转增股本 1.96 亿元 \\
2015-12 & 向特定对象非公开发行A股股票（定增） & 发行 1.18 亿股，\allowbreak{}发行价 16.88 元/\allowbreak{}股；募集资金总额 20.00 亿元 & 新增股份于 2015 年 12 月 18 日在深圳证券交易所上市，\allowbreak{}以资本公积金向全体股东每 10 股转增 10 股 \\
2016-03 & 控股股东以所持公司股票为标的非公开发行可交换公司债券（股东层面融资） & 控股股东福建山田实业发展有限公司质押公司 5,\allowbreak{}800 万股作为担保 & 该事项已获深圳证券交易所深证函【2016】106 号确认 \\
\bottomrule
\end{tabular}
\end{adjustbox}
\end{center}

\pfl{历史融投资情况} 从现金流量表看，2025 年公司取得借款收到的现金 0.26 亿元、偿还债务支付的现金 0.04 亿元、吸收权益性投资 0.00 亿元，筹资活动现金流净额 0.05 亿元。 2026 年上半年筹资活动现金流净额为 -0.01 亿元（取得借款 — 亿元、偿还债务 — 亿元，未年化），已转为净流出，处于净还债状态。 2025 年分配股利、利润或偿付利息支付的现金合计 0.00 亿元。 公开披露的融资记录共 4 笔，工具类型以 IPO、定向增发为主；金额与用途见下表。 其中规模最大的两笔为：1996 年 3 月，IPO，发行价 4.33 元/股，首发前总股本 11，540.00 万股，实际发行 3。 股本方面，近三年披露股本变动 4 次（股权分置受限股份上市 1 次）；总股本自 2021 年末以来未发生变化（19.32 亿股）。 分红方面，近三年未实施现金分红，与重整期间的资金约束一致。 近三年“再融资”类公告 1 条，涉及发行股份购买资产。 综合看，公司融资结构上以债务滚动为主、股权融资为补充；2025 年筹资活动现金流净额 0.05 亿元，年度内仍为净融入，与前述经营与投资现金流的方向基本一致。

\pfl{未来潜在融资需求分析} 2026 年 6 月末货币资金 4.17 亿元，而短期债务 5.27 亿元（现金短债比 0.79 倍），2026 年上半年经营现金流净额 0.36 亿元。按“短期债务减去账面货币资金”的滚动偿付口径，静态缺口约 1.10 亿元；若再计入上半年亏损的年化额（1.39 亿元），维持一年正常经营所需的外部资金约 1.10---2.49 亿元。 公司 2026 年上半年归母净利润为负（-0.69 亿元），但 6 月末资产负债率 39.3\% 尚未进入第一档（亏损且负债率 $>$65\%）的区间，当前融资需求以补充流动资金、支撑低谷期经营性支出为主。 公开披露的后续资金安排线索包括：股权融资线索：2026 年 9 月，公司未公告定向增发预案、配股、可转债或公司债等再融资计划。 综合看，公司的外部资金需求以补充营运资金为主、项目投入为辅，滚动偿付口径下的静态缺口约 1.10 亿元，融资的时点与规模更多取决于债权人与重整投资人的进度。




\subsubsection{中宠股份（002891）}
\label{co:002891}

\pfl{公司基本情况} 山东烟台的宠物食品企业，全称烟台中宠食品股份有限公司，2014 年 11 月 24 日由烟台中宠食品有限公司整体变更设立，2017 年 8 月 21 日在深圳证券交易所上市；主营宠物食品和用品的研发、生产与销售，采取自主品牌与 OEM/ODM 双轨并行模式，自有品牌包括“Wanpy 顽皮”“TOPTREES 领先”“新西兰 ZEAL 真致”等。2025 年归母净利润 3.65 亿元、同比下降 7.28\%；2026 年上半年营收 32.806 亿元、同比增长 34.88\%，但归母净利润 1.27 亿元、同比下降 37.63\%。公司近年连续通过可转债与非公开发行融资，2026 年 3 月完成可转债提前赎回。 公司于 2017-08-21 在深交所主板首发上市，注册资本 3.22 亿元，总股本 2.94 亿股。 截至 2026 年 9 月 24 日收盘价 27.83 元，流通市值 89.6 亿元，近一年-47.2\%，流通市值在三级行业“宠物食品”的 3 家公司中按由大到小排第\zhnumber{1}位，近一年涨跌幅低于全样本中位数（-10.1\%）37.2 个百分点。 按 2026 年半年报披露的行业构成，收入占比前两位为宠物食品及用品 93.9\%；其他(补充) 6.1\%。

\pfl{历史股价} 自 2021 年 1 月至 2026 年 9 月，股价由 48.72 元变为 27.83 元，累计 -42.9\%，区间最高 62.86 元（2025 年 5 月）、最低 19.16 元（2024 年 7 月）；当前价相当于区间最高价的 44.3\%，为区间最低价的 1.45 倍。 月度收益的年化波动率 40.7\%，高于全样本中位数 37.4\%，在三级行业“宠物食品”的 3 家公司中波动率由高到低排第\zhnumber{2}位。

\begin{center}
\begin{minipage}{\textwidth}\centering
\begin{tikzpicture}
\begin{axis}[
  width=12.6cm, height=3.9cm, scale only axis,
  xmin=2020.922, xmax=2026.888, ymin=17.82, ymax=67.26,
  xtick={2021,2022,2023,2024,2025,2026},
  yearticks,
  yticklabel style={font=\scriptsize},
  ylabel={元/股}, ylabel style={font=\scriptsize},
  grid=major, grid style={LightGray, line width=0.3pt},
  axis line style={InkGray, line width=0.5pt},
  every axis plot/.append style={line width=0.9pt},
]
\addplot[AgriGreen] table[x=x,y=y]{data/clean/profiles/px_002891.dat};
\addplot[SoftRed, dashed, line width=0.7pt] table[x=x,y=y]{data/clean/profiles/pxzz_002891.dat};
\addplot[only marks, mark=*, mark size=1.1pt, SoftRed] table[x=x,y=y]{data/clean/profiles/pxzz_002891.dat};
\end{axis}
\end{tikzpicture}

\captionof{figure}[中宠股份（002891）股价与周期骨架]{中宠股份（002891）月度收盘价（元/股）与 ZigZag 周期骨架（阈值 25\%，圆点为周期转折点）}
\end{minipage}
\end{center}

\pfl{过去周期分析} 以 25\% 的 ZigZag 阈值划分，样本区间内股价形成 4 个主升段与 4 个主跌段。 最大主升段出现在 2024 年 7 月 至 2025 年 5 月，10 个月+228.1\%；最大主跌段为 2021 年 4 月 至 2022 年 4 月，12 个月-59.6\%。 最近一次转折出现在 2026 年 6 月（下跌段自 2025 年 5 月起，13 个月-57.5\%），当前处于该段的延续之中。 公司股价较申万农林牧渔指数提前约 22 个月见底，是板块中较早定价周期反转的公司之一。 宠物食品属于消费型农业（第 \ref{sec:financing-strategy} 节归入“向消费端延伸”），收入增长由宠物数量与单宠消费驱动，与农产品价格周期的联动最弱，公司 2026 年上半年实现盈利、半年 ROE 4.0\%（未年化），好于全样本中位数（0.75\%），下行周期对其报表的冲击小于同业。 综合区间形态看，该股单段涨跌幅度偏大、以趋势段为主（最大主升 +228.1\%、最大主跌 -59.6\%），年化波动率 40.7\% 与全样本中位数 37.4\% 相比波动与样本中位数接近。

\pfl{盈利情况} 2026 年上半年营业收入 32.81 亿元（同比 +34.9\%），归母净利润 1.27 亿元，同比 -37.6\%。 以年报为趋势对照，2023---2025 年营业收入由 37.47 亿元变为 52.21 亿元（两年年均复合增速 +18.0\%），归母净利润由 2.33 亿元变为 3.65 亿元。 按半年报口径，2026 年上半年的 ROE 为 4.0\%（未年化）、销售净利率 4.2\%、毛利率 25.9\%，ROE 在“宠物食品”3 家公司中按数值由高到低排第\zhnumber{2}位。 同期全样本半年 ROE 中位数 0.75\%，公司与之相差 3.25 个百分点（略好于中位数公司）。 公司披露的应对举措集中在：2026 年 3 月，公司实施“中宠转 2”可转债提前赎回。 综合看，公司在报告期维持盈利（高于全样本半年 ROE 中位数 0.75\%），利润的可持续性取决于宠物食品景气的延续与成本管控的执行。

\begin{center}\small\setlength{\tabcolsep}{3.4pt}
\captionof{table}[中宠股份（002891）关键财务指标]{中宠股份（002891）关键财务指标（合并报表口径；两个半年报期均未年化；现金短债比 $=$ 货币资金 $\div$ 短期债务）}
\begin{adjustbox}{max width=\textwidth}
\begin{tabular}{@{}lrrrrr@{}}
\toprule
指标 & 2023 年 & 2024 年 & 2025 年 & 2025 年半年报 & 2026 年半年报 \\
\midrule
营业收入（亿元） & 37.47 & 44.65 & 52.21 & 24.32 & 32.81 \\
归母净利润（亿元） & 2.33 & 3.94 & 3.65 & 2.03 & 1.27 \\
毛利率（\%） & 26.3 & 28.2 & 29.5 & 31.4 & 25.9 \\
ROE（\%） & 10.9 & 16.8 & 13.5 & 7.8 & 4.0 \\
资产负债率（\%） & 44.9 & 41.2 & 42.7 & 40.9 & 44.0 \\
经营活动现金流净额（亿元） & 4.47 & 4.96 & 5.34 & 2.35 & -1.93 \\
货币资金（亿元） & 5.27 & 5.13 & 12.73 & 13.87 & 13.68 \\
有息负债合计（亿元） & 14.34 & 11.68 & 15.19 & 14.94 & 18.56 \\
现金短债比（倍） & 0.77 & 1.22 & 1.27 & 1.38 & 0.85 \\
\bottomrule
\end{tabular}
\end{adjustbox}
\end{center}

\begin{center}
\begin{minipage}{\textwidth}\centering
\begin{tikzpicture}
\begin{groupplot}[
  group style={group size=2 by 1, horizontal sep=0.60cm},
  width=6.85cm, height=4.60cm,
  tick label style={font=\tiny},
  title style={font=\scriptsize\heiti},
  yticklabel style={font=\tiny},
  ylabel style={font=\tiny},
  grid=major, grid style={LightGray, line width=0.3pt},
  axis line style={InkGray, line width=0.5pt},
  enlarge x limits={abs=0.8},
  legend style={font=\scriptsize, at={(0.5,-0.20)}, anchor=north,
                legend columns=2, draw=none, fill=none, column sep=6pt},
]
\nextgroupplot[
  ybar, bar width=4pt,
  ymin=-5.22, ymax=58.48,
  xtick={0,1,2,3,4,5.7},
  xticklabels={2021,2022,2023,2024,2025,2026H1},
  ylabel={亿元},
]
\addplot[fill=AgriGreen!75, draw=AgriGreen, bar shift=-2.6pt] table[x=i, y=rev]{data/clean/profiles/fin_002891.dat};
\addplot[fill=SoftRed!60, draw=SoftRed, bar shift=2.6pt] table[x=i, y=ni]{data/clean/profiles/fin_002891.dat};
\addplot[fill=AgriGreen!30, draw=AgriGreen, pattern=north east lines, bar shift=-2.6pt] table[x=x, y=rev]{data/clean/profiles/finh_002891.dat};
\addplot[fill=SoftRed!20, draw=SoftRed, pattern=north east lines, bar shift=2.6pt] table[x=x, y=ni]{data/clean/profiles/finh_002891.dat};
\legend{营业收入（年/半年）, 归母净利润（年/半年）}
\nextgroupplot[
  ymin=-3.6, ymax=49.4,
  xtick={0,1,2,3,4,5.7},
  xticklabels={2021,2022,2023,2024,2025,2026H1},
  ylabel={\%},
]
\addplot[AgriGold, mark=*, mark size=1.3pt] table[x=i, y=roe]{data/clean/profiles/fin_002891.dat};
\addplot[OilBlue, mark=square*, mark size=1.3pt] table[x=i, y=debt]{data/clean/profiles/fin_002891.dat};
\addplot[only marks, AgriGold, mark=*, mark size=2.0pt] table[x=x, y=roe]{data/clean/profiles/finh_002891.dat};
\addplot[only marks, OilBlue, mark=square*, mark size=2.0pt] table[x=x, y=debt]{data/clean/profiles/finh_002891.dat};
\legend{ROE, 资产负债率}
\end{groupplot}
\end{tikzpicture}
\begin{tikzpicture}
\begin{axis}[
  name=finbot,
  width=13.75cm, height=3.10cm, scale only axis,
  ymin=-1.86, ymax=20.79,
  xtick={0,1,2,3,4,5.7},
  xticklabels={2021,2022,2023,2024,2025,2026H1},
  tick label style={font=\tiny},
  yticklabel style={font=\tiny},
  ylabel={亿元}, ylabel style={font=\tiny},
  ybar, bar width=4pt,
  grid=major, grid style={LightGray, line width=0.3pt},
  axis line style={InkGray, line width=0.5pt},
  enlarge x limits={abs=0.8},
  legend style={font=\scriptsize, at={(0.5,-0.26)}, anchor=north,
                legend columns=2, draw=none, fill=none, column sep=10pt},
]
\addplot[fill=AgriGreen!55, draw=AgriGreen, bar shift=-3.0pt] table[x=i, y=cash]{data/clean/profiles/fin_002891.dat};
\addplot[fill=SoftRed!45, draw=SoftRed, bar shift=3.0pt] table[x=i, y=ibd]{data/clean/profiles/fin_002891.dat};
\addplot[fill=AgriGreen!25, draw=AgriGreen, pattern=north east lines, bar shift=-3.0pt] table[x=x, y=cash]{data/clean/profiles/finh_002891.dat};
\addplot[fill=SoftRed!20, draw=SoftRed, pattern=north east lines, bar shift=3.0pt] table[x=x, y=ibd]{data/clean/profiles/finh_002891.dat};
\legend{货币资金（左轴）, 有息负债（左轴）}
\end{axis}
\begin{axis}[
  at={(finbot.south east)}, anchor=south east,
  width=13.75cm, height=3.10cm, scale only axis,
  axis y line*=right, axis x line*=none, xtick=\empty,
  ymin=-3.82, ymax=7.53,
  tick label style={font=\tiny},
  yticklabel style={font=\tiny}, scaled y ticks=false,
  legend style={font=\scriptsize, at={(0.5,-0.44)}, anchor=north,
                legend columns=1, draw=none, fill=none},
]
\addplot[OilBlue, mark=*, mark size=2.2pt, line width=1.3pt] table[x=i, y=ocf]{data/clean/profiles/fin_002891.dat};
\addplot[only marks, OilBlue, mark=*, mark size=3.2pt] table[x=x, y=ocf]{data/clean/profiles/finh_002891.dat};
\legend{经营活动现金流净额（右轴，亿元，蓝点）}
\end{axis}
\end{tikzpicture}

\captionof{figure}[中宠股份（002891）近年主要财务指标]{中宠股份（002891）近年主要财务指标（左上：营业收入与归母净利润；右上：ROE 与资产负债率；下：货币资金、有息负债为左轴柱状，经营活动现金流净额为其右轴的蓝点折线。
2021---2025 年为年报口径，2026H1 为 2026 年半年报/半年末，与其前一年报之间留出间隔、以斜纹或空心点区分；半年报数据未年化）}
\end{minipage}
\end{center}

\pfl{财务分析} 2026 年 6 月末资产负债率 44.0\%，较 2025 年末上升 1.3 个百分点（样本中位数 52.2\%，所处三级行业内按由低到高排第\zhnumber{3}位）。 2026 年 6 月末货币资金 13.68 亿元、短期债务（短期借款与一年内到期非流动负债）16.07 亿元、长期债务 2.50 亿元，有息负债合计 18.56 亿元，现金短债比 0.85 倍——覆盖偏紧。 净有息负债 4.88 亿元，相当于其 2026 年上半年营业收入的 15\%。 流动比率 1.29、速动比率 0.97，短期偿债能力处于正常区间。 2026 年上半年经营活动现金流净额 -1.93 亿元，净利润为正但经营现金流为负，利润的现金含量偏低。 2026 年 6 月末在建工程 4.39 亿元，较上年末增加 1.75 亿元，仍有在建投入。 2026 年 6 月末商誉 2.26 亿元，占归母净资产的 6.3\%，是净资产中最脆弱的一块。 综合看，现金短债比 0.85 倍，短期资金安排偏紧；资产负债率低于样本中位数 8.2 个百分点；有息负债中短期债务占比更高，期限结构仍以滚动续作为主。 公司披露的股东与质押风险为：2026 年 6 月，控股股东烟台中幸生物科技有限公司所持公司股份中存在质押 1，090.00 万股的情形。

\pfl{重大战略转型} 近三年公司的战略动作集中在：融资与资本运作（2 项）：2019 年，可转债募集资金投向年产 3 万吨宠物湿粮项目，提升湿粮品类产能；并购与重组（1 项）：2026 年 4 月，公司通过“自主建厂+并购整合”建设横跨五大洲的生产基地网络。 公告口径上，近三年（2023 年 9 月---2026 年 9 月）公司披露重大重组与并购类公告 0 条、对外投资与转型类公告 1 条、回购与股权激励类公告 34 条，合计公告 422 条。 第一大行业仍是“宠物食品及用品”，收入占比 93.9\%，与 2021 年年报相比没有方向性变化。 综合判断，报告期内公司有 1 条与战略相关的公告落地（重大重组与并购 0 条、对外投资与转型 1 条），资本运作与主业调整之间存在直接关联。

\begin{center}\small\setlength{\tabcolsep}{4pt}
\captionof{table}[中宠股份（002891）历史融资与资本运作]{中宠股份（002891）历史融资与资本运作（据公司公告与公开报道整理；金额为披露值，含首次公开发行、再融资、债券、银行借款与重整投资等）}
\begin{adjustbox}{max width=\textwidth}
\begin{tabular}{@{}p{1.45cm}p{2.55cm}p{4.05cm}p{4.95cm}@{}}
\toprule
时间 & 方式 & 规模 & 用途与说明 \\
\midrule
2017-08 & 首次公开发行A股股票并在深圳证券交易所上市（IPO） & 发行价 15.46 元/\allowbreak{}股；首发前总股本 7，\allowbreak{}500.00 万股，\allowbreak{}实际发行 2 & 限责任公司，\allowbreak{}发行方式为网下询价配售与网上定价发行 \\
2019-03 & 公开发行可转换公司债券（中宠转债，债券代码128054） & 发行 1.94 亿元，\allowbreak{}扣除不含税发行费用后实际募集资金净额 1.84 亿元 & 用于年产 3 万吨宠物湿粮项目，\allowbreak{}该项目投资总额 2.70 亿元。；债券于 2019 年 3 月 14 日在深交所上市；公司后续实施提前赎回 \\
2019-07-31 & 重大资产重组（终止） & --- & 不适用。；公司终止重大资产重组并自 2019 年 7 月 31 日开市起复牌 \\
2020-10 & 非公开发行A股股票（2020年非公开发行） & 发行 1，\allowbreak{}737.33 万股，\allowbreak{}发行价格 37.50 元/\allowbreak{}股；募集资金总额 6.51 亿元 & 拟投入年产 6 万吨宠物干粮项目、\allowbreak{}年产 2 万吨宠物湿粮新西兰项目 2.30 亿元、\allowbreak{}营销中心建设及营销渠道智能化升级项目 1.14 亿元、\allowbreak{}补充流动资金 7 \\
2022-11 & 公开发行可转换公司债券（中宠转2，债券代码127076） & 发行 7.69 亿元 & 投向年产 6 万吨高品质宠物干粮项目 2.30 亿元、\allowbreak{}年产 4 万吨新型宠物湿粮项目 2.56 亿元、\allowbreak{}年产 2 \\
2026-01-23 & 股份回购方案 & 拟以集中竞价交易方式回购公司 A 股股票，\allowbreak{}回购价格不超过 78.00 元/\allowbreak{}股（含） & 使用公司自有资金和自筹资金回购公司股份。；公司于 2026 年 5 月 19 日公告取得金融机构股票回购专项贷款承诺函 \\
2026-05-19 & 股票回购专项贷款 & 取得金融机构股票回购专项贷款承诺函 & 专项用于回购公司股份。；回购方案已于 2026 年 1 月 23 日经第四届董事会第二十次会议审议通过 \\
\bottomrule
\end{tabular}
\end{adjustbox}
\end{center}

\pfl{历史融投资情况} 从现金流量表看，2025 年公司取得借款收到的现金 16.85 亿元、偿还债务支付的现金 10.75 亿元、吸收权益性投资 0.04 亿元，筹资活动现金流净额 3.43 亿元（2023 年为取得借款 6.82 亿元、偿还债务 3.66 亿元）。 2026 年上半年筹资活动现金流净额为 6.89 亿元（取得借款 14.30 亿元、偿还债务 7.43 亿元，未年化），仍处于净融入状态。 2025 年分配股利、利润或偿付利息支付的现金合计 1.37 亿元。 公开披露的融资记录共 7 笔，工具类型以 IPO、可转债、定向增发、银行借款为主；金额与用途见下表。 其中规模最大的两笔为：2026 年 1 月，股份回购方案，拟以集中竞价交易方式回购公司 A 股股票，回购价格不超过 78.00 元/股（含）；2022 年 11 月，可转债，发行 7.69 亿元。 股本方面，近三年披露股本变动 16 次（可转债转股 13 次、其他 1 次）；总股本自 2021 年末以来未发生变化（2.94 亿股）。 分红方面，近三年现金分红 7 次、合计每股派现 1.12 元，最近一次实施在 2026 年 6 月。 近三年“再融资”类公告 7 条，涉及可转债。 综合看，公司融资结构上以债务滚动为主、股权融资为补充；2025 年筹资活动现金流净额 3.43 亿元，年度内仍为净融入，与前述经营与投资现金流的方向基本一致。

\pfl{未来潜在融资需求分析} 2026 年 6 月末货币资金 13.68 亿元，而短期债务 16.07 亿元（现金短债比 0.85 倍），2026 年上半年经营现金流净额 -1.93 亿元。按“短期债务减去账面货币资金”的滚动偿付口径，静态缺口约 2.39 亿元；上半年盈利或接近盈亏平衡，该缺口不随经营亏损进一步扩大。 公司 2026 年 6 月末资产负债率 44.0\%、2026 年上半年 ROE 4.0\%（未年化），现金短债比 0.85 倍、短债覆盖不足，融资需求首先用于置换与滚动存量债务、补充营运资金，属于“补血型”而非扩张型。 公开披露的后续资金安排线索包括：股权融资线索：2026 年 4 月，公司 2022 年可转债募投项目中“年产 6 万吨高品质宠物干粮项目”“年产 4 万吨新型宠物湿粮项目”截至 2025 年末投资进度分别为 0.77\% 与 1.06\%；2026 年 8 月，公司未公告新的定向增发预案、配股或债券发行计划；债务与授信安排：2026 年 5 月，公司取得金融机构股票回购专项贷款承诺函，为最新一笔公开披露的银行融资安排。 资金需求还要叠加在建工程：期末在建工程 4.39 亿元、2026 年上半年购建固定资产等支出 2.81 亿元（未年化），这部分支出在周期底部通常会被主动放缓。 综合看，公司的外部资金需求以补充流动性与项目投入为主，滚动偿付口径下的静态缺口约 2.39 亿元，能否落地取决于授信续作与发行窗口的配合。




\subsubsection{乖宝宠物（301498）}
\label{co:301498}

\pfl{公司基本情况} 山东聊城的宠物食品企业，全称乖宝宠物食品集团股份有限公司，2006 年 6 月 26 日注册成立，2023 年 8 月 16 日在深圳证券交易所创业板上市；主营犬用和猫用多品类宠物食品的研发、生产和销售，产品包括主粮系列、零食系列与保健品系列，旗下拥有“麦富迪”“弗列加特”等自有品牌，并在泰国、新西兰布局高端产能。控股股东、实际控制人为秦华。2025 年归母净利润 6.73 亿元、同比增长 7.75\%；2026 年上半年营业收入 35.430 亿元、同比增长 10.01\%，但归母净利润 1.91 亿元、同比下降 49.57\%。2026 年 8 月公司推出向不特定对象发行可转债预案，拟募集资金不超过 11 亿元。 公司于 2023-08-16 在深交所创业板首发上市，注册资本 4.00 亿元，总股本 4.00 亿股。 截至 2026 年 9 月 24 日收盘价 39.67 元，流通市值 75.2 亿元，近一年-57.1\%，流通市值在三级行业“宠物食品”的 3 家公司中按由大到小排第\zhnumber{2}位，近一年涨跌幅低于全样本中位数（-10.1\%）47.0 个百分点。 按 2026 年半年报披露的行业构成，收入占比前两位为宠物食品及用品 99.6\%；其他(补充) 0.4\%。

\pfl{历史股价} 以 2023 年 8 月为起点、2026 年 9 月为终点，股价由 45.24 元变为 39.67 元，累计 -12.3\%，区间最高 109.37 元（2025 年 6 月）、最低 36.30 元（2024 年 1 月）；当前价相当于区间最高价的 36.3\%，为区间最低价的 1.09 倍。 月度收益的年化波动率 44.6\%，高于全样本中位数 37.4\%，在三级行业“宠物食品”的 3 家公司中波动率由高到低排第\zhnumber{1}位。 把股价阶段与公司事件对齐看，2025 年 6 月 至 2026 年 9 月股价下跌-63.7\%，同期（2026 年 8 月）公司同日披露可转债募集资金使用可行性分析报告。 整体上看，区间的几处大级别拐点都能在公司的资本运作、股权变动或风险事件上找到对应，股价波动有明显的公司层面驱动，而非单纯的行业 beta。

\begin{center}
\begin{minipage}{\textwidth}\centering
\begin{tikzpicture}
\begin{axis}[
  width=12.6cm, height=3.9cm, scale only axis,
  xmin=2023.505, xmax=2026.888, ymin=33.76, ymax=117.03,
  xtick={2023,2024,2025,2026},
  yearticks,
  yticklabel style={font=\scriptsize},
  ylabel={元/股}, ylabel style={font=\scriptsize},
  grid=major, grid style={LightGray, line width=0.3pt},
  axis line style={InkGray, line width=0.5pt},
  every axis plot/.append style={line width=0.9pt},
]
\addplot[AgriGreen] table[x=x,y=y]{data/clean/profiles/px_301498.dat};
\addplot[SoftRed, dashed, line width=0.7pt] table[x=x,y=y]{data/clean/profiles/pxzz_301498.dat};
\addplot[only marks, mark=*, mark size=1.1pt, SoftRed] table[x=x,y=y]{data/clean/profiles/pxzz_301498.dat};
\end{axis}
\end{tikzpicture}

\captionof{figure}[乖宝宠物（301498）股价与周期骨架]{乖宝宠物（301498）月度收盘价（元/股）与 ZigZag 周期骨架（阈值 25\%，圆点为周期转折点）}
\end{minipage}
\end{center}

\pfl{过去周期分析} 以 25\% 的 ZigZag 阈值划分，样本区间内股价形成 1 个主升段与 2 个主跌段。 最大主跌段为 2025 年 6 月 至 2026 年 9 月，15 个月-63.7\%；最大主升段出现在 2024 年 1 月 至 2025 年 6 月，17 个月+201.3\%。 最近一次转折出现在 2026 年 9 月（下跌段自 2025 年 6 月起，15 个月-63.7\%），当前处于该段的延续之中。 公司股价较申万农林牧渔指数提前约 28 个月见底，是板块中较早定价周期反转的公司之一。 宠物食品属于消费型农业（第 \ref{sec:financing-strategy} 节归入“向消费端延伸”），收入增长由宠物数量与单宠消费驱动，与农产品价格周期的联动最弱，公司 2026 年上半年实现盈利、半年 ROE 4.1\%（未年化），好于全样本中位数（0.75\%），下行周期对其报表的冲击小于同业。 综合区间形态看，该股单段涨跌幅度偏大、以趋势段为主（最大主升 +201.3\%、最大主跌 -63.7\%），年化波动率 44.6\% 与全样本中位数 37.4\% 相比波动明显大于样本中位数。

\pfl{盈利情况} 2026 年上半年营业收入 35.43 亿元（同比 +10.0\%），归母净利润 1.91 亿元，同比 -49.6\%。 以年报为趋势对照，2023---2025 年营业收入由 43.27 亿元变为 67.69 亿元（两年年均复合增速 +25.1\%），归母净利润由 4.29 亿元变为 6.73 亿元。 按半年报口径，2026 年上半年的 ROE 为 4.1\%（未年化）、销售净利率 5.4\%、毛利率 40.1\%，ROE 在“宠物食品”3 家公司中按数值由高到低排第\zhnumber{1}位。 按同一口径，三级行业“宠物食品”3 家公司的半年 ROE 中位数为 4.00\%，该公司高于其中位数 0.06 个百分点。 从公司披露的原因看，业绩变动主要来自：公司持续扩大国内智能工厂产能；2026 年 4 月，公司加码泰国、新西兰高端产能布局并搭建 WCM 世界级制造体系。 综合看，公司在报告期维持盈利（高于全样本半年 ROE 中位数 0.75\%），利润的可持续性取决于宠物食品景气的延续与成本管控的执行。

\begin{center}\small\setlength{\tabcolsep}{3.4pt}
\captionof{table}[乖宝宠物（301498）关键财务指标]{乖宝宠物（301498）关键财务指标（合并报表口径；两个半年报期均未年化；现金短债比 $=$ 货币资金 $\div$ 短期债务）}
\begin{adjustbox}{max width=\textwidth}
\begin{tabular}{@{}lrrrrr@{}}
\toprule
指标 & 2023 年 & 2024 年 & 2025 年 & 2025 年半年报 & 2026 年半年报 \\
\midrule
营业收入（亿元） & 43.27 & 52.45 & 67.69 & 32.21 & 35.43 \\
归母净利润（亿元） & 4.29 & 6.25 & 6.73 & 3.78 & 1.91 \\
毛利率（\%） & 36.8 & 42.3 & 42.5 & 42.8 & 40.1 \\
ROE（\%） & 17.4 & 15.8 & 15.3 & 8.7 & 4.1 \\
资产负债率（\%） & 11.4 & 17.1 & 20.5 & 14.9 & 31.4 \\
经营活动现金流净额（亿元） & 6.17 & 7.20 & 8.42 & 3.50 & 1.51 \\
货币资金（亿元） & 10.09 & 6.95 & 8.19 & 6.33 & 9.85 \\
有息负债合计（亿元） & 0.25 & 1.86 & 3.70 & 1.48 & 11.94 \\
现金短债比（倍） & 216.75 & 5.71 & 2.61 & 7.68 & 1.06 \\
\bottomrule
\end{tabular}
\end{adjustbox}
\end{center}

\begin{center}
\begin{minipage}{\textwidth}\centering
\begin{tikzpicture}
\begin{groupplot}[
  group style={group size=2 by 1, horizontal sep=0.60cm},
  width=6.85cm, height=4.60cm,
  tick label style={font=\tiny},
  title style={font=\scriptsize\heiti},
  yticklabel style={font=\tiny},
  ylabel style={font=\tiny},
  grid=major, grid style={LightGray, line width=0.3pt},
  axis line style={InkGray, line width=0.5pt},
  enlarge x limits={abs=0.8},
  legend style={font=\scriptsize, at={(0.5,-0.20)}, anchor=north,
                legend columns=2, draw=none, fill=none, column sep=6pt},
]
\nextgroupplot[
  ybar, bar width=4pt,
  ymin=-6.77, ymax=75.81,
  xtick={0,1,2,3,4,5.7},
  xticklabels={2021,2022,2023,2024,2025,2026H1},
  ylabel={亿元},
]
\addplot[fill=AgriGreen!75, draw=AgriGreen, bar shift=-2.6pt] table[x=i, y=rev]{data/clean/profiles/fin_301498.dat};
\addplot[fill=SoftRed!60, draw=SoftRed, bar shift=2.6pt] table[x=i, y=ni]{data/clean/profiles/fin_301498.dat};
\addplot[fill=AgriGreen!30, draw=AgriGreen, pattern=north east lines, bar shift=-2.6pt] table[x=x, y=rev]{data/clean/profiles/finh_301498.dat};
\addplot[fill=SoftRed!20, draw=SoftRed, pattern=north east lines, bar shift=2.6pt] table[x=x, y=ni]{data/clean/profiles/finh_301498.dat};
\legend{营业收入（年/半年）, 归母净利润（年/半年）}
\nextgroupplot[
  ymin=-2.9, ymax=39.5,
  xtick={0,1,2,3,4,5.7},
  xticklabels={2021,2022,2023,2024,2025,2026H1},
  ylabel={\%},
]
\addplot[AgriGold, mark=*, mark size=1.3pt] table[x=i, y=roe]{data/clean/profiles/fin_301498.dat};
\addplot[OilBlue, mark=square*, mark size=1.3pt] table[x=i, y=debt]{data/clean/profiles/fin_301498.dat};
\addplot[only marks, AgriGold, mark=*, mark size=2.0pt] table[x=x, y=roe]{data/clean/profiles/finh_301498.dat};
\addplot[only marks, OilBlue, mark=square*, mark size=2.0pt] table[x=x, y=debt]{data/clean/profiles/finh_301498.dat};
\legend{ROE, 资产负债率}
\end{groupplot}
\end{tikzpicture}
\begin{tikzpicture}
\begin{axis}[
  name=finbot,
  width=13.75cm, height=3.10cm, scale only axis,
  ymin=-1.19, ymax=13.38,
  xtick={0,1,2,3,4,5.7},
  xticklabels={2021,2022,2023,2024,2025,2026H1},
  tick label style={font=\tiny},
  yticklabel style={font=\tiny},
  ylabel={亿元}, ylabel style={font=\tiny},
  ybar, bar width=4pt,
  grid=major, grid style={LightGray, line width=0.3pt},
  axis line style={InkGray, line width=0.5pt},
  enlarge x limits={abs=0.8},
  legend style={font=\scriptsize, at={(0.5,-0.26)}, anchor=north,
                legend columns=2, draw=none, fill=none, column sep=10pt},
]
\addplot[fill=AgriGreen!55, draw=AgriGreen, bar shift=-3.0pt] table[x=i, y=cash]{data/clean/profiles/fin_301498.dat};
\addplot[fill=SoftRed!45, draw=SoftRed, bar shift=3.0pt] table[x=i, y=ibd]{data/clean/profiles/fin_301498.dat};
\addplot[fill=AgriGreen!25, draw=AgriGreen, pattern=north east lines, bar shift=-3.0pt] table[x=x, y=cash]{data/clean/profiles/finh_301498.dat};
\addplot[fill=SoftRed!20, draw=SoftRed, pattern=north east lines, bar shift=3.0pt] table[x=x, y=ibd]{data/clean/profiles/finh_301498.dat};
\legend{货币资金（左轴）, 有息负债（左轴）}
\end{axis}
\begin{axis}[
  at={(finbot.south east)}, anchor=south east,
  width=13.75cm, height=3.10cm, scale only axis,
  axis y line*=right, axis x line*=none, xtick=\empty,
  ymin=-2.19, ymax=10.95,
  tick label style={font=\tiny},
  yticklabel style={font=\tiny}, scaled y ticks=false,
  legend style={font=\scriptsize, at={(0.5,-0.44)}, anchor=north,
                legend columns=1, draw=none, fill=none},
]
\addplot[OilBlue, mark=*, mark size=2.2pt, line width=1.3pt] table[x=i, y=ocf]{data/clean/profiles/fin_301498.dat};
\addplot[only marks, OilBlue, mark=*, mark size=3.2pt] table[x=x, y=ocf]{data/clean/profiles/finh_301498.dat};
\legend{经营活动现金流净额（右轴，亿元，蓝点）}
\end{axis}
\end{tikzpicture}

\captionof{figure}[乖宝宠物（301498）近年主要财务指标]{乖宝宠物（301498）近年主要财务指标（左上：营业收入与归母净利润；右上：ROE 与资产负债率；下：货币资金、有息负债为左轴柱状，经营活动现金流净额为其右轴的蓝点折线。
2021---2025 年为年报口径，2026H1 为 2026 年半年报/半年末，与其前一年报之间留出间隔、以斜纹或空心点区分；半年报数据未年化）}
\end{minipage}
\end{center}

\pfl{财务分析} 2026 年 6 月末资产负债率 31.4\%，较 2025 年末上升 11.0 个百分点（样本中位数 52.2\%，所处三级行业内按由低到高排第\zhnumber{1}位）。 2026 年 6 月末货币资金 9.85 亿元、短期债务（短期借款与一年内到期非流动负债）9.29 亿元、长期债务 2.65 亿元，有息负债合计 11.94 亿元，现金短债比 1.06 倍——覆盖充分。 净有息负债 2.10 亿元，相当于其 2026 年上半年营业收入的 6\%。 流动比率 2.13、速动比率 1.59，短期偿债能力处于正常区间。 2026 年上半年经营活动现金流净额 1.51 亿元，经营现金流与利润同向为正，内生资金可以支撑运营。 2026 年 6 月末在建工程 3.78 亿元，较上年末增加 3.19 亿元，仍有在建投入。 综合看，现金短债比 1.06 倍，短期偿付压力可控；资产负债率低于样本中位数 20.8 个百分点；有息负债中短期债务占比更高，期限结构仍以滚动续作为主。

\pfl{重大战略转型} 公告口径上，近三年（2023 年 9 月---2026 年 9 月）公司披露重大重组与并购类公告 0 条、对外投资与转型类公告 2 条、回购与股权激励类公告 60 条，合计公告 430 条。 第一大行业仍是“宠物食品及用品”，收入占比 99.6\%，与 2021 年年报相比没有方向性变化。 综合判断，报告期内公司有 2 条与战略相关的公告落地（重大重组与并购 0 条、对外投资与转型 2 条），资本运作与主业调整之间存在直接关联。

\begin{center}\small\setlength{\tabcolsep}{4pt}
\captionof{table}[乖宝宠物（301498）历史融资与资本运作]{乖宝宠物（301498）历史融资与资本运作（据公司公告与公开报道整理；金额为披露值，含首次公开发行、再融资、债券、银行借款与重整投资等）}
\begin{adjustbox}{max width=\textwidth}
\begin{tabular}{@{}p{1.45cm}p{2.55cm}p{4.05cm}p{4.95cm}@{}}
\toprule
时间 & 方式 & 规模 & 用途与说明 \\
\midrule
2023-08 & 首次公开发行A股股票并在深圳证券交易所创业板上市（IPO） & 发行价 39.99 元/\allowbreak{}股；首发前总股本 36，\allowbreak{}004.00 万股，\allowbreak{}实际发行 4 & 投向公司首次公开发行募集资金投资项目；超募资金部分用于投资建设新项目及永久补充流动资金。司与中泰证券股份有限公司；发行市盈率 61.28 倍；招股公告日 2023 年 7 月 27 日 \\
2023-09-19 & 超募资金使用（投资建设新项目及永久补充流动资金） & 拟使用部分超募资金投资建设新项目并用部分超募资金永久补充流动资金 & 新增募投项目及补充流动资金。核查意见 \\
2024-04-16 & 超募资金使用（永久补充流动资金） & 使用部分超募资金永久补充流动资金 & 永久补充流动资金。核查意见 \\
2025-08-10 & 股权激励（第二期限制性股票激励计划） & 公司董事会审议通过第二期限制性股票激励计划（草案）及其摘要等议案 & 用于核心人员激励。；2025 年 4 月 19 日、\allowbreak{}7 月 7 日曾两次调整第一期限制性股票激励计划授予价格并向激励对象授予预留部分限制性股票 \\
2026-04-23 & 募集资金使用进展（IPO募投项目） & 截至 2025 年 12 月 31 日公司累计使用募集资金 13.49 亿元 & 用于首发募投项目及超募资金安排项目。；公司同时披露 2025 年度募集资金存放与使用情况的专项报告及中泰证券核查意见；2026 年 4 月 23 日另公告使用部分闲置募集资金进行现金管理 \\
2026-08-27 & 拟向不特定对象发行可转换公司债券 & 募集资金总额不超过人民币 11.00 亿元，\allowbreak{}尚需履行相应审议与注册程序 & 年产 30 万吨高端宠物主粮食品项目（拟投入 4.50 亿元）、\allowbreak{}智能仓储及数智化分拣中心项目（拟投入 4.00 亿元）、\allowbreak{}数智化运营与决策平台建设项目（拟投入 1.00 亿元）、\allowbreak{}补充流动资金（拟投入 1.50 亿元）。；公司同日披露前次募集资金使用情况报告与德勤华永会计师事务所出具的前次募集资金使用情况审核报告 \\
\bottomrule
\end{tabular}
\end{adjustbox}
\end{center}

\pfl{历史融投资情况} 从现金流量表看，2025 年公司取得借款收到的现金 4.20 亿元、偿还债务支付的现金 2.29 亿元、吸收权益性投资 0.11 亿元，筹资活动现金流净额 -0.16 亿元（2023 年为取得借款 0.98 亿元、偿还债务 2.94 亿元）。 2026 年上半年筹资活动现金流净额为 5.06 亿元（取得借款 9.22 亿元、偿还债务 3.01 亿元，未年化），仍处于净融入状态。 2025 年分配股利、利润或偿付利息支付的现金合计 2.02 亿元。 公开披露的融资记录共 6 笔，工具类型以 IPO、可转债为主；金额与用途见下表。 其中规模最大的两笔为：2026 年 8 月，可转债，募集资金总额不超过人民币 11.00 亿元，尚需履行相应审议与注册程序；2026 年 4 月，募集资金使用进展，截至 2025 年 12 月 31 日公司累计使用募集资金 13.49 亿元。 股本方面，近三年披露股本变动 5 次（限售股份上市 2 次、A 股上市 1 次）；总股本自 2021 年末以来未发生变化（4.00 亿股）。 分红方面，近三年现金分红 7 次、合计每股派现 1.42 元，最近一次实施在 2026 年 9 月。 近三年“再融资”类公告 7 条，涉及可转债、募集资金使用。 综合看，公司融资结构上以债务滚动为主、股权融资为补充；2025 年筹资活动现金流净额 -0.16 亿元，融资节奏仍以维持存量债务周转为主，与前述经营与投资现金流的方向基本一致。

\pfl{未来潜在融资需求分析} 2026 年 6 月末货币资金 9.85 亿元，而短期债务 9.29 亿元（现金短债比 1.06 倍），2026 年上半年经营现金流净额 1.51 亿元。账面货币资金可以覆盖短期债务，缺口不体现在静态偿付层面。 公司 2026 年 6 月末资产负债率 31.4\%、2026 年上半年 ROE 4.1\%（未年化），资产负债表尚有余量，融资需求不是“补血型”的刚性需求，而是配套在建项目投入的主动性安排。 公开披露的后续资金安排线索包括：股权融资线索：2026 年 8 月，公司同日披露可转债募集资金使用可行性分析报告；债务与授信安排：2026 年 6 月，公司未公告银行综合授信额度、公司债或中票发行计划；2026 年 8 月，公司已披露向不特定对象发行可转换公司债券预案。 资金需求还要叠加在建工程：期末在建工程 3.78 亿元、2026 年上半年购建固定资产等支出 4.40 亿元（未年化），这部分支出在周期底部通常会被主动放缓。 以现有货币资金与经营现金流衡量，公司不存在刚性补血的紧迫性，融资动作更多是主动型的。




\subsubsection{中水渔业（000798）}
\label{co:000798}

\pfl{公司基本情况} 公司前身可追溯至 1997 年，中国农发集团旗下的远洋渔业上市平台，1998 年 2 月在深交所上市（原中水集团远洋股份有限公司），以金枪鱼等大洋性捕捞为基本盘；2023 年以 17.15 亿元现金收购中渔环球 51\% 等农发集团渔业板块资产并配套 11 亿元并购贷款，2025 年扭亏（归母净利润 8,375.42 万元），2026 年上半年归母净利润 1.02 亿元、同比增长 18.60\%。 公司于 1998-02-12 在深交所主板首发上市，注册资本 3.66 亿元，总股本 3.66 亿股。 截至 2026 年 9 月 24 日收盘价 9.98 元，流通市值 36.5 亿元，近一年+27.6\%，流通市值在三级行业“海洋捕捞”的 2 家公司中按由大到小排第\zhnumber{1}位，近一年涨跌幅高于全样本中位数（-10.1\%）37.7 个百分点。 按 2026 年半年报披露的行业构成，收入占比前三位为渔业服务 47.2\%；捕捞 41.7\%；零售及加工贸易 11.0\%。 与 2021 年年报相比，第一大行业口径由“捕捞收入”（77.4\%）变为“渔业服务”（47.2\%）；两期披露的分类口径有调整，占比差异中同时包含分类因素。 股权结构方面，实际控制人为中国农业发展集团有限公司。

\pfl{历史股价} 自 2021 年 1 月至 2026 年 9 月，股价由 8.19 元变为 9.98 元，累计 +21.9\%，区间最高 12.14 元（2022 年 8 月）、最低 5.63 元（2024 年 6 月）；当前价相当于区间最高价的 82.2\%，为区间最低价的 1.77 倍。 月度收益的年化波动率 45.4\%，高于全样本中位数 37.4\%，在三级行业“海洋捕捞”的 2 家公司中波动率由高到低排第\zhnumber{1}位。

\begin{center}
\begin{minipage}{\textwidth}\centering
\begin{tikzpicture}
\begin{axis}[
  width=12.6cm, height=3.9cm, scale only axis,
  xmin=2020.922, xmax=2026.888, ymin=5.24, ymax=12.99,
  xtick={2021,2022,2023,2024,2025,2026},
  yearticks,
  yticklabel style={font=\scriptsize},
  ylabel={元/股}, ylabel style={font=\scriptsize},
  grid=major, grid style={LightGray, line width=0.3pt},
  axis line style={InkGray, line width=0.5pt},
  every axis plot/.append style={line width=0.9pt},
]
\addplot[AgriGreen] table[x=x,y=y]{data/clean/profiles/px_000798.dat};
\addplot[SoftRed, dashed, line width=0.7pt] table[x=x,y=y]{data/clean/profiles/pxzz_000798.dat};
\addplot[only marks, mark=*, mark size=1.1pt, SoftRed] table[x=x,y=y]{data/clean/profiles/pxzz_000798.dat};
\end{axis}
\end{tikzpicture}

\captionof{figure}[中水渔业（000798）股价与周期骨架]{中水渔业（000798）月度收盘价（元/股）与 ZigZag 周期骨架（阈值 25\%，圆点为周期转折点）}
\end{minipage}
\end{center}

\pfl{过去周期分析} 以 25\% 的 ZigZag 阈值划分，样本区间内股价形成 3 个主升段与 3 个主跌段。 最大主跌段为 2022 年 8 月 至 2024 年 6 月，22 个月-53.6\%；最大主升段出现在 2024 年 6 月 至 2025 年 11 月，17 个月+115.6\%。 最近一次转折出现在 2026 年 9 月（上涨段自 2026 年 6 月起，3 个月+41.4\%），当前处于该段之后的震荡之中。 公司股价较申万农林牧渔指数提前约 23 个月见底，是板块中较早定价周期反转的公司之一。 远洋与近海捕捞的产量受资源与配额约束，成本端则跟随燃油与人工，公司 2026 年 6 月末资产负债率 71.6\%，高于全样本中位数 52.2\%，其股价因此更多反映资金面与信用风险，而不是价格弹性本身。 综合区间形态看，该股单段涨跌幅度偏大、以趋势段为主（最大主升 +115.6\%、最大主跌 -53.6\%），年化波动率 45.4\% 与全样本中位数 37.4\% 相比波动明显大于样本中位数。

\pfl{盈利情况} 2026 年上半年营业收入 20.78 亿元（同比 +18.9\%），归母净利润 1.02 亿元，同比 +18.6\%。 以年报为趋势对照，2023---2025 年营业收入由 40.42 亿元变为 43.42 亿元（两年年均复合增速 +3.6\%），归母净利润由 -1.18 亿元变为 0.84 亿元。 按半年报口径，2026 年上半年的 ROE 为 20.9\%（未年化）、销售净利率 6.8\%、毛利率 5.5\%，ROE 在“海洋捕捞”2 家公司中按数值由高到低排第\zhnumber{1}位。 按同一口径，三级行业“海洋捕捞”2 家公司的半年 ROE 中位数为 9.79\%，该公司高于其中位数 11.12 个百分点。 从公司披露的原因看，业绩变动主要来自：2025 年，归属于上市公司股东的净利润 8，375.00 万元、同比增加 1.88 亿元。 综合看，公司在报告期维持盈利（高于全样本半年 ROE 中位数 0.75\%），利润的可持续性取决于海洋捕捞景气的延续与成本管控的执行。

\begin{center}\small\setlength{\tabcolsep}{3.4pt}
\captionof{table}[中水渔业（000798）关键财务指标]{中水渔业（000798）关键财务指标（合并报表口径；两个半年报期均未年化；现金短债比 $=$ 货币资金 $\div$ 短期债务）}
\begin{adjustbox}{max width=\textwidth}
\begin{tabular}{@{}lrrrrr@{}}
\toprule
指标 & 2023 年 & 2024 年 & 2025 年 & 2025 年半年报 & 2026 年半年报 \\
\midrule
营业收入（亿元） & 40.42 & 49.33 & 43.42 & 17.47 & 20.78 \\
归母净利润（亿元） & -1.18 & -1.04 & 0.84 & 0.86 & 1.02 \\
毛利率（\%） & 5.1 & 3.2 & 8.3 & 5.2 & 5.5 \\
ROE（\%） & -8.3 & -25.4 & 21.1 & 21.3 & 20.9 \\
资产负债率（\%） & 72.5 & 73.8 & 72.9 & 72.4 & 71.6 \\
经营活动现金流净额（亿元） & -1.87 & 5.25 & 4.70 & 3.78 & -0.17 \\
货币资金（亿元） & 6.86 & 8.92 & 8.28 & 11.53 & 8.75 \\
有息负债合计（亿元） & 24.60 & 28.07 & 26.28 & 27.57 & 27.86 \\
现金短债比（倍） & 0.81 & 0.56 & 0.55 & 0.70 & 0.52 \\
\bottomrule
\end{tabular}
\end{adjustbox}
\end{center}

\begin{center}
\begin{minipage}{\textwidth}\centering
\begin{tikzpicture}
\begin{groupplot}[
  group style={group size=2 by 1, horizontal sep=0.60cm},
  width=6.85cm, height=4.60cm,
  tick label style={font=\tiny},
  title style={font=\scriptsize\heiti},
  yticklabel style={font=\tiny},
  ylabel style={font=\tiny},
  grid=major, grid style={LightGray, line width=0.3pt},
  axis line style={InkGray, line width=0.5pt},
  enlarge x limits={abs=0.8},
  legend style={font=\scriptsize, at={(0.5,-0.20)}, anchor=north,
                legend columns=2, draw=none, fill=none, column sep=6pt},
]
\nextgroupplot[
  ybar, bar width=4pt,
  ymin=-6.23, ymax=55.39,
  xtick={0,1,2,3,4,5.7},
  xticklabels={2021,2022,2023,2024,2025,2026H1},
  ylabel={亿元},
]
\addplot[fill=AgriGreen!75, draw=AgriGreen, bar shift=-2.6pt] table[x=i, y=rev]{data/clean/profiles/fin_000798.dat};
\addplot[fill=SoftRed!60, draw=SoftRed, bar shift=2.6pt] table[x=i, y=ni]{data/clean/profiles/fin_000798.dat};
\addplot[fill=AgriGreen!30, draw=AgriGreen, pattern=north east lines, bar shift=-2.6pt] table[x=x, y=rev]{data/clean/profiles/finh_000798.dat};
\addplot[fill=SoftRed!20, draw=SoftRed, pattern=north east lines, bar shift=2.6pt] table[x=x, y=ni]{data/clean/profiles/finh_000798.dat};
\legend{营业收入（年/半年）, 归母净利润（年/半年）}
\nextgroupplot[
  ymin=-33.4, ymax=83.7,
  xtick={0,1,2,3,4,5.7},
  xticklabels={2021,2022,2023,2024,2025,2026H1},
  ylabel={\%},
]
\addplot[AgriGold, mark=*, mark size=1.3pt] table[x=i, y=roe]{data/clean/profiles/fin_000798.dat};
\addplot[OilBlue, mark=square*, mark size=1.3pt] table[x=i, y=debt]{data/clean/profiles/fin_000798.dat};
\addplot[only marks, AgriGold, mark=*, mark size=2.0pt] table[x=x, y=roe]{data/clean/profiles/finh_000798.dat};
\addplot[only marks, OilBlue, mark=square*, mark size=2.0pt] table[x=x, y=debt]{data/clean/profiles/finh_000798.dat};
\legend{ROE, 资产负债率}
\end{groupplot}
\end{tikzpicture}
\begin{tikzpicture}
\begin{axis}[
  name=finbot,
  width=13.75cm, height=3.10cm, scale only axis,
  ymin=-2.81, ymax=31.44,
  xtick={0,1,2,3,4,5.7},
  xticklabels={2021,2022,2023,2024,2025,2026H1},
  tick label style={font=\tiny},
  yticklabel style={font=\tiny},
  ylabel={亿元}, ylabel style={font=\tiny},
  ybar, bar width=4pt,
  grid=major, grid style={LightGray, line width=0.3pt},
  axis line style={InkGray, line width=0.5pt},
  enlarge x limits={abs=0.8},
  legend style={font=\scriptsize, at={(0.5,-0.26)}, anchor=north,
                legend columns=2, draw=none, fill=none, column sep=10pt},
]
\addplot[fill=AgriGreen!55, draw=AgriGreen, bar shift=-3.0pt] table[x=i, y=cash]{data/clean/profiles/fin_000798.dat};
\addplot[fill=SoftRed!45, draw=SoftRed, bar shift=3.0pt] table[x=i, y=ibd]{data/clean/profiles/fin_000798.dat};
\addplot[fill=AgriGreen!25, draw=AgriGreen, pattern=north east lines, bar shift=-3.0pt] table[x=x, y=cash]{data/clean/profiles/finh_000798.dat};
\addplot[fill=SoftRed!20, draw=SoftRed, pattern=north east lines, bar shift=3.0pt] table[x=x, y=ibd]{data/clean/profiles/finh_000798.dat};
\legend{货币资金（左轴）, 有息负债（左轴）}
\end{axis}
\begin{axis}[
  at={(finbot.south east)}, anchor=south east,
  width=13.75cm, height=3.10cm, scale only axis,
  axis y line*=right, axis x line*=none, xtick=\empty,
  ymin=-3.72, ymax=7.39,
  tick label style={font=\tiny},
  yticklabel style={font=\tiny}, scaled y ticks=false,
  legend style={font=\scriptsize, at={(0.5,-0.44)}, anchor=north,
                legend columns=1, draw=none, fill=none},
]
\addplot[OilBlue, mark=*, mark size=2.2pt, line width=1.3pt] table[x=i, y=ocf]{data/clean/profiles/fin_000798.dat};
\addplot[only marks, OilBlue, mark=*, mark size=3.2pt] table[x=x, y=ocf]{data/clean/profiles/finh_000798.dat};
\legend{经营活动现金流净额（右轴，亿元，蓝点）}
\end{axis}
\end{tikzpicture}

\captionof{figure}[中水渔业（000798）近年主要财务指标]{中水渔业（000798）近年主要财务指标（左上：营业收入与归母净利润；右上：ROE 与资产负债率；下：货币资金、有息负债为左轴柱状，经营活动现金流净额为其右轴的蓝点折线。
2021---2025 年为年报口径，2026H1 为 2026 年半年报/半年末，与其前一年报之间留出间隔、以斜纹或空心点区分；半年报数据未年化）}
\end{minipage}
\end{center}

\pfl{财务分析} 2026 年 6 月末资产负债率 71.6\%，较 2025 年末下降 1.3 个百分点（样本中位数 52.2\%，所处三级行业内按由低到高排第\zhnumber{2}位）。 2026 年 6 月末货币资金 8.75 亿元、短期债务（短期借款与一年内到期非流动负债）16.75 亿元、长期债务 11.11 亿元，有息负债合计 27.86 亿元，现金短债比 0.52 倍——覆盖偏紧。 净有息负债 19.11 亿元，相当于其 2026 年上半年营业收入的 92\%。 流动比率 1.02、速动比率 0.44，短期偿债能力处于正常区间。 2026 年上半年经营活动现金流净额 -0.17 亿元，净利润为正但经营现金流为负，利润的现金含量偏低。 综合看，现金短债比 0.52 倍，短期资金安排偏紧；资产负债率高于样本中位数 19.4 个百分点；有息负债中短期债务占比更高，期限结构仍以滚动续作为主。 公司披露的重整与债务违约风险为：2026 年 6 月，公司因张福赐未完成 2015 年业绩承诺要求张福赐与张福庆承担违约责任。

\pfl{重大战略转型} 近三年公司的战略动作集中在：并购与重组（2 项）：2023 年 6 月，上市公司在远洋捕捞、水产品贸易的基础上新增水产品加工、渔业服务等多元化业务；2020 年，公司收购超低温金枪鱼项目、丰汇项目后与原低温金枪鱼项目实现深度融合。 公告口径上，近三年（2023 年 9 月---2026 年 9 月）公司披露重大重组与并购类公告 4 条、对外投资与转型类公告 0 条、回购与股权激励类公告 0 条，合计公告 332 条。 主营结构的口径变化已经落到报表上：2021 年年报的第一大行业为“捕捞收入”（77.4\%），2026 年半年报为“渔业服务”（47.2\%）；公司在这两期的分部划分方式有过调整。 综合判断，报告期内公司有 4 条与战略相关的公告落地（重大重组与并购 4 条、对外投资与转型 0 条），资本运作与主业调整之间存在直接关联。

\begin{center}\small\setlength{\tabcolsep}{4pt}
\captionof{table}[中水渔业（000798）历史融资与资本运作]{中水渔业（000798）历史融资与资本运作（据公司公告与公开报道整理；金额为披露值，含首次公开发行、再融资、债券、银行借款与重整投资等）}
\begin{adjustbox}{max width=\textwidth}
\begin{tabular}{@{}p{1.45cm}p{2.55cm}p{4.05cm}p{4.95cm}@{}}
\toprule
时间 & 方式 & 规模 & 用途与说明 \\
\midrule
1997-11-24 & IPO（深圳证券交易所） & 发行 6,\allowbreak{}300 万股，\allowbreak{}发行价 7.24 元/\allowbreak{}股，\allowbreak{}实际募集资金 4.56 亿元 & 经证监发字[1997]480 号、\allowbreak{}证监发字[1997]481 号文批准 \\
2014-12 & 现金收购股权（并购） & 以现金购买张福赐持有的厦门新阳洲水产品工贸有限公司 55\% 股权 & 延伸水产品加工与贸易业务；交易对方张福赐承诺新阳洲 2014—2017 年净利润分别不低于 3 \\
2018-10 & 对外投资（关联交易） & 项目总投资 4.20 亿元，\allowbreak{}中国农发集团以现金出资 1.00 亿元占 23.81\% 股权 & 与控股股东及其关联方共同投资渔业相关项目；中水渔业与舟渔公司均系中国农发集团直接控制的企业 \\
2023 & 并购贷款 & 拟向合作银行申请总额 11.00 亿元人民币的并购贷款，\allowbreak{}用于支付本次交易的部分转让现金对价 & 支付重大资产购买对价；根据备考审阅报告，\allowbreak{}2022 年末公司模拟账面货币资金为 6.03 亿元；本次交易完成后公司资产负债率将从 28.97\% 大幅上升至 68.57\% \\
2023-06-15 & 重大资产购买（现金）暨关联交易 & 交易价格合计 17.15 亿元 & 整合中国农发集团渔业板块 \\
\bottomrule
\end{tabular}
\end{adjustbox}
\end{center}

\pfl{历史融投资情况} 从现金流量表看，2025 年公司取得借款收到的现金 20.94 亿元、偿还债务支付的现金 21.22 亿元，筹资活动现金流净额 -2.03 亿元（2023 年为取得借款 33.27 亿元、偿还债务 12.32 亿元）。 2026 年上半年筹资活动现金流净额为 1.23 亿元（取得借款 9.68 亿元、偿还债务 7.89 亿元，未年化），仍处于净融入状态。 2025 年分配股利、利润或偿付利息支付的现金合计 1.16 亿元。 公开披露的融资记录共 5 笔，工具类型以 IPO、银行借款为主；金额与用途见下表。 其中规模最大的两笔为：2023 年，银行借款，拟向合作银行申请总额 11.00 亿元人民币的并购贷款，用于支付本次交易的部分转让现金对价；1997 年 11 月，IPO，发行 6,300 万股，发行价 7.24 元/股，实际募集资金 4.56 亿元。 股本方面，近三年披露股本变动 4 次（限售股份上市 1 次）；总股本由 2021 年末的 3.19 亿股变为 3.66 亿股（+14.53\%）。 分红方面，近三年未实施现金分红，在全部样本中属于分红能力偏弱的一端。 综合看，公司融资结构上股权与债务融资并举；2025 年筹资活动现金流净额 -2.03 亿元，融资节奏仍以维持存量债务周转为主，与前述经营与投资现金流的方向基本一致。

\pfl{未来潜在融资需求分析} 2026 年 6 月末货币资金 8.75 亿元，而短期债务 16.75 亿元（现金短债比 0.52 倍），2026 年上半年经营现金流净额 -0.17 亿元。按“短期债务减去账面货币资金”的滚动偿付口径，静态缺口约 8.00 亿元；上半年盈利或接近盈亏平衡，该缺口不随经营亏损进一步扩大。 公司 2026 年 6 月末资产负债率 71.6\%、2026 年上半年 ROE 20.9\%（未年化），资产负债率偏高（71.6\%），融资需求首先用于置换与滚动存量债务、补充营运资金，属于“补血型”而非扩张型。 公开披露的后续资金安排线索包括：债务与授信安排：2023 年，公司以支付现金方式实施 17.15 亿元重大资产购买并申请 11.00 亿元并购贷款。 综合看，公司的外部资金需求以营运资金周转与在建项目投入为主，滚动偿付口径下的静态缺口约 8.00 亿元，能否落地取决于授信续作与发行窗口的配合。




\subsubsection{好当家（600467）}
\label{co:600467}

\pfl{公司基本情况} 山东荣成的海参养殖与海洋食品加工企业，2004 年 4 月在上交所上市，拥有全国最大海参育苗与养殖基地之一；2025 年受海参价格与需求影响营业收入 10.51 亿元、同比下降 25.76\%，归母净利润 2,987.23 万元、同比下降 31.21\%，2026 年上半年营收 8.03 亿元、归母净利润 2,997.88 万元，同时推进不超过 1 亿美元的境外债券发行与不超过 35 亿元的年度融资计划。 公司于 2004-04-05 在上交所首发上市，注册资本 14.61 亿元，总股本 14.61 亿股。 截至 2026 年 9 月 24 日收盘价 2.38 元，流通市值 34.8 亿元，近一年-0.8\%，流通市值在三级行业“水产养殖”的 4 家公司中按由大到小排第\zhnumber{1}位，近一年涨跌幅高于全样本中位数（-10.1\%）9.3 个百分点。 按 2025 年年报披露的行业构成，收入占比前三位为食品加工业 57.1\%；养殖业 30.4\%；捕捞业 11.7\%。

\pfl{历史股价} 自 2021 年 1 月至 2026 年 9 月，股价由 2.50 元变为 2.38 元，累计 -4.8\%，区间最高 3.16 元（2022 年 3 月）、最低 1.42 元（2024 年 6 月）；当前价相当于区间最高价的 75.3\%，为区间最低价的 1.68 倍。 月度收益的年化波动率 29.0\%，低于全样本中位数 37.4\%，在三级行业“水产养殖”的 4 家公司中波动率由高到低排第\zhnumber{4}位。

\begin{center}
\begin{minipage}{\textwidth}\centering
\begin{tikzpicture}
\begin{axis}[
  width=12.6cm, height=3.9cm, scale only axis,
  xmin=2020.922, xmax=2026.888, ymin=1.32, ymax=3.38,
  xtick={2021,2022,2023,2024,2025,2026},
  yearticks,
  yticklabel style={font=\scriptsize},
  ylabel={元/股}, ylabel style={font=\scriptsize},
  grid=major, grid style={LightGray, line width=0.3pt},
  axis line style={InkGray, line width=0.5pt},
  every axis plot/.append style={line width=0.9pt},
]
\addplot[AgriGreen] table[x=x,y=y]{data/clean/profiles/px_600467.dat};
\addplot[SoftRed, dashed, line width=0.7pt] table[x=x,y=y]{data/clean/profiles/pxzz_600467.dat};
\addplot[only marks, mark=*, mark size=1.1pt, SoftRed] table[x=x,y=y]{data/clean/profiles/pxzz_600467.dat};
\end{axis}
\end{tikzpicture}

\captionof{figure}[好当家（600467）股价与周期骨架]{好当家（600467）月度收盘价（元/股）与 ZigZag 周期骨架（阈值 25\%，圆点为周期转折点）}
\end{minipage}
\end{center}

\pfl{过去周期分析} 以 25\% 的 ZigZag 阈值划分，样本区间内股价形成 2 个主升段与 2 个主跌段。 最大主升段出现在 2024 年 6 月 至 2026 年 1 月，19 个月+93.7\%；最大主跌段为 2022 年 3 月 至 2024 年 6 月，27 个月-55.1\%。 最近一次转折出现在 2026 年 6 月（下跌段自 2026 年 1 月起，5 个月-28.4\%），当前处于该段的延续之中。 公司股价较申万农林牧渔指数提前约 23 个月见底，是板块中较早定价周期反转的公司之一。 水产养殖的价格周期由投苗量与消费需求共同决定（第 \ref{sec:other-livestock} 节），与猪周期的同步性弱于禽链，公司 2026 年上半年实现盈利、半年 ROE 0.9\%（未年化），好于全样本中位数（0.75\%），下行周期对其报表的冲击小于同业。 综合区间形态看，该股单段涨跌幅度偏大、以趋势段为主（最大主升 +93.7\%、最大主跌 -55.1\%），年化波动率 29.0\% 与全样本中位数 37.4\% 相比波动小于样本中位数。

\pfl{盈利情况} 2026 年上半年营业收入 8.03 亿元（同比 +13.0\%），归母净利润 0.30 亿元，同比 +21.7\%。 以年报为趋势对照，2023---2025 年营业收入由 15.64 亿元变为 10.51 亿元（两年年均复合增速 -18.0\%），归母净利润由 0.49 亿元变为 0.30 亿元。 按半年报口径，2026 年上半年的 ROE 为 0.9\%（未年化）、销售净利率 3.8\%、毛利率 17.2\%，ROE 在“水产养殖”4 家公司中按数值由高到低排第\zhnumber{4}位。 按同一口径，三级行业“水产养殖”4 家公司的半年 ROE 中位数为 6.05\%，该公司低于其中位数 5.18 个百分点。 综合看，公司在报告期维持盈利（高于全样本半年 ROE 中位数 0.75\%），利润的可持续性取决于水产养殖景气的延续与成本管控的执行。

\begin{center}\small\setlength{\tabcolsep}{3.4pt}
\captionof{table}[好当家（600467）关键财务指标]{好当家（600467）关键财务指标（合并报表口径；两个半年报期均未年化；现金短债比 $=$ 货币资金 $\div$ 短期债务）}
\begin{adjustbox}{max width=\textwidth}
\begin{tabular}{@{}lrrrrr@{}}
\toprule
指标 & 2023 年 & 2024 年 & 2025 年 & 2025 年半年报 & 2026 年半年报 \\
\midrule
营业收入（亿元） & 15.64 & 14.16 & 10.51 & 7.11 & 8.03 \\
归母净利润（亿元） & 0.49 & 0.43 & 0.30 & 0.25 & 0.30 \\
毛利率（\%） & 16.5 & 15.6 & 23.9 & 20.1 & 17.2 \\
ROE（\%） & 1.5 & 1.3 & 0.9 & 0.8 & 0.9 \\
资产负债率（\%） & 49.1 & 49.1 & 48.5 & 51.4 & 49.1 \\
经营活动现金流净额（亿元） & 8.68 & 5.88 & 4.51 & 0.59 & 1.08 \\
货币资金（亿元） & 4.80 & 4.14 & 4.40 & 6.25 & 5.64 \\
有息负债合计（亿元） & 23.87 & 25.92 & 27.10 & 29.28 & 27.09 \\
现金短债比（倍） & 0.21 & 0.18 & 0.18 & 0.25 & 0.24 \\
\bottomrule
\end{tabular}
\end{adjustbox}
\end{center}

\begin{center}
\begin{minipage}{\textwidth}\centering
\begin{tikzpicture}
\begin{groupplot}[
  group style={group size=2 by 1, horizontal sep=0.60cm},
  width=6.85cm, height=4.60cm,
  tick label style={font=\tiny},
  title style={font=\scriptsize\heiti},
  yticklabel style={font=\tiny},
  ylabel style={font=\tiny},
  grid=major, grid style={LightGray, line width=0.3pt},
  axis line style={InkGray, line width=0.5pt},
  enlarge x limits={abs=0.8},
  legend style={font=\scriptsize, at={(0.5,-0.20)}, anchor=north,
                legend columns=2, draw=none, fill=none, column sep=6pt},
]
\nextgroupplot[
  ybar, bar width=4pt,
  ymin=-1.56, ymax=17.51,
  xtick={0,1,2,3,4,5.7},
  xticklabels={2021,2022,2023,2024,2025,2026H1},
  ylabel={亿元},
]
\addplot[fill=AgriGreen!75, draw=AgriGreen, bar shift=-2.6pt] table[x=i, y=rev]{data/clean/profiles/fin_600467.dat};
\addplot[fill=SoftRed!60, draw=SoftRed, bar shift=2.6pt] table[x=i, y=ni]{data/clean/profiles/fin_600467.dat};
\addplot[fill=AgriGreen!30, draw=AgriGreen, pattern=north east lines, bar shift=-2.6pt] table[x=x, y=rev]{data/clean/profiles/finh_600467.dat};
\addplot[fill=SoftRed!20, draw=SoftRed, pattern=north east lines, bar shift=2.6pt] table[x=x, y=ni]{data/clean/profiles/finh_600467.dat};
\legend{营业收入（年/半年）, 归母净利润（年/半年）}
\nextgroupplot[
  ymin=-4.0, ymax=55.6,
  xtick={0,1,2,3,4,5.7},
  xticklabels={2021,2022,2023,2024,2025,2026H1},
  ylabel={\%},
]
\addplot[AgriGold, mark=*, mark size=1.3pt] table[x=i, y=roe]{data/clean/profiles/fin_600467.dat};
\addplot[OilBlue, mark=square*, mark size=1.3pt] table[x=i, y=debt]{data/clean/profiles/fin_600467.dat};
\addplot[only marks, AgriGold, mark=*, mark size=2.0pt] table[x=x, y=roe]{data/clean/profiles/finh_600467.dat};
\addplot[only marks, OilBlue, mark=square*, mark size=2.0pt] table[x=x, y=debt]{data/clean/profiles/finh_600467.dat};
\legend{ROE, 资产负债率}
\end{groupplot}
\end{tikzpicture}
\begin{tikzpicture}
\begin{axis}[
  name=finbot,
  width=13.75cm, height=3.10cm, scale only axis,
  ymin=-2.78, ymax=31.17,
  xtick={0,1,2,3,4,5.7},
  xticklabels={2021,2022,2023,2024,2025,2026H1},
  tick label style={font=\tiny},
  yticklabel style={font=\tiny},
  ylabel={亿元}, ylabel style={font=\tiny},
  ybar, bar width=4pt,
  grid=major, grid style={LightGray, line width=0.3pt},
  axis line style={InkGray, line width=0.5pt},
  enlarge x limits={abs=0.8},
  legend style={font=\scriptsize, at={(0.5,-0.26)}, anchor=north,
                legend columns=2, draw=none, fill=none, column sep=10pt},
]
\addplot[fill=AgriGreen!55, draw=AgriGreen, bar shift=-3.0pt] table[x=i, y=cash]{data/clean/profiles/fin_600467.dat};
\addplot[fill=SoftRed!45, draw=SoftRed, bar shift=3.0pt] table[x=i, y=ibd]{data/clean/profiles/fin_600467.dat};
\addplot[fill=AgriGreen!25, draw=AgriGreen, pattern=north east lines, bar shift=-3.0pt] table[x=x, y=cash]{data/clean/profiles/finh_600467.dat};
\addplot[fill=SoftRed!20, draw=SoftRed, pattern=north east lines, bar shift=3.0pt] table[x=x, y=ibd]{data/clean/profiles/finh_600467.dat};
\legend{货币资金（左轴）, 有息负债（左轴）}
\end{axis}
\begin{axis}[
  at={(finbot.south east)}, anchor=south east,
  width=13.75cm, height=3.10cm, scale only axis,
  axis y line*=right, axis x line*=none, xtick=\empty,
  ymin=-2.26, ymax=11.28,
  tick label style={font=\tiny},
  yticklabel style={font=\tiny}, scaled y ticks=false,
  legend style={font=\scriptsize, at={(0.5,-0.44)}, anchor=north,
                legend columns=1, draw=none, fill=none},
]
\addplot[OilBlue, mark=*, mark size=2.2pt, line width=1.3pt] table[x=i, y=ocf]{data/clean/profiles/fin_600467.dat};
\addplot[only marks, OilBlue, mark=*, mark size=3.2pt] table[x=x, y=ocf]{data/clean/profiles/finh_600467.dat};
\legend{经营活动现金流净额（右轴，亿元，蓝点）}
\end{axis}
\end{tikzpicture}

\captionof{figure}[好当家（600467）近年主要财务指标]{好当家（600467）近年主要财务指标（左上：营业收入与归母净利润；右上：ROE 与资产负债率；下：货币资金、有息负债为左轴柱状，经营活动现金流净额为其右轴的蓝点折线。
2021---2025 年为年报口径，2026H1 为 2026 年半年报/半年末，与其前一年报之间留出间隔、以斜纹或空心点区分；半年报数据未年化）}
\end{minipage}
\end{center}

\pfl{财务分析} 2026 年 6 月末资产负债率 49.1\%，较 2025 年末上升 0.6 个百分点（样本中位数 52.2\%，所处三级行业内按由低到高排第\zhnumber{2}位）。 2026 年 6 月末货币资金 5.64 亿元、短期债务（短期借款与一年内到期非流动负债）23.80 亿元、长期债务 3.29 亿元，有息负债合计 27.09 亿元，现金短债比 0.24 倍——缺口明显。 净有息负债 21.45 亿元，相当于其 2026 年上半年营业收入的 267\%。 流动比率 0.51、速动比率 0.26，短期偿债指标偏紧。 2026 年上半年经营活动现金流净额 1.08 亿元，经营现金流与利润同向为正，内生资金可以支撑运营。 综合看，账面货币资金对有息负债与短期债务的覆盖不足，滚动偿付对新增融资的依赖度较高；资产负债率低于样本中位数 3.1 个百分点；有息负债中短期债务占比更高，期限结构仍以滚动续作为主。 公司披露的股东与质押风险为：2026 年 1 月，控股股东好当家集团有限公司持有公司股份累计质押数量（含本次）3.46 亿股。

\pfl{重大战略转型} 近三年公司的战略动作集中在：产能与项目（1 项）：2024 年 12 月，公司建设并投产 130 多万平方水体现代海参育苗车间；融资与资本运作（1 项）：2026 年 4 月，公司拟委任天风国际证券与期货有限公司担任本次境外债券的和牵头安排行。 公告口径上，近三年（2023 年 9 月---2026 年 9 月）公司披露重大重组与并购类公告 0 条、对外投资与转型类公告 0 条、回购与股权激励类公告 0 条，合计公告 195 条。 第一大行业“食品加工业”的收入占比由 2021 年年报的 44.0\% 变为 57.1\%（13.0 个百分点），收入结构出现再平衡。 综合判断，公司报告期内的变化主要发生在收入结构层面，属于口径与业务重心的再平衡，而非外延式扩张。

\begin{center}\small\setlength{\tabcolsep}{4pt}
\captionof{table}[好当家（600467）历史融资与资本运作]{好当家（600467）历史融资与资本运作（据公司公告与公开报道整理；金额为披露值，含首次公开发行、再融资、债券、银行借款与重整投资等）}
\begin{adjustbox}{max width=\textwidth}
\begin{tabular}{@{}p{1.45cm}p{2.55cm}p{4.05cm}p{4.95cm}@{}}
\toprule
时间 & 方式 & 规模 & 用途与说明 \\
\midrule
2004-04-05 & IPO（上海证券交易所主板） & 发行 6，\allowbreak{}000 万股 & 经证监发行字[2004]26 号文批准 \\
2011-11-23 & 非公开发行股票（向特定对象） & 发行 9，\allowbreak{}689.72 万股 & 公司于 2011 年 4 月 29 日至 2011 年 11 月 23 日期间以自有资金先期投入募投项目 \\
2015 & 资本公积转增股本（权益融资之外股本扩张） & 每 10 股转增 10 股，\allowbreak{}共计转增股本 7.30 亿元 & 扩大股本规模；转增后公司注册资本为人民币 14.61 亿元 \\
2026-01-13 & 控股股东股权质押融资 & 好当家集团有限公司持有公司股份 5.73 亿股（占 39.25\%）；2025 年 1 月 10 日其将 730.00 万股无限售流通股质押给交通银行股份有限公司威海分行提供质押担保 & 为控股股东融资提供质押担保；2025 年 12 月控股股东再将部分股份质押给交通银行威海分行（占公司总股本 0.68\%） \\
2026-04-22 & 年度融资计划与境外债券（董事会授权） & 母公司及控股子公司向国内外银行、\allowbreak{}融资租赁公司、\allowbreak{}商业保理公司及国内外资本市场融资静态余额拟不超过 35.00 亿元 & 满足 2026 年度生产经营需要 \\
\bottomrule
\end{tabular}
\end{adjustbox}
\end{center}

\pfl{历史融投资情况} 从现金流量表看，2025 年公司取得借款收到的现金 35.09 亿元、偿还债务支付的现金 33.44 亿元，筹资活动现金流净额 -1.23 亿元（2023 年为取得借款 27.54 亿元、偿还债务 29.74 亿元）。 2026 年上半年筹资活动现金流净额为 -0.11 亿元（取得借款 14.73 亿元、偿还债务 12.95 亿元，未年化），已转为净流出，处于净还债状态。 2025 年分配股利、利润或偿付利息支付的现金合计 1.75 亿元。 公开披露的融资记录共 5 笔，工具类型以 IPO、定向增发为主；金额与用途见下表。 其中规模最大的两笔为：2026 年 4 月，年度融资计划与境外债，母公司及控股子公司向国内外银行、融资租赁公司、商业保理公司及国内外资本市场融资静态余额拟不超过 35.00 亿元；2011 年 11 月，定向增发，发行 9，689.72 万股。 股本方面，近三年披露股本变动 3 次；总股本自 2021 年末以来未发生变化（14.61 亿股）。 分红方面，近三年现金分红 4 次、合计每股派现 0.04 元，最近一次实施在 2025 年 12 月。 综合看，公司融资结构上股权与债务融资并举；2025 年筹资活动现金流净额 -1.23 亿元，融资节奏仍以维持存量债务周转为主，与前述经营与投资现金流的方向基本一致。

\pfl{未来潜在融资需求分析} 2026 年 6 月末货币资金 5.64 亿元，而短期债务 23.80 亿元（现金短债比 0.24 倍），2026 年上半年经营现金流净额 1.08 亿元。按“短期债务减去账面货币资金”的滚动偿付口径，静态缺口约 18.16 亿元；上半年盈利或接近盈亏平衡，该缺口不随经营亏损进一步扩大。 公司 2026 年 6 月末资产负债率 49.1\%、2026 年上半年 ROE 0.9\%（未年化），现金短债比 0.24 倍、短债覆盖不足，融资需求首先用于置换与滚动存量债务、补充营运资金，属于“补血型”而非扩张型。 公开披露的后续资金安排线索包括：股权融资线索：2026 年 4 月，公司 2026-2027 年度融资计划：银行金融机构、融资租赁公司、商业保理公司及国内外资本市场融资静态余额拟不超过 35.00 亿元。 综合看，公司的外部资金需求以补充营运资金为主、项目投入为辅，滚动偿付口径下的静态缺口约 18.16 亿元，能否落地取决于授信续作与发行窗口的配合。




\subsubsection{国联水产（300094）}
\label{co:300094}

\pfl{公司基本情况} 湛江水产品加工与预制菜龙头，2010 年 7 月在深交所创业板上市，2010 年 IPO 募资 11.50 亿元、2022 年定增募资 10 亿元；2023—2025 年连续亏损、2025 年归母净利润 -16.31 亿元，2026 年上半年扭亏为盈（归母净利润 1,003.92 万元）但资产负债率达 95.19\%，公司通过出售子公司与资产、引入益阳市国资股东等方式压降负债与补充流动性。 公司于 2010-07-08 在深交所创业板首发上市，注册资本 11.28 亿元，总股本 11.28 亿股。 截至 2026 年 9 月 24 日收盘价 2.90 元，流通市值 32.1 亿元，近一年-21.8\%，流通市值在三级行业“水产养殖”的 4 家公司中按由大到小排第\zhnumber{2}位，近一年涨跌幅低于全样本中位数（-10.1\%）11.7 个百分点。 股权结构方面，控股股东为新余国通投资管理有限公司（持股 14.00\%）。

\pfl{历史股价} 以 2021 年 1 月为起点、2026 年 9 月为终点，股价由 3.88 元变为 2.90 元，累计 -25.3\%，区间最高 6.62 元（2022 年 3 月）、最低 2.45 元（2026 年 6 月）；当前价相当于区间最高价的 43.8\%，为区间最低价的 1.18 倍。 月度收益的年化波动率 44.6\%，高于全样本中位数 37.4\%，在三级行业“水产养殖”的 4 家公司中波动率由高到低排第\zhnumber{2}位。

\begin{center}
\begin{minipage}{\textwidth}\centering
\begin{tikzpicture}
\begin{axis}[
  width=12.6cm, height=3.9cm, scale only axis,
  xmin=2020.922, xmax=2026.888, ymin=2.28, ymax=7.08,
  xtick={2021,2022,2023,2024,2025,2026},
  yearticks,
  yticklabel style={font=\scriptsize},
  ylabel={元/股}, ylabel style={font=\scriptsize},
  grid=major, grid style={LightGray, line width=0.3pt},
  axis line style={InkGray, line width=0.5pt},
  every axis plot/.append style={line width=0.9pt},
]
\addplot[AgriGreen] table[x=x,y=y]{data/clean/profiles/px_300094.dat};
\addplot[SoftRed, dashed, line width=0.7pt] table[x=x,y=y]{data/clean/profiles/pxzz_300094.dat};
\addplot[only marks, mark=*, mark size=1.1pt, SoftRed] table[x=x,y=y]{data/clean/profiles/pxzz_300094.dat};
\end{axis}
\end{tikzpicture}

\captionof{figure}[国联水产（300094）股价与周期骨架]{国联水产（300094）月度收盘价（元/股）与 ZigZag 周期骨架（阈值 25\%，圆点为周期转折点）}
\end{minipage}
\end{center}

\pfl{过去周期分析} 以 25\% 的 ZigZag 阈值划分，样本区间内股价形成 5 个主升段与 6 个主跌段。 最大主跌段为 2023 年 11 月 至 2024 年 6 月，7 个月-55.7\%；最大主升段出现在 2024 年 6 月 至 2024 年 11 月，5 个月+83.9\%。 最近一次转折出现在 2026 年 6 月（下跌段自 2024 年 11 月起，19 个月-47.5\%），当前处于该段的延续之中。 公司股价与申万农林牧渔指数几乎同步见底，个股并未走出独立行情。 水产养殖的价格周期由投苗量与消费需求共同决定（第 \ref{sec:other-livestock} 节），与猪周期的同步性弱于禽链，公司 2026 年 6 月末资产负债率 95.2\%，高于全样本中位数 52.2\%，其股价因此更多反映资金面与信用风险，而不是价格弹性本身。 综合区间形态看，该股单段涨跌幅度偏大、以趋势段为主（最大主升 +83.9\%、最大主跌 -55.7\%），年化波动率 44.6\% 与全样本中位数 37.4\% 相比波动明显大于样本中位数。

\pfl{盈利情况} 2026 年上半年营业收入 14.45 亿元（同比 -12.5\%），归母净利润 0.10 亿元，实现扭亏。 以年报为趋势对照，2023---2025 年营业收入由 46.17 亿元变为 31.90 亿元（两年年均复合增速 -16.9\%），归母净利润由 -5.32 亿元变为 -16.31 亿元。 按半年报口径，2026 年上半年的 ROE 为 10.0\%（未年化）、销售净利率 0.7\%、毛利率 11.4\%，ROE 在“水产养殖”4 家公司中按数值由高到低排第\zhnumber{2}位。 同期全样本半年 ROE 中位数 0.75\%，公司与之相差 9.26 个百分点（略好于中位数公司）。 针对上述压力，公司披露的应对举措包括：2026 年 6 月，存货规模由前期高位压降至 3.07 亿元（同比下降 80.78\%）。 综合看，公司在报告期维持盈利（高于全样本半年 ROE 中位数 0.75\%），利润的可持续性取决于水产养殖景气的延续与成本管控的执行。

\begin{center}\small\setlength{\tabcolsep}{3.4pt}
\captionof{table}[国联水产（300094）关键财务指标]{国联水产（300094）关键财务指标（合并报表口径；两个半年报期均未年化；现金短债比 $=$ 货币资金 $\div$ 短期债务）}
\begin{adjustbox}{max width=\textwidth}
\begin{tabular}{@{}lrrrrr@{}}
\toprule
指标 & 2023 年 & 2024 年 & 2025 年 & 2025 年半年报 & 2026 年半年报 \\
\midrule
营业收入（亿元） & 46.17 & 34.09 & 31.90 & 16.51 & 14.45 \\
归母净利润（亿元） & -5.32 & -7.42 & -16.31 & -5.40 & 0.10 \\
毛利率（\%） & 7.1 & 9.8 & -6.1 & -4.7 & 11.4 \\
ROE（\%） & -18.7 & -33.6 & -160.2 & -34.6 & 10.0 \\
资产负债率（\%） & 47.6 & 54.1 & 96.0 & 63.2 & 95.2 \\
经营活动现金流净额（亿元） & 1.71 & 0.29 & 1.06 & 0.98 & 1.92 \\
货币资金（亿元） & 4.82 & 3.78 & 2.58 & 3.05 & 1.98 \\
有息负债合计（亿元） & 16.79 & 14.52 & 13.52 & 14.45 & 10.84 \\
现金短债比（倍） & 0.34 & 0.29 & 0.20 & 0.23 & 0.19 \\
\bottomrule
\end{tabular}
\end{adjustbox}
\end{center}

\begin{center}
\begin{minipage}{\textwidth}\centering
\begin{tikzpicture}
\begin{groupplot}[
  group style={group size=2 by 1, horizontal sep=0.60cm},
  width=6.85cm, height=4.60cm,
  tick label style={font=\tiny},
  title style={font=\scriptsize\heiti},
  yticklabel style={font=\tiny},
  ylabel style={font=\tiny},
  grid=major, grid style={LightGray, line width=0.3pt},
  axis line style={InkGray, line width=0.5pt},
  enlarge x limits={abs=0.8},
  legend style={font=\scriptsize, at={(0.5,-0.20)}, anchor=north,
                legend columns=2, draw=none, fill=none, column sep=6pt},
]
\nextgroupplot[
  ybar, bar width=4pt,
  ymin=-22.79, ymax=56.27,
  xtick={0,1,2,3,4,5.7},
  xticklabels={2021,2022,2023,2024,2025,2026H1},
  ylabel={亿元},
]
\addplot[fill=AgriGreen!75, draw=AgriGreen, bar shift=-2.6pt] table[x=i, y=rev]{data/clean/profiles/fin_300094.dat};
\addplot[fill=SoftRed!60, draw=SoftRed, bar shift=2.6pt] table[x=i, y=ni]{data/clean/profiles/fin_300094.dat};
\addplot[fill=AgriGreen!30, draw=AgriGreen, pattern=north east lines, bar shift=-2.6pt] table[x=x, y=rev]{data/clean/profiles/finh_300094.dat};
\addplot[fill=SoftRed!20, draw=SoftRed, pattern=north east lines, bar shift=2.6pt] table[x=x, y=ni]{data/clean/profiles/finh_300094.dat};
\legend{营业收入（年/半年）, 归母净利润（年/半年）}
\nextgroupplot[
  ymin=-180.7, ymax=121.6,
  xtick={0,1,2,3,4,5.7},
  xticklabels={2021,2022,2023,2024,2025,2026H1},
  ylabel={\%},
]
\addplot[AgriGold, mark=*, mark size=1.3pt] table[x=i, y=roe]{data/clean/profiles/fin_300094.dat};
\addplot[OilBlue, mark=square*, mark size=1.3pt] table[x=i, y=debt]{data/clean/profiles/fin_300094.dat};
\addplot[only marks, AgriGold, mark=*, mark size=2.0pt] table[x=x, y=roe]{data/clean/profiles/finh_300094.dat};
\addplot[only marks, OilBlue, mark=square*, mark size=2.0pt] table[x=x, y=debt]{data/clean/profiles/finh_300094.dat};
\legend{ROE, 资产负债率}
\end{groupplot}
\end{tikzpicture}
\begin{tikzpicture}
\begin{axis}[
  name=finbot,
  width=13.75cm, height=3.10cm, scale only axis,
  ymin=-1.78, ymax=19.98,
  xtick={0,1,2,3,4,5.7},
  xticklabels={2021,2022,2023,2024,2025,2026H1},
  tick label style={font=\tiny},
  yticklabel style={font=\tiny},
  ylabel={亿元}, ylabel style={font=\tiny},
  ybar, bar width=4pt,
  grid=major, grid style={LightGray, line width=0.3pt},
  axis line style={InkGray, line width=0.5pt},
  enlarge x limits={abs=0.8},
  legend style={font=\scriptsize, at={(0.5,-0.26)}, anchor=north,
                legend columns=2, draw=none, fill=none, column sep=10pt},
]
\addplot[fill=AgriGreen!55, draw=AgriGreen, bar shift=-3.0pt] table[x=i, y=cash]{data/clean/profiles/fin_300094.dat};
\addplot[fill=SoftRed!45, draw=SoftRed, bar shift=3.0pt] table[x=i, y=ibd]{data/clean/profiles/fin_300094.dat};
\addplot[fill=AgriGreen!25, draw=AgriGreen, pattern=north east lines, bar shift=-3.0pt] table[x=x, y=cash]{data/clean/profiles/finh_300094.dat};
\addplot[fill=SoftRed!20, draw=SoftRed, pattern=north east lines, bar shift=3.0pt] table[x=x, y=ibd]{data/clean/profiles/finh_300094.dat};
\legend{货币资金（左轴）, 有息负债（左轴）}
\end{axis}
\begin{axis}[
  at={(finbot.south east)}, anchor=south east,
  width=13.75cm, height=3.10cm, scale only axis,
  axis y line*=right, axis x line*=none, xtick=\empty,
  ymin=-0.53, ymax=2.66,
  tick label style={font=\tiny},
  yticklabel style={font=\tiny}, scaled y ticks=false,
  legend style={font=\scriptsize, at={(0.5,-0.44)}, anchor=north,
                legend columns=1, draw=none, fill=none},
]
\addplot[OilBlue, mark=*, mark size=2.2pt, line width=1.3pt] table[x=i, y=ocf]{data/clean/profiles/fin_300094.dat};
\addplot[only marks, OilBlue, mark=*, mark size=3.2pt] table[x=x, y=ocf]{data/clean/profiles/finh_300094.dat};
\legend{经营活动现金流净额（右轴，亿元，蓝点）}
\end{axis}
\end{tikzpicture}

\captionof{figure}[国联水产（300094）近年主要财务指标]{国联水产（300094）近年主要财务指标（左上：营业收入与归母净利润；右上：ROE 与资产负债率；下：货币资金、有息负债为左轴柱状，经营活动现金流净额为其右轴的蓝点折线。
2021---2025 年为年报口径，2026H1 为 2026 年半年报/半年末，与其前一年报之间留出间隔、以斜纹或空心点区分；半年报数据未年化）}
\end{minipage}
\end{center}

\pfl{财务分析} 2026 年 6 月末资产负债率 95.2\%，较 2025 年末下降 0.8 个百分点（样本中位数 52.2\%，所处三级行业内按由低到高排第\zhnumber{3}位）。 2026 年 6 月末货币资金 1.98 亿元、短期债务（短期借款与一年内到期非流动负债）10.29 亿元、长期债务 0.56 亿元，有息负债合计 10.84 亿元，现金短债比 0.19 倍——缺口明显。 净有息负债 8.86 亿元，相当于其 2026 年上半年营业收入的 61\%。 流动比率 0.65、速动比率 0.47，短期偿债指标偏紧。 2026 年上半年经营活动现金流净额 1.92 亿元，经营现金流与利润同向为正，内生资金可以支撑运营。 综合看，账面货币资金对有息负债与短期债务的覆盖不足，滚动偿付对新增融资的依赖度较高；资产负债率高于样本中位数 43.0 个百分点；有息负债中短期债务占比更高，期限结构仍以滚动续作为主。 公司披露的股东与质押风险为：2026 年 5 月，控股股东关联方冠联国际投资有限公司及实际控制人李忠合计质押公司股份 5，896.68 万股；2026 年 9 月，第三方风险评级显示公司整体质押股份 2.08 亿股、占总股本 18.43\%。

\pfl{重大战略转型} 从公开披露看，公司的战略推进集中在：并购与重组（2 项）：2026 年 8 月，控股股东新余国通向益阳高新农业开发有限公司（实际控制人为益阳市国资委）协议转让 5.67\% 公司股份；融资与资本运作（1 项）：2022 年 12 月，定增募资 10.00 亿元中的 7.00 亿元投向国美水产中央厨房项目等预制菜产能与渠道建设。 公告口径上，近三年（2023 年 9 月---2026 年 9 月）公司披露重大重组与并购类公告 4 条、对外投资与转型类公告 6 条、回购与股权激励类公告 29 条，合计公告 405 条。 综合判断，报告期内公司有 10 条与战略相关的公告落地（重大重组与并购 4 条、对外投资与转型 6 条），资本运作与主业调整之间存在直接关联。

\begin{center}\small\setlength{\tabcolsep}{4pt}
\captionof{table}[国联水产（300094）历史融资与资本运作]{国联水产（300094）历史融资与资本运作（据公司公告与公开报道整理；金额为披露值，含首次公开发行、再融资、债券、银行借款与重整投资等）}
\begin{adjustbox}{max width=\textwidth}
\begin{tabular}{@{}p{1.45cm}p{2.55cm}p{4.05cm}p{4.95cm}@{}}
\toprule
时间 & 方式 & 规模 & 用途与说明 \\
\midrule
2010-06-28 & IPO（深圳证券交易所创业板） & 发行 8，\allowbreak{}000 万股（占发行后总股本 25\%），\allowbreak{}发行价 14.38 元/\allowbreak{}股 & 罗非鱼加工厂建设项目（拟使用募集资金 7，\allowbreak{}科研中心建设项目及其他与主营业务相关的营运资金；股票于 2010 年 7 月 8 日挂牌 \\
2022-12-12 & 向特定对象发行股票（定增） & 向 15 名特定对象发行 2.21 亿股，\allowbreak{}发行价 4.52 元/\allowbreak{}股 & 不超 10.00 亿元募资中 7.00 亿元投向全资子公司国美水产中央厨房项目、\allowbreak{}其余补充流动资金（后两个募投项目均已终止 \\
2025-12-09 & 出售全资子公司股权（关联交易、承债式收购） & 向控股股东新余国通投资管理有限公司转让全资子公司广东新盈食品科技有限公司 100\% 股权 & 盘活存量资产、\allowbreak{}将部分上游闲置资产及存货转为现金资产、\allowbreak{}优化资产结构与存货周转效率；经 2025 年 12 月 8 日第六届董事会第十七次会议审议通过（公告编号：2025-062） \\
2026-05-29 & 出售资产（孙公司资产打包出售） & 以人民币 1，\allowbreak{}配套设施一次性打包出售给湖北鸿展仓储服务有限公司 & 优化闲置资产、\allowbreak{}快速回笼资金改善现金流；经 2026 年 5 月 29 日第六届董事会第二十三次会议审议通过 \\
2026-07-15 & 出售境外房产资产 & 全资子公司 Liancheng Investments，\allowbreak{}700.00 万美元（折合人民币 1.15 亿元）向 LBA-Logistics 出售位于美国加利福尼亚州尤尼恩城法伯街 2910 号的房地产 & 交易所得款项用于公司业务发展；经 2026 年 7 月 15 日第六届董事会第二十四次会议审议通过 \\
\bottomrule
\end{tabular}
\end{adjustbox}
\end{center}

\pfl{历史融投资情况} 从现金流量表看，2025 年公司取得借款收到的现金 22.27 亿元、偿还债务支付的现金 23.79 亿元，筹资活动现金流净额 -2.76 亿元（2023 年为取得借款 22.56 亿元、偿还债务 23.45 亿元）。 2026 年上半年筹资活动现金流净额为 -1.99 亿元（取得借款 4.63 亿元、偿还债务 6.50 亿元，未年化），已转为净流出，处于净还债状态。 2025 年分配股利、利润或偿付利息支付的现金合计 0.53 亿元。 公开披露的融资记录共 5 笔，工具类型以 IPO、定向增发为主；金额与用途见下表。 其中规模最大的两笔为：2010 年 6 月，IPO，发行 8，000 万股（占发行后总股本 25\%），发行价 14.38 元/股；2026 年 7 月，出售境外房产资产，全资子公司 Liancheng Investments，700.00 万美元（折合人民币 1.15 亿元）向 LBA-Logistics 出售位于美国加利福尼亚州尤尼恩城法伯街 2910 号的房地产。 股本方面，近三年披露股本变动 5 次（限售股份上市 1 次、股份回购 1 次）；总股本由 2021 年末的 9.12 亿股变为 11.28 亿股（+23.66\%）。 分红方面，近三年未实施现金分红，在全部样本中属于分红能力偏弱的一端。 近三年“再融资”类公告 1 条，涉及定向增发。 综合看，公司融资结构上股权与债务融资并举；2025 年筹资活动现金流净额 -2.76 亿元，融资节奏仍以维持存量债务周转为主，与前述经营与投资现金流的方向基本一致。

\pfl{未来潜在融资需求分析} 2026 年 6 月末货币资金 1.98 亿元，而短期债务 10.29 亿元（现金短债比 0.19 倍），2026 年上半年经营现金流净额 1.92 亿元。按“短期债务减去账面货币资金”的滚动偿付口径，静态缺口约 8.30 亿元；上半年盈利或接近盈亏平衡，该缺口不随经营亏损进一步扩大。 公司 2026 年 6 月末资产负债率 95.2\%、2026 年上半年 ROE 10.0\%（未年化），资产负债率偏高（95.2\%），融资需求首先用于置换与滚动存量债务、补充营运资金，属于“补血型”而非扩张型。 截至目前公司仍处于预重整/重整程序中，重整投资人的招募与投资款到位是未来一年最主要的外部资金来源。 综合看，公司的外部资金需求以补充流动性与项目投入为主，滚动偿付口径下的静态缺口约 8.30 亿元，融资的时点与规模更多取决于信用条件与资本市场窗口。




\subsubsection{*ST景谷（600265）}
\label{co:600265}

\pfl{公司基本情况} 云南景谷的林业与木材加工上市壳资源型公司（云南景谷林业股份有限公司），2000 年 8 月在上交所上市，2008 年后多次易主、三度被实施退市风险警示；2018 年 8 月周大福投资入主（最终控制方郑家纯），2023 年现金 2.703 亿元收购汇银木业 51\% 后业绩暴雷、2025 年以 13,336.60 万元将汇银木业 51\% 股权回售给控股股东，2025 年 10 月起无偿受赠算力服务公司博达数科 100\% 股权；2025 年净利润 -2.46 亿元且扣除后营业收入低于 3 亿元，2026 年 4 月 27 日起简称变更为“*ST 景谷”。 公司于 2000-08-25 在上交所首发上市，注册资本 1.30 亿元，总股本 1.30 亿股。 截至 2026 年 9 月 24 日收盘价 23.41 元，流通市值 30.4 亿元，近一年-7.9\%，流通市值在三级行业“林业Ⅲ”的 4 家公司中按由大到小排第\zhnumber{2}位，近一年涨跌幅高于全样本中位数（-10.1\%）2.1 个百分点。 按 2025 年年报披露的行业构成，收入占比前两位为制造业行业 71.7\%；其他(补充) 28.3\%。 与 2021 年年报相比，第一大行业口径由“林业行业”（94.4\%）变为“制造业行业”（71.7\%）；两期披露的分类口径有调整，占比差异中同时包含分类因素。

\pfl{历史股价} 自 2021 年 1 月至 2026 年 9 月，股价由 17.27 元变为 23.41 元，累计 +35.6\%，区间最高 24.45 元（2026 年 5 月）、最低 13.30 元（2022 年 6 月）；当前价相当于区间最高价的 95.7\%，为区间最低价的 1.76 倍。 月度收益的年化波动率 29.5\%，低于全样本中位数 37.4\%，在三级行业“林业Ⅲ”的 4 家公司中波动率由高到低排第\zhnumber{4}位。 把股价阶段与公司事件对齐看，2024 年 8 月 至 2026 年 5 月股价上涨+81.1\%，同期（2026 年 5 月）提示除已披露事项外不存在筹划重大资产重组、股份发行、破产重整、引进战略投资者等应披露而未披露的重大事项。 把这些阶段与公司事件对齐后可以看到，几轮大涨大跌多由公司自身的资本运作与风险事件触发，行业景气只解释了其中一部分。

\begin{center}
\begin{minipage}{\textwidth}\centering
\begin{tikzpicture}
\begin{axis}[
  width=12.6cm, height=3.9cm, scale only axis,
  xmin=2020.922, xmax=2026.888, ymin=12.37, ymax=26.16,
  xtick={2021,2022,2023,2024,2025,2026},
  yearticks,
  yticklabel style={font=\scriptsize},
  ylabel={元/股}, ylabel style={font=\scriptsize},
  grid=major, grid style={LightGray, line width=0.3pt},
  axis line style={InkGray, line width=0.5pt},
  every axis plot/.append style={line width=0.9pt},
]
\addplot[AgriGreen] table[x=x,y=y]{data/clean/profiles/px_600265.dat};
\addplot[SoftRed, dashed, line width=0.7pt] table[x=x,y=y]{data/clean/profiles/pxzz_600265.dat};
\addplot[only marks, mark=*, mark size=1.1pt, SoftRed] table[x=x,y=y]{data/clean/profiles/pxzz_600265.dat};
\end{axis}
\end{tikzpicture}

\captionof{figure}[*ST景谷（600265）股价与周期骨架]{*ST景谷（600265）月度收盘价（元/股）与 ZigZag 周期骨架（阈值 25\%，圆点为周期转折点）}
\end{minipage}
\end{center}

\pfl{过去周期分析} 以 25\% 的 ZigZag 阈值划分，样本区间内股价形成 3 个主升段与 3 个主跌段。 最大主升段出现在 2024 年 8 月 至 2026 年 5 月，21 个月+81.1\%；最大主跌段为 2021 年 8 月 至 2022 年 6 月，10 个月-43.9\%。 最近一次转折出现在 2026 年 5 月（上涨段自 2024 年 8 月起，21 个月+81.1\%），当前处于该段之后的震荡之中。 公司股价较申万农林牧渔指数提前约 47 个月见底，是板块中较早定价周期反转的公司之一。 林业的现金流依赖林木采伐与政策补贴，周期长、变现慢，公司 2026 年 6 月末资产负债率 26.7\%、半年 ROE -0.0\%（未年化），在本轮下行中属于随行业同步调整的一类。 综合区间形态看，该股单段涨跌幅度偏大、以趋势段为主（最大主升 +81.1\%、最大主跌 -43.9\%），年化波动率 29.5\% 与全样本中位数 37.4\% 相比波动小于样本中位数。

\pfl{盈利情况} 2026 年上半年营业收入 0.56 亿元（同比 -54.6\%），归母净利润 -0.00 亿元，亏损同比收窄 1.24 亿元。 以年报为趋势对照，2023---2025 年营业收入由 5.90 亿元变为 1.95 亿元（两年年均复合增速 -42.5\%），归母净利润由 0.06 亿元变为 -2.46 亿元。 按半年报口径，2026 年上半年的 ROE 为 -0.0\%（未年化）、销售净利率 5.4\%、毛利率 32.9\%，ROE 在“林业Ⅲ”4 家公司中按数值由高到低排第\zhnumber{1}位。 同期全样本半年 ROE 中位数 0.75\%，公司与之相差 0.75 个百分点（弱于中位数公司）。 综合看，公司仍处亏损区间、亏损同比收窄，半年 ROE 在三级行业内按由高到低排第\zhnumber{1}位，单位盈利能力的修复依赖行业景气与自身成本端的改善，报表弹性主要体现为价格与销量的方向性变化。

\begin{center}\small\setlength{\tabcolsep}{3.4pt}
\captionof{table}[*ST景谷（600265）关键财务指标]{*ST景谷（600265）关键财务指标（合并报表口径；两个半年报期均未年化；现金短债比 $=$ 货币资金 $\div$ 短期债务）}
\begin{adjustbox}{max width=\textwidth}
\begin{tabular}{@{}lrrrrr@{}}
\toprule
指标 & 2023 年 & 2024 年 & 2025 年 & 2025 年半年报 & 2026 年半年报 \\
\midrule
营业收入（亿元） & 5.90 & 4.47 & 1.95 & 1.23 & 0.56 \\
归母净利润（亿元） & 0.06 & -0.73 & -2.46 & -1.24 & -0.00 \\
毛利率（\%） & 10.3 & 2.3 & 0.1 & -11.9 & 32.9 \\
ROE（\%） & 4.0 & -56.1 & -4.7 & -1.3 & -0.0 \\
资产负债率（\%） & 63.4 & 73.5 & 29.9 & 77.7 & 26.7 \\
经营活动现金流净额（亿元） & 0.92 & 0.03 & -0.26 & -0.08 & 0.29 \\
货币资金（亿元） & 0.59 & 0.40 & 0.71 & 0.25 & 0.92 \\
有息负债合计（亿元） & 2.03 & 2.28 & 0.05 & 2.25 & 0.04 \\
现金短债比（倍） & 2.18 & 0.53 & 34.26 & 0.11 & 55.11 \\
\bottomrule
\end{tabular}
\end{adjustbox}
\end{center}

\begin{center}
\begin{minipage}{\textwidth}\centering
\begin{tikzpicture}
\begin{groupplot}[
  group style={group size=2 by 1, horizontal sep=0.60cm},
  width=6.85cm, height=4.60cm,
  tick label style={font=\tiny},
  title style={font=\scriptsize\heiti},
  yticklabel style={font=\tiny},
  ylabel style={font=\tiny},
  grid=major, grid style={LightGray, line width=0.3pt},
  axis line style={InkGray, line width=0.5pt},
  enlarge x limits={abs=0.8},
  legend style={font=\scriptsize, at={(0.5,-0.20)}, anchor=north,
                legend columns=2, draw=none, fill=none, column sep=6pt},
]
\nextgroupplot[
  ybar, bar width=4pt,
  ymin=-3.29, ymax=6.90,
  xtick={0,1,2,3,4,5.7},
  xticklabels={2021,2022,2023,2024,2025,2026H1},
  ylabel={亿元},
]
\addplot[fill=AgriGreen!75, draw=AgriGreen, bar shift=-2.6pt] table[x=i, y=rev]{data/clean/profiles/fin_600265.dat};
\addplot[fill=SoftRed!60, draw=SoftRed, bar shift=2.6pt] table[x=i, y=ni]{data/clean/profiles/fin_600265.dat};
\addplot[fill=AgriGreen!30, draw=AgriGreen, pattern=north east lines, bar shift=-2.6pt] table[x=x, y=rev]{data/clean/profiles/finh_600265.dat};
\addplot[fill=SoftRed!20, draw=SoftRed, pattern=north east lines, bar shift=2.6pt] table[x=x, y=ni]{data/clean/profiles/finh_600265.dat};
\legend{营业收入（年/半年）, 归母净利润（年/半年）}
\nextgroupplot[
  ymin=-113.3, ymax=90.8,
  xtick={0,1,2,3,4,5.7},
  xticklabels={2021,2022,2023,2024,2025,2026H1},
  ylabel={\%},
]
\addplot[AgriGold, mark=*, mark size=1.3pt] table[x=i, y=roe]{data/clean/profiles/fin_600265.dat};
\addplot[OilBlue, mark=square*, mark size=1.3pt] table[x=i, y=debt]{data/clean/profiles/fin_600265.dat};
\addplot[only marks, AgriGold, mark=*, mark size=2.0pt] table[x=x, y=roe]{data/clean/profiles/finh_600265.dat};
\addplot[only marks, OilBlue, mark=square*, mark size=2.0pt] table[x=x, y=debt]{data/clean/profiles/finh_600265.dat};
\legend{ROE, 资产负债率}
\end{groupplot}
\end{tikzpicture}
\begin{tikzpicture}
\begin{axis}[
  name=finbot,
  width=13.75cm, height=3.10cm, scale only axis,
  ymin=-0.23, ymax=2.55,
  xtick={0,1,2,3,4,5.7},
  xticklabels={2021,2022,2023,2024,2025,2026H1},
  tick label style={font=\tiny},
  yticklabel style={font=\tiny},
  ylabel={亿元}, ylabel style={font=\tiny},
  ybar, bar width=4pt,
  grid=major, grid style={LightGray, line width=0.3pt},
  axis line style={InkGray, line width=0.5pt},
  enlarge x limits={abs=0.8},
  legend style={font=\scriptsize, at={(0.5,-0.26)}, anchor=north,
                legend columns=2, draw=none, fill=none, column sep=10pt},
]
\addplot[fill=AgriGreen!55, draw=AgriGreen, bar shift=-3.0pt] table[x=i, y=cash]{data/clean/profiles/fin_600265.dat};
\addplot[fill=SoftRed!45, draw=SoftRed, bar shift=3.0pt] table[x=i, y=ibd]{data/clean/profiles/fin_600265.dat};
\addplot[fill=AgriGreen!25, draw=AgriGreen, pattern=north east lines, bar shift=-3.0pt] table[x=x, y=cash]{data/clean/profiles/finh_600265.dat};
\addplot[fill=SoftRed!20, draw=SoftRed, pattern=north east lines, bar shift=3.0pt] table[x=x, y=ibd]{data/clean/profiles/finh_600265.dat};
\legend{货币资金（左轴）, 有息负债（左轴）}
\end{axis}
\begin{axis}[
  at={(finbot.south east)}, anchor=south east,
  width=13.75cm, height=3.10cm, scale only axis,
  axis y line*=right, axis x line*=none, xtick=\empty,
  ymin=-0.86, ymax=1.35,
  tick label style={font=\tiny},
  yticklabel style={font=\tiny}, scaled y ticks=false,
  legend style={font=\scriptsize, at={(0.5,-0.44)}, anchor=north,
                legend columns=1, draw=none, fill=none},
]
\addplot[OilBlue, mark=*, mark size=2.2pt, line width=1.3pt] table[x=i, y=ocf]{data/clean/profiles/fin_600265.dat};
\addplot[only marks, OilBlue, mark=*, mark size=3.2pt] table[x=x, y=ocf]{data/clean/profiles/finh_600265.dat};
\legend{经营活动现金流净额（右轴，亿元，蓝点）}
\end{axis}
\end{tikzpicture}

\captionof{figure}[*ST景谷（600265）近年主要财务指标]{*ST景谷（600265）近年主要财务指标（左上：营业收入与归母净利润；右上：ROE 与资产负债率；下：货币资金、有息负债为左轴柱状，经营活动现金流净额为其右轴的蓝点折线。
2021---2025 年为年报口径，2026H1 为 2026 年半年报/半年末，与其前一年报之间留出间隔、以斜纹或空心点区分；半年报数据未年化）}
\end{minipage}
\end{center}

\pfl{财务分析} 2026 年 6 月末资产负债率 26.7\%，较 2025 年末下降 3.2 个百分点（样本中位数 52.2\%，所处三级行业内按由低到高排第\zhnumber{1}位）。 2026 年 6 月末货币资金 0.92 亿元、短期债务（短期借款与一年内到期非流动负债）0.02 亿元、长期债务 0.02 亿元，有息负债合计 0.04 亿元，现金短债比 55.11 倍——覆盖充分。 净有息负债 -0.87 亿元，相当于其 2026 年上半年营业收入的 -156\%。 流动比率 3.58、速动比率 1.72，短期流动性无明显缺口。 2026 年上半年经营活动现金流净额 0.29 亿元，净利润为负而经营现金流为正，差额主要来自折旧摊销与营运资本变动，说明亏损尚未完全转化为现金流出。 综合看，现金短债比 55.11 倍，短期偿付压力可控；资产负债率低于样本中位数 25.5 个百分点；有息负债中长期债务占比更高，期限结构对短债滚动的压力较小。 公司披露的退市与持续经营风险为：2026 年 4 月，公司股票被实施退市风险警示；2025 年 8 月，2024 年营业收入占比达 87.02\% 本次出售完成后主营业务规模将急剧下降。

\pfl{重大战略转型} 从公开披露看，公司的战略推进集中在：并购与重组（4 项）：2025 年 12 月，公司完成向控股股东出售汇银木业 51\% 股权的重大资产出售；2023 年 2 月，公司以现金 2.70 亿元收购上述项目权；其他动作（1 项）：2026 年 4 月，控股股东无偿赠与上海博达数智科技有限公司 51\% 股权（2026 年 4 月再受赠剩余 49\%）。 公告口径上，近三年（2023 年 9 月---2026 年 9 月）公司披露重大重组与并购类公告 23 条、对外投资与转型类公告 1 条、回购与股权激励类公告 0 条，合计公告 478 条。 主营结构的口径变化已经落到报表上：2021 年年报的第一大行业为“林业行业”（94.4\%），2025 年年报为“制造业行业”（71.7\%）；公司在这两期的分部划分方式有过调整。 综合判断，报告期内公司有 24 条与战略相关的公告落地（重大重组与并购 23 条、对外投资与转型 1 条），资本运作与主业调整之间存在直接关联。

\begin{center}\small\setlength{\tabcolsep}{4pt}
\captionof{table}[*ST景谷（600265）历史融资与资本运作]{*ST景谷（600265）历史融资与资本运作（据公司公告与公开报道整理；金额为披露值，含首次公开发行、再融资、债券、银行借款与重整投资等）}
\begin{adjustbox}{max width=\textwidth}
\begin{tabular}{@{}p{1.45cm}p{2.55cm}p{4.05cm}p{4.95cm}@{}}
\toprule
时间 & 方式 & 规模 & 用途与说明 \\
\midrule
2000-07-21 & IPO（上海证券交易所） & 发行 4,\allowbreak{}000 万股，\allowbreak{}发行价 5.19 元/\allowbreak{}股，\allowbreak{}募集资金总额 1.96 亿元 & 经证监发行字[2000]99 号文核准，\allowbreak{}发行后注册资本 1.05 亿元 \\
2023-02 & 现金收购重大资产（并购） & 以支付现金方式购买崔会军、\allowbreak{}王兰存、\allowbreak{}京保基金、\allowbreak{}技改基金合计持有的唐县汇银木业有限公司 51\% 股权 & 新增人造板制造业务、\allowbreak{}改善上市公司盈利能力；评估机构亚超评估采用收益法评估结论（北京亚超评报字（2023）第 A003 号），\allowbreak{}增值率 83.7\% \\
2025-08-16 & 重大资产出售暨关联交易（向控股股东出售子公司股权） & 以现金交易方式向控股股东周大福投资有限公司转让持有的汇银木业 51\% 股权 & 剥离严重亏损且面临诉讼风险的子公司、\allowbreak{}化解风险、\allowbreak{}为业务转型打基础；2025 年 8 月 16 日披露筹划出售资产涉及重大资产重组（公告编号：2025-066） \\
2025-10 & 无偿受赠资产（控股股东资产赠与） & 控股股东周大福投资有限公司将其持有的上海博达数智科技有限公司 51\% 股权无偿赠与公司 & 提高公司持续经营能力、\allowbreak{}优化业务结构，\allowbreak{}注入算力服务业务；2026 年 4 月控股股东将博达数科剩余 49\% 股权无偿赠与公司 \\
2025-12-05 & 关联方资金占用清收（重大资产出售配套安排） & 截至上交所问询函回复出具日，\allowbreak{}标的公司汇银木业应偿还上市公司拆借资金的本金及利息为 6，\allowbreak{}172.22 万元 & 明确标的公司对上市公司借款的本息偿付安排；公司回复上海证券交易所云南景谷林业股份有限公司重大资产重组草案的信息披露的问询函 \\
\bottomrule
\end{tabular}
\end{adjustbox}
\end{center}

\pfl{历史融投资情况} 从现金流量表看，2025 年公司取得借款收到的现金 0.23 亿元、偿还债务支付的现金 0.70 亿元、吸收权益性投资 1.95 亿元，筹资活动现金流净额 0.39 亿元（2023 年为取得借款 2.00 亿元、偿还债务 0.06 亿元）。 2026 年上半年筹资活动现金流净额为 -0.04 亿元（取得借款 — 亿元、偿还债务 — 亿元，未年化），已转为净流出，处于净还债状态。 2025 年分配股利、利润或偿付利息支付的现金合计 0.17 亿元。 公开披露的融资记录共 5 笔，工具类型以 IPO 为主；金额与用途见下表。 其中规模最大的两笔为：2000 年 7 月，IPO，发行 4,000 万股，发行价 5.19 元/股，募集资金总额 1.96 亿元；2025 年 12 月，关联方资金占用清收，截至上交所问询函回复出具日，标的公司汇银木业应偿还上市公司拆借资金的本金及利息为 6，172.22 万元。 股本方面，近三年披露股本变动 3 次；总股本自 2021 年末以来未发生变化（1.30 亿股）。 分红方面，近三年未实施现金分红，在全部样本中属于分红能力偏弱的一端。 近三年“再融资”类公告 16 条，涉及定向增发、募集资金使用。 综合看，公司融资结构上股权与债务融资并举；2025 年筹资活动现金流净额 0.39 亿元，年度内仍为净融入，与前述经营与投资现金流的方向基本一致。

\pfl{未来潜在融资需求分析} 2026 年 6 月末货币资金 0.92 亿元，而短期债务 0.02 亿元（现金短债比 55.11 倍），2026 年上半年经营现金流净额 0.29 亿元。账面货币资金可以覆盖短期债务，缺口不体现在静态偿付层面。 按第 \ref{sec:financing-need} 节的三档划分，公司属于第二档“保壳型”：近三年重大重组与并购类公告 23 条，2026 年上半年归母净利润 -0.00 亿元，融资需求更多体现为“资产置换 + 维持上市地位”，而不是扩产。 公开披露的后续资金安排线索包括：资本开支安排：2026 年半年报，提出“立足林业主业精益管理+完成博达数科剩余 49\% 股权受赠”的工作主线。 综合看，公司在静态偿付层面不存在缺口，外部融资更多是配合扩张节奏与并购整合的主动性安排，而非偿债压力驱动。




\subsubsection{福建金森（002679）}
\label{co:002679}

\pfl{公司基本情况} 福建三明将乐县的国有控股林业上市公司（福建金森林业股份有限公司），2012 年 6 月在深交所中小板上市，是全国首家纯森林资源培育型国有控股上市公司、南方集体林区 FSC 国际森林认证用材林蓄积量最大的企业；主营森林经营与木材采运，2025 年营业收入 1.487 亿元、归母净利润 923.06 万元，2026 年上半年受季节性因素影响归母净利润 -1,358.18 万元（同比减亏 30.03\%），近年无新增股权融资、以存货抵押贷款为主要融资方式。 公司于 2012-06-05 在深交所主板首发上市，注册资本 2.36 亿元，总股本 2.36 亿股。 截至 2026 年 9 月 24 日收盘价 11.58 元，流通市值 27.3 亿元，近一年+9.4\%，流通市值在三级行业“林业Ⅲ”的 4 家公司中按由大到小排第\zhnumber{3}位，近一年涨跌幅高于全样本中位数（-10.1\%）19.5 个百分点。 按 2026 年半年报披露的行业构成，收入占比前三位为林业 74.4\%；生物质燃料 19.2\%；木材加工 4.3\%。 股权结构方面，控股股东为福建金森集团有限公司（持股 63.63\%）；实际控制人为将乐县财政局。

\pfl{历史股价} 以 2021 年 1 月为起点、2026 年 9 月为终点，股价由 9.90 元变为 11.58 元，累计 +17.0\%，区间最高 24.22 元（2021 年 5 月）、最低 6.67 元（2024 年 8 月）；当前价相当于区间最高价的 47.8\%，为区间最低价的 1.74 倍。 月度收益的年化波动率 55.8\%，高于全样本中位数 37.4\%，在三级行业“林业Ⅲ”的 4 家公司中波动率由高到低排第\zhnumber{2}位。 把股价阶段与公司事件对齐看，2024 年 8 月 至 2025 年 10 月股价上涨+109.9\%，同期（2025 年 6 月）同一报告期受限资产还包括 4，250 元 ETC 保证金冻结货币资金。 从时间对应关系看，公司层面的资本运作与股权、风险事件解释了区间内多数急涨急跌，与行业指数的同步性相对有限。

\begin{center}
\begin{minipage}{\textwidth}\centering
\begin{tikzpicture}
\begin{axis}[
  width=12.6cm, height=3.9cm, scale only axis,
  xmin=2020.922, xmax=2026.888, ymin=6.20, ymax=25.92,
  xtick={2021,2022,2023,2024,2025,2026},
  yearticks,
  yticklabel style={font=\scriptsize},
  ylabel={元/股}, ylabel style={font=\scriptsize},
  grid=major, grid style={LightGray, line width=0.3pt},
  axis line style={InkGray, line width=0.5pt},
  every axis plot/.append style={line width=0.9pt},
]
\addplot[AgriGreen] table[x=x,y=y]{data/clean/profiles/px_002679.dat};
\addplot[SoftRed, dashed, line width=0.7pt] table[x=x,y=y]{data/clean/profiles/pxzz_002679.dat};
\addplot[only marks, mark=*, mark size=1.1pt, SoftRed] table[x=x,y=y]{data/clean/profiles/pxzz_002679.dat};
\end{axis}
\end{tikzpicture}

\captionof{figure}[福建金森（002679）股价与周期骨架]{福建金森（002679）月度收盘价（元/股）与 ZigZag 周期骨架（阈值 25\%，圆点为周期转折点）}
\end{minipage}
\end{center}

\pfl{过去周期分析} 以 25\% 的 ZigZag 阈值划分，样本区间内股价形成 4 个主升段与 3 个主跌段。 最大主跌段为 2021 年 5 月 至 2022 年 4 月，11 个月-62.3\%；最大主升段出现在 2021 年 1 月 至 2021 年 5 月，4 个月+144.6\%。 最近一次转折出现在 2026 年 9 月（上涨段自 2026 年 6 月起，3 个月+47.0\%），当前处于该段之后的震荡之中。 公司股价较申万农林牧渔指数提前约 21 个月见底，是板块中较早定价周期反转的公司之一。 林业的现金流依赖林木采伐与政策补贴，周期长、变现慢，公司 2026 年 6 月末资产负债率 61.9\%、半年 ROE -1.8\%（未年化），在本轮下行中属于随行业同步调整的一类。 综合区间形态看，该股单段涨跌幅度偏大、以趋势段为主（最大主升 +144.6\%、最大主跌 -62.3\%），年化波动率 55.8\% 与全样本中位数 37.4\% 相比波动明显大于样本中位数。

\pfl{盈利情况} 2026 年上半年营业收入 0.45 亿元（同比 -9.5\%），归母净利润 -0.14 亿元，亏损同比收窄 0.06 亿元。 以年报为趋势对照，2023---2025 年营业收入由 1.48 亿元变为 1.49 亿元（两年年均复合增速 +0.4\%），归母净利润由 0.08 亿元变为 0.09 亿元。 按半年报口径，2026 年上半年的 ROE 为 -1.8\%（未年化）、销售净利率 -29.5\%、毛利率 40.5\%，ROE 在“林业Ⅲ”4 家公司中按数值由高到低排第\zhnumber{2}位。 同期全样本半年 ROE 中位数 0.75\%，公司与之相差 2.55 个百分点（弱于中位数公司）。 综合看，公司仍处亏损区间、亏损同比收窄，半年 ROE 在三级行业内按由高到低排第\zhnumber{2}位，单位盈利能力的修复依赖行业景气与自身成本端的改善，报表弹性主要体现为价格与销量的方向性变化。

\begin{center}\small\setlength{\tabcolsep}{3.4pt}
\captionof{table}[福建金森（002679）关键财务指标]{福建金森（002679）关键财务指标（合并报表口径；两个半年报期均未年化；现金短债比 $=$ 货币资金 $\div$ 短期债务）}
\begin{adjustbox}{max width=\textwidth}
\begin{tabular}{@{}lrrrrr@{}}
\toprule
指标 & 2023 年 & 2024 年 & 2025 年 & 2025 年半年报 & 2026 年半年报 \\
\midrule
营业收入（亿元） & 1.48 & 1.51 & 1.49 & 0.50 & 0.45 \\
归母净利润（亿元） & 0.08 & 0.10 & 0.09 & -0.19 & -0.14 \\
毛利率（\%） & 52.7 & 57.1 & 38.5 & 38.5 & 40.5 \\
ROE（\%） & 1.1 & 1.4 & 1.2 & -2.6 & -1.8 \\
资产负债率（\%） & 62.8 & 61.2 & 61.7 & 62.6 & 61.9 \\
经营活动现金流净额（亿元） & 0.70 & 0.77 & 0.71 & -0.33 & -0.46 \\
货币资金（亿元） & 2.55 & 1.83 & 1.83 & 1.53 & 1.34 \\
有息负债合计（亿元） & 12.12 & 11.31 & 11.32 & 11.66 & 11.56 \\
现金短债比（倍） & 1.49 & 1.04 & 1.04 & 0.90 & 0.77 \\
\bottomrule
\end{tabular}
\end{adjustbox}
\end{center}

\begin{center}
\begin{minipage}{\textwidth}\centering
\begin{tikzpicture}
\begin{groupplot}[
  group style={group size=2 by 1, horizontal sep=0.60cm},
  width=6.85cm, height=4.60cm,
  tick label style={font=\tiny},
  title style={font=\scriptsize\heiti},
  yticklabel style={font=\tiny},
  ylabel style={font=\tiny},
  grid=major, grid style={LightGray, line width=0.3pt},
  axis line style={InkGray, line width=0.5pt},
  enlarge x limits={abs=0.8},
  legend style={font=\scriptsize, at={(0.5,-0.20)}, anchor=north,
                legend columns=2, draw=none, fill=none, column sep=6pt},
]
\nextgroupplot[
  ybar, bar width=4pt,
  ymin=-0.34, ymax=2.14,
  xtick={0,1,2,3,4,5.7},
  xticklabels={2021,2022,2023,2024,2025,2026H1},
  ylabel={亿元},
]
\addplot[fill=AgriGreen!75, draw=AgriGreen, bar shift=-2.6pt] table[x=i, y=rev]{data/clean/profiles/fin_002679.dat};
\addplot[fill=SoftRed!60, draw=SoftRed, bar shift=2.6pt] table[x=i, y=ni]{data/clean/profiles/fin_002679.dat};
\addplot[fill=AgriGreen!30, draw=AgriGreen, pattern=north east lines, bar shift=-2.6pt] table[x=x, y=rev]{data/clean/profiles/finh_002679.dat};
\addplot[fill=SoftRed!20, draw=SoftRed, pattern=north east lines, bar shift=2.6pt] table[x=x, y=ni]{data/clean/profiles/finh_002679.dat};
\legend{营业收入（年/半年）, 归母净利润（年/半年）}
\nextgroupplot[
  ymin=-7.0, ymax=69.3,
  xtick={0,1,2,3,4,5.7},
  xticklabels={2021,2022,2023,2024,2025,2026H1},
  ylabel={\%},
]
\addplot[AgriGold, mark=*, mark size=1.3pt] table[x=i, y=roe]{data/clean/profiles/fin_002679.dat};
\addplot[OilBlue, mark=square*, mark size=1.3pt] table[x=i, y=debt]{data/clean/profiles/fin_002679.dat};
\addplot[only marks, AgriGold, mark=*, mark size=2.0pt] table[x=x, y=roe]{data/clean/profiles/finh_002679.dat};
\addplot[only marks, OilBlue, mark=square*, mark size=2.0pt] table[x=x, y=debt]{data/clean/profiles/finh_002679.dat};
\legend{ROE, 资产负债率}
\end{groupplot}
\end{tikzpicture}
\begin{tikzpicture}
\begin{axis}[
  name=finbot,
  width=13.75cm, height=3.10cm, scale only axis,
  ymin=-1.21, ymax=13.58,
  xtick={0,1,2,3,4,5.7},
  xticklabels={2021,2022,2023,2024,2025,2026H1},
  tick label style={font=\tiny},
  yticklabel style={font=\tiny},
  ylabel={亿元}, ylabel style={font=\tiny},
  ybar, bar width=4pt,
  grid=major, grid style={LightGray, line width=0.3pt},
  axis line style={InkGray, line width=0.5pt},
  enlarge x limits={abs=0.8},
  legend style={font=\scriptsize, at={(0.5,-0.26)}, anchor=north,
                legend columns=2, draw=none, fill=none, column sep=10pt},
]
\addplot[fill=AgriGreen!55, draw=AgriGreen, bar shift=-3.0pt] table[x=i, y=cash]{data/clean/profiles/fin_002679.dat};
\addplot[fill=SoftRed!45, draw=SoftRed, bar shift=3.0pt] table[x=i, y=ibd]{data/clean/profiles/fin_002679.dat};
\addplot[fill=AgriGreen!25, draw=AgriGreen, pattern=north east lines, bar shift=-3.0pt] table[x=x, y=cash]{data/clean/profiles/finh_002679.dat};
\addplot[fill=SoftRed!20, draw=SoftRed, pattern=north east lines, bar shift=3.0pt] table[x=x, y=ibd]{data/clean/profiles/finh_002679.dat};
\legend{货币资金（左轴）, 有息负债（左轴）}
\end{axis}
\begin{axis}[
  at={(finbot.south east)}, anchor=south east,
  width=13.75cm, height=3.10cm, scale only axis,
  axis y line*=right, axis x line*=none, xtick=\empty,
  ymin=-0.80, ymax=1.22,
  tick label style={font=\tiny},
  yticklabel style={font=\tiny}, scaled y ticks=false,
  legend style={font=\scriptsize, at={(0.5,-0.44)}, anchor=north,
                legend columns=1, draw=none, fill=none},
]
\addplot[OilBlue, mark=*, mark size=2.2pt, line width=1.3pt] table[x=i, y=ocf]{data/clean/profiles/fin_002679.dat};
\addplot[only marks, OilBlue, mark=*, mark size=3.2pt] table[x=x, y=ocf]{data/clean/profiles/finh_002679.dat};
\legend{经营活动现金流净额（右轴，亿元，蓝点）}
\end{axis}
\end{tikzpicture}

\captionof{figure}[福建金森（002679）近年主要财务指标]{福建金森（002679）近年主要财务指标（左上：营业收入与归母净利润；右上：ROE 与资产负债率；下：货币资金、有息负债为左轴柱状，经营活动现金流净额为其右轴的蓝点折线。
2021---2025 年为年报口径，2026H1 为 2026 年半年报/半年末，与其前一年报之间留出间隔、以斜纹或空心点区分；半年报数据未年化）}
\end{minipage}
\end{center}

\pfl{财务分析} 2026 年 6 月末资产负债率 61.9\%，较 2025 年末上升 0.2 个百分点（样本中位数 52.2\%，所处三级行业内按由低到高排第\zhnumber{4}位）。 2026 年 6 月末货币资金 1.34 亿元、短期债务（短期借款与一年内到期非流动负债）1.75 亿元、长期债务 9.81 亿元，有息负债合计 11.56 亿元，现金短债比 0.77 倍——覆盖偏紧。 净有息负债 10.22 亿元，相当于其 2026 年上半年营业收入的 2,247\%。 流动比率 7.63、速动比率 1.15，短期流动性无明显缺口。 2026 年上半年经营活动现金流净额 -0.46 亿元，经营现金流与利润同时为负，现金消耗较快。 综合看，现金短债比 0.77 倍，短期资金安排偏紧；资产负债率高于样本中位数 9.7 个百分点；有息负债中长期债务占比更高，期限结构对短债滚动的压力较小。 公司披露的风险集中在监管与合规、股东与质押两类：一是 2025 年，公司未在 2025 年年度报告中披露重大诉讼、仲裁及处罚事项；二是 2026 年，一季报显示控股股东福建金森集团有限公司持股 1.50 亿股、占 63.63\%，其中质押 3。

\pfl{重大战略转型} 近三年公司的战略动作集中在：其他动作（3 项）：2023 年 4 月，公司对外投资设立全资子公司及注销分公司；2023 年 2 月，公司对外投资暨设立控股子公司。 公告口径上，近三年（2023 年 9 月---2026 年 9 月）公司披露重大重组与并购类公告 0 条、对外投资与转型类公告 5 条、回购与股权激励类公告 0 条，合计公告 235 条。 第一大行业“林业”的收入占比由 2021 年年报的 98.3\% 变为 74.4\%（23.8 个百分点），收入结构出现再平衡。 综合判断，报告期内公司有 5 条与战略相关的公告落地（重大重组与并购 0 条、对外投资与转型 5 条），资本运作与主业调整之间存在直接关联。

\begin{center}\small\setlength{\tabcolsep}{4pt}
\captionof{table}[福建金森（002679）历史融资与资本运作]{福建金森（002679）历史融资与资本运作（据公司公告与公开报道整理；金额为披露值，含首次公开发行、再融资、债券、银行借款与重整投资等）}
\begin{adjustbox}{max width=\textwidth}
\begin{tabular}{@{}p{1.45cm}p{2.55cm}p{4.05cm}p{4.95cm}@{}}
\toprule
时间 & 方式 & 规模 & 用途与说明 \\
\midrule
2012-06-05 & IPO（深圳证券交易所中小板） & 发行 3，\allowbreak{}468 万股，\allowbreak{}发行价 12.00 元/\allowbreak{}股，\allowbreak{}募集资金总额 4.16 亿元 & 商品材基地建设林木资源资产并购项目（原计划投入 2.57 亿元 \\
2015-12-02 & 资产转让（向控股股东出售资产） & 公司与控股股东福建金森集团有限公司签订资产转让协议（附生效条款），\allowbreak{}423.66 万元 & 优化资产结构、\allowbreak{}收缩非核心经营资产；该公告于 2015 年 12 月 4 日披露；交易对手为公司控股股东 \\
2023-12-31 & 现金分红 & 以 2.36 亿股为基数，\allowbreak{}向全体股东每 10 股派发现金红利 0.26 元（含税） & 回报股东；经公司董事会审议通过的 2023 年度利润分配预案 \\
2025-06-30 & 存货抵押贷款（银行融资） & 账面价值 7.84 亿元的存货全部抵押用于银行借款 & 生产经营周转资金；同一报告期受限资产还包括 4，\allowbreak{}250 元 ETC 保证金冻结货币资金；公司未披露具体授信额度 \\
\bottomrule
\end{tabular}
\end{adjustbox}
\end{center}

\pfl{历史融投资情况} 从现金流量表看，2025 年公司取得借款收到的现金 1.80 亿元、偿还债务支付的现金 1.79 亿元，筹资活动现金流净额 -0.48 亿元（2023 年为取得借款 5.58 亿元、偿还债务 4.19 亿元）。 2026 年上半年筹资活动现金流净额为 0.04 亿元（取得借款 1.00 亿元、偿还债务 0.77 亿元，未年化），仍处于净融入状态。 2025 年分配股利、利润或偿付利息支付的现金合计 0.49 亿元。 公开披露的融资记录共 4 笔，工具类型以 IPO、银行借款为主；金额与用途见下表。 其中规模最大的两笔为：2012 年 6 月，IPO，发行 3，468 万股，发行价 12.00 元/股，募集资金总额 4.16 亿元；2015 年 12 月，资产转让，公司与控股股东福建金森集团有限公司签订资产转让协议（附生效条款），423.66 万元。 股本方面，近三年披露股本变动 3 次；总股本自 2021 年末以来未发生变化（2.36 亿股）。 分红方面，近三年现金分红 4 次、合计每股派现 0.13 元，最近一次实施在 2025 年 12 月。 近三年“再融资”类公告 1 条，涉及定向增发。 综合看，公司融资结构上股权与债务融资并举；2025 年筹资活动现金流净额 -0.48 亿元，融资节奏仍以维持存量债务周转为主，与前述经营与投资现金流的方向基本一致。

\pfl{未来潜在融资需求分析} 2026 年 6 月末货币资金 1.34 亿元，而短期债务 1.75 亿元（现金短债比 0.77 倍），2026 年上半年经营现金流净额 -0.46 亿元。账面货币资金可以覆盖短期债务，缺口不体现在静态偿付层面。 公司 2026 年上半年归母净利润为负（-0.14 亿元），但 6 月末资产负债率 61.9\% 尚未进入第一档（亏损且负债率 $>$65\%）的区间，当前融资需求以补充流动资金、支撑低谷期经营性支出为主。 公开披露的后续资金安排线索包括：债务与授信安排：2025 年 6 月，公司主要融资方式为以存货（林木资产）抵押取得银行借款。 以现有货币资金与经营现金流衡量，公司不存在刚性补血的紧迫性，融资动作更多是主动型的。




\subsubsection{大湖股份（600257）}
\label{co:600257}

\pfl{公司基本情况} 湖南常德起家的淡水养殖上市公司（曾名“洞庭水殖”“大湖水殖”，自称“中国淡水渔业第一股”），2000 年 6 月在上交所上市；2016 年定增募资 5.5 亿元后由水产养殖向水产加工、白酒和康复护理医疗延伸，2020—2021 年分步收购并增资东方华康医疗管理有限公司取得 60\% 股权（2026 年 4 月增至 61.1103\%），2025 年 10 月公司全称变更为“大湖健康产业股份有限公司”；2025 年营业收入 9.20 亿元、归母净利润 506.53 万元，2026 年上半年归母净利润 1,688.13 万元、同比扭亏。 公司于 2000-06-12 在上交所首发上市，注册资本 4.81 亿元，总股本 4.81 亿股。 截至 2026 年 9 月 24 日收盘价 5.48 元，流通市值 26.4 亿元，近一年+4.8\%，流通市值在三级行业“水产养殖”的 4 家公司中按由大到小排第\zhnumber{3}位，近一年涨跌幅高于全样本中位数（-10.1\%）14.9 个百分点。 按 2025 年年报披露的行业构成，收入占比前三位为农业 47.7\%；其他 41.9\%；工业 8.8\%。

\pfl{历史股价} 自 2021 年 1 月至 2026 年 9 月，股价由 5.02 元变为 5.48 元，累计 +9.2\%，区间最高 9.23 元（2023 年 8 月）、最低 4.29 元（2024 年 8 月）；当前价相当于区间最高价的 59.4\%，为区间最低价的 1.28 倍。 月度收益的年化波动率 48.1\%，高于全样本中位数 37.4\%，在三级行业“水产养殖”的 4 家公司中波动率由高到低排第\zhnumber{1}位。 把股价阶段与公司事件对齐看，2022 年 10 月 至 2023 年 8 月股价上涨+105.1\%，同期（2023 年 7 月）公司此前已于 2023 年 1 月 19 日、2 月 6 日审议通过延长该次非公开发行股东大会决议有效期及授权有效期至 2024 年 2 月 13 日；2024 年 8 月 至 2026 年 1 月股价上涨+72.3\%，同期（2026 年 1 月）公司以 700.00 万元购买徐培麒持有的嘉兴智康 35\% 股权、以 100.00 万元购买上海麒聚瑆伙机器人科技合伙企业（有限合伙）持有的 5\% 股权（合计 800.00 万元受让 40\% 股权）。 整体上看，区间的几处大级别拐点都能在公司的资本运作、股权变动或风险事件上找到对应，股价波动有明显的公司层面驱动，而非单纯的行业 beta。

\begin{center}
\begin{minipage}{\textwidth}\centering
\begin{tikzpicture}
\begin{axis}[
  width=12.6cm, height=3.9cm, scale only axis,
  xmin=2020.922, xmax=2026.888, ymin=3.99, ymax=9.88,
  xtick={2021,2022,2023,2024,2025,2026},
  yearticks,
  yticklabel style={font=\scriptsize},
  ylabel={元/股}, ylabel style={font=\scriptsize},
  grid=major, grid style={LightGray, line width=0.3pt},
  axis line style={InkGray, line width=0.5pt},
  every axis plot/.append style={line width=0.9pt},
]
\addplot[AgriGreen] table[x=x,y=y]{data/clean/profiles/px_600257.dat};
\addplot[SoftRed, dashed, line width=0.7pt] table[x=x,y=y]{data/clean/profiles/pxzz_600257.dat};
\addplot[only marks, mark=*, mark size=1.1pt, SoftRed] table[x=x,y=y]{data/clean/profiles/pxzz_600257.dat};
\end{axis}
\end{tikzpicture}

\captionof{figure}[大湖股份（600257）股价与周期骨架]{大湖股份（600257）月度收盘价（元/股）与 ZigZag 周期骨架（阈值 25\%，圆点为周期转折点）}
\end{minipage}
\end{center}

\pfl{过去周期分析} 以 25\% 的 ZigZag 阈值划分，样本区间内股价形成 5 个主升段与 5 个主跌段。 最大主跌段为 2023 年 8 月 至 2024 年 2 月，6 个月-45.9\%；最大主升段出现在 2022 年 10 月 至 2023 年 8 月，10 个月+105.1\%。 最近一次转折出现在 2026 年 9 月（上涨段自 2026 年 6 月起，3 个月+26.6\%），当前处于该段之后的震荡之中。 公司股价较申万农林牧渔指数提前约 21 个月见底，是板块中较早定价周期反转的公司之一。 水产养殖的价格周期由投苗量与消费需求共同决定（第 \ref{sec:other-livestock} 节），与猪周期的同步性弱于禽链，公司 2026 年上半年实现盈利、半年 ROE 2.1\%（未年化），好于全样本中位数（0.75\%），下行周期对其报表的冲击小于同业。 综合区间形态看，该股单段涨跌幅度偏大、以趋势段为主（最大主升 +105.1\%、最大主跌 -45.9\%），年化波动率 48.1\% 与全样本中位数 37.4\% 相比波动明显大于样本中位数。

\pfl{盈利情况} 2026 年上半年营业收入 4.34 亿元（同比 +2.0\%），归母净利润 0.17 亿元，实现扭亏。 以年报为趋势对照，2023---2025 年营业收入由 11.93 亿元变为 9.20 亿元（两年年均复合增速 -12.2\%），归母净利润由 -0.08 亿元变为 0.05 亿元。 按半年报口径，2026 年上半年的 ROE 为 2.1\%（未年化）、销售净利率 7.8\%、毛利率 29.4\%，ROE 在“水产养殖”4 家公司中按数值由高到低排第\zhnumber{3}位。 同期全样本半年 ROE 中位数 0.75\%，公司与之相差 1.33 个百分点（略好于中位数公司）。 面对这一局面，公司披露的举措包括：2025 年 10 月，公司全称变更为“大湖健康产业股份有限公司”。 综合看，公司在报告期维持盈利（高于全样本半年 ROE 中位数 0.75\%），利润的可持续性取决于水产养殖景气的延续与成本管控的执行。

\begin{center}\small\setlength{\tabcolsep}{3.4pt}
\captionof{table}[大湖股份（600257）关键财务指标]{大湖股份（600257）关键财务指标（合并报表口径；两个半年报期均未年化；现金短债比 $=$ 货币资金 $\div$ 短期债务）}
\begin{adjustbox}{max width=\textwidth}
\begin{tabular}{@{}lrrrrr@{}}
\toprule
指标 & 2023 年 & 2024 年 & 2025 年 & 2025 年半年报 & 2026 年半年报 \\
\midrule
营业收入（亿元） & 11.93 & 10.44 & 9.20 & 4.26 & 4.34 \\
归母净利润（亿元） & -0.08 & -0.77 & 0.05 & -0.03 & 0.17 \\
毛利率（\%） & 24.5 & 23.1 & 27.6 & 25.0 & 29.4 \\
ROE（\%） & -0.9 & -9.2 & 0.6 & -0.3 & 2.1 \\
资产负债率（\%） & 51.0 & 48.6 & 46.5 & 48.0 & 46.4 \\
经营活动现金流净额（亿元） & 1.62 & 0.67 & 2.18 & 0.84 & 0.56 \\
货币资金（亿元） & 3.36 & 1.87 & 3.22 & 2.33 & 2.61 \\
有息负债合计（亿元） & 8.68 & 7.30 & 6.90 & 7.32 & 6.87 \\
现金短债比（倍） & 1.04 & 0.67 & 1.08 & 0.74 & 1.04 \\
\bottomrule
\end{tabular}
\end{adjustbox}
\end{center}

\begin{center}
\begin{minipage}{\textwidth}\centering
\begin{tikzpicture}
\begin{groupplot}[
  group style={group size=2 by 1, horizontal sep=0.60cm},
  width=6.85cm, height=4.60cm,
  tick label style={font=\tiny},
  title style={font=\scriptsize\heiti},
  yticklabel style={font=\tiny},
  ylabel style={font=\tiny},
  grid=major, grid style={LightGray, line width=0.3pt},
  axis line style={InkGray, line width=0.5pt},
  enlarge x limits={abs=0.8},
  legend style={font=\scriptsize, at={(0.5,-0.20)}, anchor=north,
                legend columns=2, draw=none, fill=none, column sep=6pt},
]
\nextgroupplot[
  ybar, bar width=4pt,
  ymin=-3.27, ymax=14.69,
  xtick={0,1,2,3,4,5.7},
  xticklabels={2021,2022,2023,2024,2025,2026H1},
  ylabel={亿元},
]
\addplot[fill=AgriGreen!75, draw=AgriGreen, bar shift=-2.6pt] table[x=i, y=rev]{data/clean/profiles/fin_600257.dat};
\addplot[fill=SoftRed!60, draw=SoftRed, bar shift=2.6pt] table[x=i, y=ni]{data/clean/profiles/fin_600257.dat};
\addplot[fill=AgriGreen!30, draw=AgriGreen, pattern=north east lines, bar shift=-2.6pt] table[x=x, y=rev]{data/clean/profiles/finh_600257.dat};
\addplot[fill=SoftRed!20, draw=SoftRed, pattern=north east lines, bar shift=2.6pt] table[x=x, y=ni]{data/clean/profiles/finh_600257.dat};
\legend{营业收入（年/半年）, 归母净利润（年/半年）}
\nextgroupplot[
  ymin=-22.9, ymax=61.2,
  xtick={0,1,2,3,4,5.7},
  xticklabels={2021,2022,2023,2024,2025,2026H1},
  ylabel={\%},
]
\addplot[AgriGold, mark=*, mark size=1.3pt] table[x=i, y=roe]{data/clean/profiles/fin_600257.dat};
\addplot[OilBlue, mark=square*, mark size=1.3pt] table[x=i, y=debt]{data/clean/profiles/fin_600257.dat};
\addplot[only marks, AgriGold, mark=*, mark size=2.0pt] table[x=x, y=roe]{data/clean/profiles/finh_600257.dat};
\addplot[only marks, OilBlue, mark=square*, mark size=2.0pt] table[x=x, y=debt]{data/clean/profiles/finh_600257.dat};
\legend{ROE, 资产负债率}
\end{groupplot}
\end{tikzpicture}
\begin{tikzpicture}
\begin{axis}[
  name=finbot,
  width=13.75cm, height=3.10cm, scale only axis,
  ymin=-0.95, ymax=10.69,
  xtick={0,1,2,3,4,5.7},
  xticklabels={2021,2022,2023,2024,2025,2026H1},
  tick label style={font=\tiny},
  yticklabel style={font=\tiny},
  ylabel={亿元}, ylabel style={font=\tiny},
  ybar, bar width=4pt,
  grid=major, grid style={LightGray, line width=0.3pt},
  axis line style={InkGray, line width=0.5pt},
  enlarge x limits={abs=0.8},
  legend style={font=\scriptsize, at={(0.5,-0.26)}, anchor=north,
                legend columns=2, draw=none, fill=none, column sep=10pt},
]
\addplot[fill=AgriGreen!55, draw=AgriGreen, bar shift=-3.0pt] table[x=i, y=cash]{data/clean/profiles/fin_600257.dat};
\addplot[fill=SoftRed!45, draw=SoftRed, bar shift=3.0pt] table[x=i, y=ibd]{data/clean/profiles/fin_600257.dat};
\addplot[fill=AgriGreen!25, draw=AgriGreen, pattern=north east lines, bar shift=-3.0pt] table[x=x, y=cash]{data/clean/profiles/finh_600257.dat};
\addplot[fill=SoftRed!20, draw=SoftRed, pattern=north east lines, bar shift=3.0pt] table[x=x, y=ibd]{data/clean/profiles/finh_600257.dat};
\legend{货币资金（左轴）, 有息负债（左轴）}
\end{axis}
\begin{axis}[
  at={(finbot.south east)}, anchor=south east,
  width=13.75cm, height=3.10cm, scale only axis,
  axis y line*=right, axis x line*=none, xtick=\empty,
  ymin=-0.57, ymax=2.83,
  tick label style={font=\tiny},
  yticklabel style={font=\tiny}, scaled y ticks=false,
  legend style={font=\scriptsize, at={(0.5,-0.44)}, anchor=north,
                legend columns=1, draw=none, fill=none},
]
\addplot[OilBlue, mark=*, mark size=2.2pt, line width=1.3pt] table[x=i, y=ocf]{data/clean/profiles/fin_600257.dat};
\addplot[only marks, OilBlue, mark=*, mark size=3.2pt] table[x=x, y=ocf]{data/clean/profiles/finh_600257.dat};
\legend{经营活动现金流净额（右轴，亿元，蓝点）}
\end{axis}
\end{tikzpicture}

\captionof{figure}[大湖股份（600257）近年主要财务指标]{大湖股份（600257）近年主要财务指标（左上：营业收入与归母净利润；右上：ROE 与资产负债率；下：货币资金、有息负债为左轴柱状，经营活动现金流净额为其右轴的蓝点折线。
2021---2025 年为年报口径，2026H1 为 2026 年半年报/半年末，与其前一年报之间留出间隔、以斜纹或空心点区分；半年报数据未年化）}
\end{minipage}
\end{center}

\pfl{财务分析} 2026 年 6 月末资产负债率 46.4\%，较 2025 年末下降 0.1 个百分点（样本中位数 52.2\%，所处三级行业内按由低到高排第\zhnumber{1}位）。 2026 年 6 月末货币资金 2.61 亿元、短期债务（短期借款与一年内到期非流动负债）2.50 亿元、长期债务 4.36 亿元，有息负债合计 6.87 亿元，现金短债比 1.04 倍——覆盖充分。 净有息负债 4.25 亿元，相当于其 2026 年上半年营业收入的 98\%。 流动比率 2.24、速动比率 1.20，短期指标尚可。 2026 年上半年经营活动现金流净额 0.56 亿元，经营现金流与利润同向为正，内生资金可以支撑运营。 2026 年 6 月末商誉 1.45 亿元，占归母净资产的 14.8\%，是净资产中最脆弱的一块。 综合看，现金短债比 1.04 倍，短期偿付压力可控；资产负债率低于样本中位数 5.8 个百分点；有息负债中长期债务占比更高，期限结构对短债滚动的压力较小。 公司披露的风险集中在重整与债务违约、股东与质押两类：一是 2024 年 10 月，公司拟追回已支付的违规费用 600.34 万元并收取 30\% 违约金 180.10 万元；二是 2025 年 12 月，子公司净资产份额 2.95 亿元用于银行借款质押。

\pfl{重大战略转型} 公司的战略动作可归为以下几类：并购与重组（5 项）：2026 年 1 月，公司以 700.00 万元购买徐培麒持有的嘉兴智康 35\% 股权、以 100.00 万元购买上海麒聚瑆伙机器人科技合伙企业（有限合伙）持有的 5\% 股权（合计 800.00 万元受让 40\% 股权）；2024 年 8 月，下属子公司皂市渔业与汉寿久达公司签订大湖水殖汉寿中华鳖有限公司的股权转让协议，以人民币 47.92 万元转让皂市渔业持有的汉寿中华鳖 72.7778\% 股权；其他动作（2 项）：2025 年 3 月，公司与子公司东方华康共同投资设立上海筷智医疗科技有限公司。 公告口径上，近三年（2023 年 9 月---2026 年 9 月）公司披露重大重组与并购类公告 1 条、对外投资与转型类公告 1 条、回购与股权激励类公告 0 条，合计公告 275 条。 第一大行业“农业”的收入占比由 2021 年年报的 58.3\% 变为 47.7\%（10.6 个百分点），收入结构出现再平衡。 综合判断，报告期内公司有 2 条与战略相关的公告落地（重大重组与并购 1 条、对外投资与转型 1 条），资本运作与主业调整之间存在直接关联。

\begin{center}\small\setlength{\tabcolsep}{4pt}
\captionof{table}[大湖股份（600257）历史融资与资本运作]{大湖股份（600257）历史融资与资本运作（据公司公告与公开报道整理；金额为披露值，含首次公开发行、再融资、债券、银行借款与重整投资等）}
\begin{adjustbox}{max width=\textwidth}
\begin{tabular}{@{}p{1.45cm}p{2.55cm}p{4.05cm}p{4.95cm}@{}}
\toprule
时间 & 方式 & 规模 & 用途与说明 \\
\midrule
2000-06-12 & IPO（上海证券交易所） & 实际发行 4，\allowbreak{}000.00 万股，\allowbreak{}发行价 8.90 元/\allowbreak{}股；实际募集资金合计 3.56 亿元 & 财证券有限责任公司；招股公告日 2000 年 5 月 11 日 \\
2016-10-18 & 非公开发行股票（定向配售、网下定价发行） & 发行价格 10.15 元，\allowbreak{}实际发行数量 5，\allowbreak{}418.7188 万股 & 公司 2025 年年度报告披露本次发行 5，\allowbreak{}419 万股、\allowbreak{}共募集资金 5.50 亿元 \\
2017-12 & 发行股份及支付现金购买资产并募集配套资金（重大资产重组，未完成） & 拟以 10.40 亿元收购西藏深万投实业有限公司 51\% 股权 & 用于支付本次交易中的现金对价、\allowbreak{}相关税费和中介机构费用（原募集配套资金方案）；标的公司核心资产为深圳市万生堂实业有限公司（日本冈本安全套中国内地总经销商） \\
2018-11-05 & 重大资产重组报告书（草案）披露 & 交易总价 10.40 亿元；募集配套资金不超过 3.50 亿元（配套股份数量不超过发行前总股本 4.81 亿股的 20\% & 支付本次交易中的现金对价、\allowbreak{}支付本次交易相关税费和中介机构费用；招商证券作为独立财务顾问披露；交易完成后公司将持有西藏深万投 51\% 股权 \\
2020-04-14 & 现金收购子公司股权及增资 & 以 2.00 亿元收购东方华康原有股东 32\% 股权 & 布局康复护理医疗产业，\allowbreak{}公司持有东方华康 60\% 股权 \\
2023-07-21 & 终止前次非公开发行股票 & --- & 第九届董事会第二次会议、\allowbreak{}第九届监事会第二次会议审议通过公司终止前次非公开发行股票事项 \\
2023-08-31 & 向特定对象发行 A 股股票（预案修订稿，未实施） & 拟募集资金总额（含发行费用）不超过 5.00 亿元 & 大湖水殖股份有限公司冰鲜冻鲜水产品深加工等项目 \\
2026-04-13 & 业绩补偿方以股权抵偿（受让子公司股权） & 业绩补偿方咖辅健康科技（上海）有限公司、\allowbreak{}上海擢英医疗管理合伙企业（有限合伙）、\allowbreak{}李爱川将其合计持有并已依法实缴的东方华康 1.1103\% 股权无偿转让给公司 & 抵偿业绩补偿方应向公司支付的业绩补偿金；交易完成后公司持有东方华康的比例由 60\% 增加至 61.1103\% \\
\bottomrule
\end{tabular}
\end{adjustbox}
\end{center}

\pfl{历史融投资情况} 从现金流量表看，2025 年公司取得借款收到的现金 2.73 亿元、偿还债务支付的现金 2.83 亿元，筹资活动现金流净额 -0.81 亿元（2023 年为取得借款 3.22 亿元、偿还债务 3.21 亿元）。 2026 年上半年筹资活动现金流净额为 -0.56 亿元（取得借款 1.79 亿元、偿还债务 1.62 亿元，未年化），已转为净流出，处于净还债状态。 2025 年分配股利、利润或偿付利息支付的现金合计 0.12 亿元。 公开披露的融资记录共 8 笔，工具类型以 IPO、定向增发、发行股份及支付现金购为主；金额与用途见下表。 其中规模最大的两笔为：2023 年 8 月，定向增发，拟募集资金总额（含发行费用）不超过 5.00 亿元；2016 年 10 月，定向增发，发行价格 10.15 元，实际发行数量 5，418.7188 万股。 股本方面，近三年披露股本变动 3 次；总股本自 2021 年末以来未发生变化（4.81 亿股）。 分红方面，近三年未实施现金分红，在全部样本中属于分红能力偏弱的一端。 近三年“再融资”类公告 2 条，涉及定向增发。 综合看，公司融资结构上股权与债务融资并举；2025 年筹资活动现金流净额 -0.81 亿元，融资节奏仍以维持存量债务周转为主，与前述经营与投资现金流的方向基本一致。

\pfl{未来潜在融资需求分析} 2026 年 6 月末货币资金 2.61 亿元，而短期债务 2.50 亿元（现金短债比 1.04 倍），2026 年上半年经营现金流净额 0.56 亿元。账面货币资金可以覆盖短期债务，缺口不体现在静态偿付层面。 公司 2026 年 6 月末资产负债率 46.4\%、2026 年上半年 ROE 2.1\%（未年化），财务结构相对宽松而盈利弹性有限，融资属于可选项而非必选项，触发条件多与产能建设或外延并购相关。 公开披露的后续资金安排线索包括：股权融资线索：2023 年 8 月，公司 2023 年度向特定对象发行 A 股股票预案（修订稿）；2024 年，第二次临时股东会审议通过 7 项议案。 以现有货币资金与经营现金流衡量，公司不存在刚性补血的紧迫性，融资动作更多是主动型的。




\subsubsection{獐子岛（002069）}
\label{co:002069}

\pfl{公司基本情况} 辽宁大连长海县獐子岛镇起家的海洋牧场企业，2006 年 9 月深交所上市，A 股“水产第一股”，因 2014 年以来多次“扇贝绝收/跑路”及 2016—2017 年年度报告虚假记载于 2020 年 6 月被证监会处罚、原董事长吴厚刚被终身市场禁入并获刑 15 年；2022 年 3 月原控股股东所持股份被拍卖、大连市国资入主，2022 年 5 月被实施其他风险警示（ST 獐子岛）、2023 年 5 月撤销；现由大连市国资委实际控制，公司以“扭亏、瘦身、聚焦、增效”为主线剥离非核心资产，2026 年上半年营业收入 7.31 亿元、归母净利润 1,287.14 万元。 公司于 2006-09-28 在深交所主板首发上市，注册资本 7.11 亿元，总股本 7.11 亿股。 截至 2026 年 9 月 24 日收盘价 3.67 元，流通市值 26.1 亿元，近一年-7.1\%，流通市值在三级行业“水产养殖”的 4 家公司中按由大到小排第\zhnumber{4}位，近一年涨跌幅高于全样本中位数（-10.1\%）3.0 个百分点。 按 2026 年半年报披露的行业构成，收入占比前三位为水产贸易业 41.3\%；水产加工业 40.5\%；水产养殖业 13.9\%。 股权结构方面，实际控制人为大连市国资委。

\pfl{历史股价} 自 2021 年 1 月至 2026 年 9 月，股价由 3.23 元变为 3.67 元，累计 +13.6\%，区间最高 4.95 元（2023 年 7 月）、最低 2.42 元（2024 年 6 月）；当前价相当于区间最高价的 74.1\%，为区间最低价的 1.52 倍。 月度收益的年化波动率 34.2\%，低于全样本中位数 37.4\%，在三级行业“水产养殖”的 4 家公司中波动率由高到低排第\zhnumber{3}位。 把股价阶段与公司事件对齐看，2024 年 6 月 至 2025 年 8 月股价上涨+80.2\%，同期（2024 年）公司另收到船厂土地及附着物收购补偿款 5，860.00 万元；2022 年 5 月 至 2023 年 7 月股价上涨+54.2\%，同期（2022 年 10 月）被判处有期徒刑十五年并处罚金 92.00 万元。 从时间对应关系看，公司层面的资本运作与股权、风险事件解释了区间内多数急涨急跌，与行业指数的同步性相对有限。

\begin{center}
\begin{minipage}{\textwidth}\centering
\begin{tikzpicture}
\begin{axis}[
  width=12.6cm, height=3.9cm, scale only axis,
  xmin=2020.922, xmax=2026.888, ymin=2.25, ymax=5.30,
  xtick={2021,2022,2023,2024,2025,2026},
  yearticks,
  yticklabel style={font=\scriptsize},
  ylabel={元/股}, ylabel style={font=\scriptsize},
  grid=major, grid style={LightGray, line width=0.3pt},
  axis line style={InkGray, line width=0.5pt},
  every axis plot/.append style={line width=0.9pt},
]
\addplot[AgriGreen] table[x=x,y=y]{data/clean/profiles/px_002069.dat};
\addplot[SoftRed, dashed, line width=0.7pt] table[x=x,y=y]{data/clean/profiles/pxzz_002069.dat};
\addplot[only marks, mark=*, mark size=1.1pt, SoftRed] table[x=x,y=y]{data/clean/profiles/pxzz_002069.dat};
\end{axis}
\end{tikzpicture}

\captionof{figure}[獐子岛（002069）股价与周期骨架]{獐子岛（002069）月度收盘价（元/股）与 ZigZag 周期骨架（阈值 25\%，圆点为周期转折点）}
\end{minipage}
\end{center}

\pfl{过去周期分析} 以 25\% 的 ZigZag 阈值划分，样本区间内股价形成 3 个主升段与 3 个主跌段。 最大主跌段为 2023 年 7 月 至 2024 年 6 月，11 个月-51.1\%；最大主升段出现在 2024 年 6 月 至 2025 年 8 月，14 个月+80.2\%。 最近一次转折出现在 2026 年 6 月（下跌段自 2025 年 8 月起，10 个月-32.6\%），当前处于该段的延续之中。 公司股价较申万农林牧渔指数提前约 23 个月见底，是板块中较早定价周期反转的公司之一。 水产养殖的价格周期由投苗量与消费需求共同决定（第 \ref{sec:other-livestock} 节），与猪周期的同步性弱于禽链，公司 2026 年 6 月末资产负债率 97.0\%，高于全样本中位数 52.2\%，其股价因此更多反映资金面与信用风险，而不是价格弹性本身。 综合区间形态看，该股单段涨跌幅度偏大、以趋势段为主（最大主升 +80.2\%、最大主跌 -51.1\%），年化波动率 34.2\% 与全样本中位数 37.4\% 相比波动与样本中位数接近。

\pfl{盈利情况} 2026 年上半年营业收入 7.31 亿元（同比 -5.3\%），归母净利润 0.13 亿元，同比 +26.8\%。 以年报为趋势对照，2023---2025 年营业收入由 16.77 亿元变为 12.78 亿元（两年年均复合增速 -12.7\%），归母净利润由 0.09 亿元变为 -0.11 亿元。 按半年报口径，2026 年上半年的 ROE 为 29.5\%（未年化）、销售净利率 2.5\%、毛利率 18.2\%，ROE 在“水产养殖”4 家公司中按数值由高到低排第\zhnumber{1}位。 同期全样本半年 ROE 中位数 0.75\%，公司与之相差 28.77 个百分点（略好于中位数公司）。 综合看，公司在报告期维持盈利（高于全样本半年 ROE 中位数 0.75\%），利润的可持续性取决于水产养殖景气的延续与成本管控的执行。

\begin{center}\small\setlength{\tabcolsep}{3.4pt}
\captionof{table}[獐子岛（002069）关键财务指标]{獐子岛（002069）关键财务指标（合并报表口径；两个半年报期均未年化；现金短债比 $=$ 货币资金 $\div$ 短期债务）}
\begin{adjustbox}{max width=\textwidth}
\begin{tabular}{@{}lrrrrr@{}}
\toprule
指标 & 2023 年 & 2024 年 & 2025 年 & 2025 年半年报 & 2026 年半年报 \\
\midrule
营业收入（亿元） & 16.77 & 15.83 & 12.78 & 7.72 & 7.31 \\
归母净利润（亿元） & 0.09 & -0.22 & -0.11 & 0.10 & 0.13 \\
毛利率（\%） & 17.4 & 11.2 & 15.9 & 16.1 & 18.2 \\
ROE（\%） & 12.3 & -32.7 & -23.2 & 17.9 & 29.5 \\
资产负债率（\%） & 94.7 & 95.8 & 96.8 & 95.2 & 97.0 \\
经营活动现金流净额（亿元） & 1.90 & 0.79 & 0.42 & 0.52 & 0.26 \\
货币资金（亿元） & 5.53 & 6.04 & 5.80 & 6.15 & 5.45 \\
有息负债合计（亿元） & 20.20 & 19.48 & 19.21 & 19.44 & 19.24 \\
现金短债比（倍） & 0.31 & 0.31 & 0.49 & 0.36 & 0.39 \\
\bottomrule
\end{tabular}
\end{adjustbox}
\end{center}

\begin{center}
\begin{minipage}{\textwidth}\centering
\begin{tikzpicture}
\begin{groupplot}[
  group style={group size=2 by 1, horizontal sep=0.60cm},
  width=6.85cm, height=4.60cm,
  tick label style={font=\tiny},
  title style={font=\scriptsize\heiti},
  yticklabel style={font=\tiny},
  ylabel style={font=\tiny},
  grid=major, grid style={LightGray, line width=0.3pt},
  axis line style={InkGray, line width=0.5pt},
  enlarge x limits={abs=0.8},
  legend style={font=\scriptsize, at={(0.5,-0.20)}, anchor=north,
                legend columns=2, draw=none, fill=none, column sep=6pt},
]
\nextgroupplot[
  ybar, bar width=4pt,
  ymin=-2.32, ymax=23.35,
  xtick={0,1,2,3,4,5.7},
  xticklabels={2021,2022,2023,2024,2025,2026H1},
  ylabel={亿元},
]
\addplot[fill=AgriGreen!75, draw=AgriGreen, bar shift=-2.6pt] table[x=i, y=rev]{data/clean/profiles/fin_002069.dat};
\addplot[fill=SoftRed!60, draw=SoftRed, bar shift=2.6pt] table[x=i, y=ni]{data/clean/profiles/fin_002069.dat};
\addplot[fill=AgriGreen!30, draw=AgriGreen, pattern=north east lines, bar shift=-2.6pt] table[x=x, y=rev]{data/clean/profiles/finh_002069.dat};
\addplot[fill=SoftRed!20, draw=SoftRed, pattern=north east lines, bar shift=2.6pt] table[x=x, y=ni]{data/clean/profiles/finh_002069.dat};
\legend{营业收入（年/半年）, 归母净利润（年/半年）}
\nextgroupplot[
  ymin=-43.0, ymax=110.0,
  xtick={0,1,2,3,4,5.7},
  xticklabels={2021,2022,2023,2024,2025,2026H1},
  ylabel={\%},
]
\addplot[AgriGold, mark=*, mark size=1.3pt] table[x=i, y=roe]{data/clean/profiles/fin_002069.dat};
\addplot[OilBlue, mark=square*, mark size=1.3pt] table[x=i, y=debt]{data/clean/profiles/fin_002069.dat};
\addplot[only marks, AgriGold, mark=*, mark size=2.0pt] table[x=x, y=roe]{data/clean/profiles/finh_002069.dat};
\addplot[only marks, OilBlue, mark=square*, mark size=2.0pt] table[x=x, y=debt]{data/clean/profiles/finh_002069.dat};
\legend{ROE, 资产负债率}
\end{groupplot}
\end{tikzpicture}
\begin{tikzpicture}
\begin{axis}[
  name=finbot,
  width=13.75cm, height=3.10cm, scale only axis,
  ymin=-2.06, ymax=23.03,
  xtick={0,1,2,3,4,5.7},
  xticklabels={2021,2022,2023,2024,2025,2026H1},
  tick label style={font=\tiny},
  yticklabel style={font=\tiny},
  ylabel={亿元}, ylabel style={font=\tiny},
  ybar, bar width=4pt,
  grid=major, grid style={LightGray, line width=0.3pt},
  axis line style={InkGray, line width=0.5pt},
  enlarge x limits={abs=0.8},
  legend style={font=\scriptsize, at={(0.5,-0.26)}, anchor=north,
                legend columns=2, draw=none, fill=none, column sep=10pt},
]
\addplot[fill=AgriGreen!55, draw=AgriGreen, bar shift=-3.0pt] table[x=i, y=cash]{data/clean/profiles/fin_002069.dat};
\addplot[fill=SoftRed!45, draw=SoftRed, bar shift=3.0pt] table[x=i, y=ibd]{data/clean/profiles/fin_002069.dat};
\addplot[fill=AgriGreen!25, draw=AgriGreen, pattern=north east lines, bar shift=-3.0pt] table[x=x, y=cash]{data/clean/profiles/finh_002069.dat};
\addplot[fill=SoftRed!20, draw=SoftRed, pattern=north east lines, bar shift=3.0pt] table[x=x, y=ibd]{data/clean/profiles/finh_002069.dat};
\legend{货币资金（左轴）, 有息负债（左轴）}
\end{axis}
\begin{axis}[
  at={(finbot.south east)}, anchor=south east,
  width=13.75cm, height=3.10cm, scale only axis,
  axis y line*=right, axis x line*=none, xtick=\empty,
  ymin=-0.49, ymax=2.47,
  tick label style={font=\tiny},
  yticklabel style={font=\tiny}, scaled y ticks=false,
  legend style={font=\scriptsize, at={(0.5,-0.44)}, anchor=north,
                legend columns=1, draw=none, fill=none},
]
\addplot[OilBlue, mark=*, mark size=2.2pt, line width=1.3pt] table[x=i, y=ocf]{data/clean/profiles/fin_002069.dat};
\addplot[only marks, OilBlue, mark=*, mark size=3.2pt] table[x=x, y=ocf]{data/clean/profiles/finh_002069.dat};
\legend{经营活动现金流净额（右轴，亿元，蓝点）}
\end{axis}
\end{tikzpicture}

\captionof{figure}[獐子岛（002069）近年主要财务指标]{獐子岛（002069）近年主要财务指标（左上：营业收入与归母净利润；右上：ROE 与资产负债率；下：货币资金、有息负债为左轴柱状，经营活动现金流净额为其右轴的蓝点折线。
2021---2025 年为年报口径，2026H1 为 2026 年半年报/半年末，与其前一年报之间留出间隔、以斜纹或空心点区分；半年报数据未年化）}
\end{minipage}
\end{center}

\pfl{财务分析} 2026 年 6 月末资产负债率 97.0\%，较 2025 年末上升 0.2 个百分点（样本中位数 52.2\%，所处三级行业内按由低到高排第\zhnumber{4}位）。 2026 年 6 月末货币资金 5.45 亿元、短期债务（短期借款与一年内到期非流动负债）14.05 亿元、长期债务 5.19 亿元，有息负债合计 19.24 亿元，现金短债比 0.39 倍——覆盖偏紧。 净有息负债 13.79 亿元，相当于其 2026 年上半年营业收入的 189\%。 流动比率 0.99、速动比率 0.47，短期偿债指标偏紧。 2026 年上半年经营活动现金流净额 0.26 亿元，经营现金流与利润同向为正，内生资金可以支撑运营。 综合看，账面货币资金对有息负债与短期债务的覆盖不足，滚动偿付对新增融资的依赖度较高；资产负债率高于样本中位数 44.8 个百分点；有息负债中短期债务占比更高，期限结构仍以滚动续作为主。 公司披露的诉讼与资产冻结风险为：2025 年 1 月，公司及控股子公司连续 12 个月内累计诉讼、仲裁事项涉案金额共计 1，082.28 万元。

\pfl{重大战略转型} 近三年公司的战略动作集中在：其他动作（3 项）：2026 年 3 月，公司拟与大连獐子岛海洋发展集团有限公司共同投资设立大连海发生态渔业发展有限公司；2022 年 3 月，原控股股东长海县獐子岛投资发展中心持有的 1.0996 亿股被大连盐化集团以 3.43 亿元拍下；并购与重组（2 项）：2025 年，公司以“扭亏、瘦身、聚焦、增效”为主线推进系统性改革；2019 年 7 月，公司将全资子公司獐子岛渔业集团香港有限公司持有的大连新中海产食品有限公司 100\% 股权、新中日本株式会社 90\% 股权出售给亚洲渔港股份有限公司。 公告口径上，近三年（2023 年 9 月---2026 年 9 月）公司披露重大重组与并购类公告 5 条、对外投资与转型类公告 2 条、回购与股权激励类公告 0 条，合计公告 307 条。 第一大行业仍是“水产贸易业”，收入占比 41.3\%，与 2021 年年报相比没有方向性变化。 综合判断，报告期内公司有 7 条与战略相关的公告落地（重大重组与并购 5 条、对外投资与转型 2 条），资本运作与主业调整之间存在直接关联。

\begin{center}\small\setlength{\tabcolsep}{4pt}
\captionof{table}[獐子岛（002069）历史融资与资本运作]{獐子岛（002069）历史融资与资本运作（据公司公告与公开报道整理；金额为披露值，含首次公开发行、再融资、债券、银行借款与重整投资等）}
\begin{adjustbox}{max width=\textwidth}
\begin{tabular}{@{}p{1.45cm}p{2.55cm}p{4.05cm}p{4.95cm}@{}}
\toprule
时间 & 方式 & 规模 & 用途与说明 \\
\midrule
2006-09-28 & IPO（深圳证券交易所） & 实际发行 2，\allowbreak{}830.00 万股 & 限责任公司，\allowbreak{}发行市盈率（按发行后总股本）18.80 倍 \\
2019-08-31 & 重大资产出售（向亚渔食品出售新中海产 100\% 股权及新中日本 90\% 股权） & 交易对价 2.35 亿元（拟出售资产于估值基准日 2019 年 5 月 31 日的估值 2.41 亿元） & 回笼资金用于集中资源加快海洋牧场重构、\allowbreak{}降低资产负债率、\allowbreak{}保障公司安全运营；交易对方亚渔食品为亚洲渔港股份有限公司全资子公司 \\
2020 & 资产处置回笼资金 & 转让位于长海县广鹿岛的 4 宗海域使用的租赁权暨海底存货，\allowbreak{}已按约定收到首付款和进度款合计 9，\allowbreak{}678.00 万元 & 盘活存量资产；经 2020 年 1 月 3 日第七届董事会第九次会议及 2020 年 2 月 3 日 2020 年第一次临时股东大会审议通过公司转让海域使用的租赁权暨海底存货 \\
2022 & 出售分公司资产 & 将大连乌蟒岛分公司的存货、\allowbreak{}固定资产、\allowbreak{}无形资产等相关资产通过大连产权交易所挂牌转让给大连智慧渔业集团有限公司，\allowbreak{}272.14 万元 & 盘活存量资产、\allowbreak{}改善财务状况；出售存货收益 2，\allowbreak{}078.73 万元 \\
2024-12-12 & 向控股股东出售资产（关联交易） & 向控股股东大连盐化集团有限公司出售公司持有的原鲍鱼厂资产、\allowbreak{}原育苗三厂资产以及部分海域使用权等，\allowbreak{}215.36 万元 & 盘活存量资产；2024 年公司另收到船厂土地及附着物收购补偿款 5，\allowbreak{}860.00 万元 \\
2025-05-22 & 向特定对象发行 A 股股票（定增预案，对象为间接控股股东） & 拟向间接股东大连獐子岛海洋发展集团有限公司定向发行，\allowbreak{}发行价格 3.09 元/\allowbreak{}股 & 扣除相关发行费用后全部用于偿还银行借款及补充流动资金；经公司第八届董事会第十九次会议、\allowbreak{}第八届监事会第十四次会议审议通过；本次发行属于董事会确定全部发行对象的再融资；尚需取得有权国有资产监督管理部门或其授权主体批准、\allowbreak{}股东大会审议通过、\allowbreak{}深交所审核通过并经中国证监会同意注册；若发行落地，\allowbreak{}国资持股比例将升至 31.69\% \\
2026 & 银行融资置换（拟申请综合授信） & 拟向大连银行股份有限公司及大连农村商业银行股份有限公司申请融资贷款 4.95 亿元 & 用于等额置换公司在中国进出口银行辽宁省分行的贷款；公司间接控股股东大连獐子岛海洋发展集团有限公司提供相关支持；2026 年度公司另拟向中国农业银行大连市分行申请综合授信 4.50 亿元、\allowbreak{}向中国建设银行大连市分行申请 3.60 亿元、\allowbreak{}向中国工商银行大连市分行申请 2.90 亿元等 \\
\bottomrule
\end{tabular}
\end{adjustbox}
\end{center}

\pfl{历史融投资情况} 从现金流量表看，2025 年公司取得借款收到的现金 24.07 亿元、偿还债务支付的现金 24.34 亿元，筹资活动现金流净额 -0.97 亿元（2023 年为取得借款 19.84 亿元、偿还债务 19.97 亿元）。 2026 年上半年筹资活动现金流净额为 -0.32 亿元（取得借款 6.51 亿元、偿还债务 6.52 亿元，未年化），已转为净流出，处于净还债状态。 2025 年分配股利、利润或偿付利息支付的现金合计 0.64 亿元。 公开披露的融资记录共 7 笔，工具类型以 IPO、定向增发为主；金额与用途见下表。 其中规模最大的两笔为：2025 年 5 月，定向增发，拟向间接股东大连獐子岛海洋发展集团有限公司定向发行，发行价格 3.09 元/股；2006 年 9 月，IPO，实际发行 2，830.00 万股。 股本方面，近三年披露股本变动 3 次；总股本自 2021 年末以来未发生变化（7.11 亿股）。 分红方面，近三年未实施现金分红，在全部样本中属于分红能力偏弱的一端。 近三年“再融资”类公告 8 条，涉及定向增发、募集资金使用。 综合看，公司融资结构上股权与债务融资并举；2025 年筹资活动现金流净额 -0.97 亿元，融资节奏仍以维持存量债务周转为主，与前述经营与投资现金流的方向基本一致。

\pfl{未来潜在融资需求分析} 2026 年 6 月末货币资金 5.45 亿元，而短期债务 14.05 亿元（现金短债比 0.39 倍），2026 年上半年经营现金流净额 0.26 亿元。按“短期债务减去账面货币资金”的滚动偿付口径，静态缺口约 8.60 亿元；上半年盈利或接近盈亏平衡，该缺口不随经营亏损进一步扩大。 公司 2026 年 6 月末资产负债率 97.0\%、2026 年上半年 ROE 29.5\%（未年化），资产负债率偏高（97.0\%），融资需求首先用于置换与滚动存量债务、补充营运资金，属于“补血型”而非扩张型。 公开披露的后续资金安排线索包括：股权融资线索：2025 年 5 月，公司拟向间接股东大连獐子岛海洋发展集团有限公司定向发行不超过 1.69 亿股；债务与授信安排：2026 年度，公司拟向下列银行申请综合授信额度并根据银行要求提供相应增信措施：中国农业银行股份有限公司大连市分行 4.50 亿元、中国建设银行股份有限公司大连市分行 3.60 亿元、中国工商银行股份有限公司大连市分行 2.90 亿元、中国民生银行股份有限公司大连分行等。 综合看，公司的外部资金需求以补充流动性与项目投入为主，滚动偿付口径下的静态缺口约 8.60 亿元，融资节奏将受制于银行续贷意愿与股权融资窗口。




\subsubsection{佩蒂股份（300673）}
\label{co:300673}

\pfl{公司基本情况} 公司前身可追溯至 2013 年，浙江温州平阳起家的宠物食品（畜皮咬胶、植物咬胶、营养肉质零食、主粮）出口企业，2017 年 7 月创业板上市，是国内宠物食品行业最早“走出去”经营的企业之一，在越南（好嚼、巴啦啦、德信）、柬埔寨（爵味）、新西兰（北岛小镇、天然纯）建有生产基地，以 ODM 模式向欧美客户提供宠物零食供应链服务；2020 年定增募资 5.31 亿元、2021 年发行 7.2 亿元“佩蒂转债”，2026 年上半年营业收入 7.35 亿元、归母净利润 2,432.07 万元、同比下降 69.25\%。 公司于 2017-07-11 在深交所创业板首发上市，注册资本 2.47 亿元，总股本 2.49 亿股。 截至 2026 年 9 月 24 日收盘价 14.82 元，流通市值 23.8 亿元，近一年-14.4\%，流通市值在三级行业“宠物食品”的 3 家公司中按由大到小排第\zhnumber{3}位，近一年涨跌幅低于全样本中位数（-10.1\%）4.3 个百分点。 按 2026 年半年报披露的行业构成，收入占比前两位为宠物食品及其他 97.4\%；其他(补充) 2.6\%。

\pfl{历史股价} 以 2021 年 1 月为起点、2026 年 9 月为终点，股价由 31.15 元变为 14.82 元，累计 -52.4\%，区间最高 34.88 元（2021 年 5 月）、最低 11.08 元（2024 年 2 月）；当前价相当于区间最高价的 42.5\%，为区间最低价的 1.34 倍。 月度收益的年化波动率 39.1\%，高于全样本中位数 37.4\%，在三级行业“宠物食品”的 3 家公司中波动率由高到低排第\zhnumber{3}位。 把股价阶段与公司事件对齐看，2024 年 2 月 至 2025 年 11 月股价上涨+69.9\%，同期 2024 年末，陈振录持股 2，599.05 万股（10.51\%），无质押。 整体上看，区间的几处大级别拐点都能在公司的资本运作、股权变动或风险事件上找到对应，股价波动有明显的公司层面驱动，而非单纯的行业 beta。

\begin{center}
\begin{minipage}{\textwidth}\centering
\begin{tikzpicture}
\begin{axis}[
  width=12.6cm, height=3.9cm, scale only axis,
  xmin=2020.922, xmax=2026.888, ymin=10.30, ymax=37.32,
  xtick={2021,2022,2023,2024,2025,2026},
  yearticks,
  yticklabel style={font=\scriptsize},
  ylabel={元/股}, ylabel style={font=\scriptsize},
  grid=major, grid style={LightGray, line width=0.3pt},
  axis line style={InkGray, line width=0.5pt},
  every axis plot/.append style={line width=0.9pt},
]
\addplot[AgriGreen] table[x=x,y=y]{data/clean/profiles/px_300673.dat};
\addplot[SoftRed, dashed, line width=0.7pt] table[x=x,y=y]{data/clean/profiles/pxzz_300673.dat};
\addplot[only marks, mark=*, mark size=1.1pt, SoftRed] table[x=x,y=y]{data/clean/profiles/pxzz_300673.dat};
\end{axis}
\end{tikzpicture}

\captionof{figure}[佩蒂股份（300673）股价与周期骨架]{佩蒂股份（300673）月度收盘价（元/股）与 ZigZag 周期骨架（阈值 25\%，圆点为周期转折点）}
\end{minipage}
\end{center}

\pfl{过去周期分析} 以 25\% 的 ZigZag 阈值划分，样本区间内股价形成 2 个主升段与 3 个主跌段。 最大主升段出现在 2024 年 2 月 至 2025 年 11 月，21 个月+69.9\%；最大主跌段为 2021 年 1 月 至 2022 年 4 月，15 个月-51.0\%。 最近一次转折出现在 2026 年 7 月（下跌段自 2025 年 11 月起，8 个月-31.9\%），当前处于该段的延续之中。 公司股价较申万农林牧渔指数提前约 27 个月见底，是板块中较早定价周期反转的公司之一。 宠物食品属于消费型农业（第 \ref{sec:financing-strategy} 节归入“向消费端延伸”），收入增长由宠物数量与单宠消费驱动，与农产品价格周期的联动最弱，公司 2026 年上半年实现盈利、半年 ROE 1.4\%（未年化），好于全样本中位数（0.75\%），下行周期对其报表的冲击小于同业。 综合区间形态看，该股单段涨跌幅度偏大、以趋势段为主（最大主升 +69.9\%、最大主跌 -51.0\%），年化波动率 39.1\% 与全样本中位数 37.4\% 相比波动与样本中位数接近。

\pfl{盈利情况} 2026 年上半年营业收入 7.35 亿元（同比 +1.0\%），归母净利润 0.24 亿元，同比 -69.3\%。 以年报为趋势对照，2023---2025 年营业收入由 14.11 亿元变为 14.49 亿元（两年年均复合增速 +1.3\%），归母净利润由 -0.11 亿元变为 1.16 亿元。 按半年报口径，2026 年上半年的 ROE 为 1.4\%（未年化）、销售净利率 3.4\%、毛利率 28.4\%，ROE 在“宠物食品”3 家公司中按数值由高到低排第\zhnumber{3}位。 按同一口径，三级行业“宠物食品”3 家公司的半年 ROE 中位数为 4.00\%，该公司低于其中位数 2.64 个百分点。 面对这一局面，公司披露的举措包括：产品涉及高端主粮、宠物零食等多个品类，其中越南、柬埔寨工厂产能合计超 3 万吨；2026 年 5 月，公司因注销部分回购股份减少注册资本暨通知债权人，并完成部分回购股份注销暨股份变动。 综合看，公司在报告期维持盈利（高于全样本半年 ROE 中位数 0.75\%），利润的可持续性取决于宠物食品景气的延续与成本管控的执行。

\begin{center}\small\setlength{\tabcolsep}{3.4pt}
\captionof{table}[佩蒂股份（300673）关键财务指标]{佩蒂股份（300673）关键财务指标（合并报表口径；两个半年报期均未年化；现金短债比 $=$ 货币资金 $\div$ 短期债务）}
\begin{adjustbox}{max width=\textwidth}
\begin{tabular}{@{}lrrrrr@{}}
\toprule
指标 & 2023 年 & 2024 年 & 2025 年 & 2025 年半年报 & 2026 年半年报 \\
\midrule
营业收入（亿元） & 14.11 & 16.59 & 14.49 & 7.28 & 7.35 \\
归母净利润（亿元） & -0.11 & 1.82 & 1.16 & 0.79 & 0.24 \\
毛利率（\%） & 19.3 & 29.4 & 31.9 & 32.0 & 28.4 \\
ROE（\%） & -0.6 & 9.6 & 5.8 & 3.9 & 1.4 \\
资产负债率（\%） & 35.9 & 34.3 & 34.0 & 33.0 & 35.3 \\
经营活动现金流净额（亿元） & 1.78 & 3.66 & 1.91 & -0.31 & -0.66 \\
货币资金（亿元） & 8.16 & 6.75 & 7.50 & 5.19 & 4.95 \\
有息负债合计（亿元） & 8.40 & 7.89 & 7.99 & 7.82 & 8.21 \\
现金短债比（倍） & 9.63 & 13.12 & 51.11 & 45.14 & 34.24 \\
\bottomrule
\end{tabular}
\end{adjustbox}
\end{center}

\begin{center}
\begin{minipage}{\textwidth}\centering
\begin{tikzpicture}
\begin{groupplot}[
  group style={group size=2 by 1, horizontal sep=0.60cm},
  width=6.85cm, height=4.60cm,
  tick label style={font=\tiny},
  title style={font=\scriptsize\heiti},
  yticklabel style={font=\tiny},
  ylabel style={font=\tiny},
  grid=major, grid style={LightGray, line width=0.3pt},
  axis line style={InkGray, line width=0.5pt},
  enlarge x limits={abs=0.8},
  legend style={font=\scriptsize, at={(0.5,-0.20)}, anchor=north,
                legend columns=2, draw=none, fill=none, column sep=6pt},
]
\nextgroupplot[
  ybar, bar width=4pt,
  ymin=-1.85, ymax=19.41,
  xtick={0,1,2,3,4,5.7},
  xticklabels={2021,2022,2023,2024,2025,2026H1},
  ylabel={亿元},
]
\addplot[fill=AgriGreen!75, draw=AgriGreen, bar shift=-2.6pt] table[x=i, y=rev]{data/clean/profiles/fin_300673.dat};
\addplot[fill=SoftRed!60, draw=SoftRed, bar shift=2.6pt] table[x=i, y=ni]{data/clean/profiles/fin_300673.dat};
\addplot[fill=AgriGreen!30, draw=AgriGreen, pattern=north east lines, bar shift=-2.6pt] table[x=x, y=rev]{data/clean/profiles/finh_300673.dat};
\addplot[fill=SoftRed!20, draw=SoftRed, pattern=north east lines, bar shift=2.6pt] table[x=x, y=ni]{data/clean/profiles/finh_300673.dat};
\legend{营业收入（年/半年）, 归母净利润（年/半年）}
\nextgroupplot[
  ymin=-3.6, ymax=40.9,
  xtick={0,1,2,3,4,5.7},
  xticklabels={2021,2022,2023,2024,2025,2026H1},
  ylabel={\%},
]
\addplot[AgriGold, mark=*, mark size=1.3pt] table[x=i, y=roe]{data/clean/profiles/fin_300673.dat};
\addplot[OilBlue, mark=square*, mark size=1.3pt] table[x=i, y=debt]{data/clean/profiles/fin_300673.dat};
\addplot[only marks, AgriGold, mark=*, mark size=2.0pt] table[x=x, y=roe]{data/clean/profiles/finh_300673.dat};
\addplot[only marks, OilBlue, mark=square*, mark size=2.0pt] table[x=x, y=debt]{data/clean/profiles/finh_300673.dat};
\legend{ROE, 资产负债率}
\end{groupplot}
\end{tikzpicture}
\begin{tikzpicture}
\begin{axis}[
  name=finbot,
  width=13.75cm, height=3.10cm, scale only axis,
  ymin=-1.03, ymax=11.53,
  xtick={0,1,2,3,4,5.7},
  xticklabels={2021,2022,2023,2024,2025,2026H1},
  tick label style={font=\tiny},
  yticklabel style={font=\tiny},
  ylabel={亿元}, ylabel style={font=\tiny},
  ybar, bar width=4pt,
  grid=major, grid style={LightGray, line width=0.3pt},
  axis line style={InkGray, line width=0.5pt},
  enlarge x limits={abs=0.8},
  legend style={font=\scriptsize, at={(0.5,-0.26)}, anchor=north,
                legend columns=2, draw=none, fill=none, column sep=10pt},
]
\addplot[fill=AgriGreen!55, draw=AgriGreen, bar shift=-3.0pt] table[x=i, y=cash]{data/clean/profiles/fin_300673.dat};
\addplot[fill=SoftRed!45, draw=SoftRed, bar shift=3.0pt] table[x=i, y=ibd]{data/clean/profiles/fin_300673.dat};
\addplot[fill=AgriGreen!25, draw=AgriGreen, pattern=north east lines, bar shift=-3.0pt] table[x=x, y=cash]{data/clean/profiles/finh_300673.dat};
\addplot[fill=SoftRed!20, draw=SoftRed, pattern=north east lines, bar shift=3.0pt] table[x=x, y=ibd]{data/clean/profiles/finh_300673.dat};
\legend{货币资金（左轴）, 有息负债（左轴）}
\end{axis}
\begin{axis}[
  at={(finbot.south east)}, anchor=south east,
  width=13.75cm, height=3.10cm, scale only axis,
  axis y line*=right, axis x line*=none, xtick=\empty,
  ymin=-1.79, ymax=4.96,
  tick label style={font=\tiny},
  yticklabel style={font=\tiny}, scaled y ticks=false,
  legend style={font=\scriptsize, at={(0.5,-0.44)}, anchor=north,
                legend columns=1, draw=none, fill=none},
]
\addplot[OilBlue, mark=*, mark size=2.2pt, line width=1.3pt] table[x=i, y=ocf]{data/clean/profiles/fin_300673.dat};
\addplot[only marks, OilBlue, mark=*, mark size=3.2pt] table[x=x, y=ocf]{data/clean/profiles/finh_300673.dat};
\legend{经营活动现金流净额（右轴，亿元，蓝点）}
\end{axis}
\end{tikzpicture}

\captionof{figure}[佩蒂股份（300673）近年主要财务指标]{佩蒂股份（300673）近年主要财务指标（左上：营业收入与归母净利润；右上：ROE 与资产负债率；下：货币资金、有息负债为左轴柱状，经营活动现金流净额为其右轴的蓝点折线。
2021---2025 年为年报口径，2026H1 为 2026 年半年报/半年末，与其前一年报之间留出间隔、以斜纹或空心点区分；半年报数据未年化）}
\end{minipage}
\end{center}

\pfl{财务分析} 2026 年 6 月末资产负债率 35.3\%，较 2025 年末上升 1.3 个百分点（样本中位数 52.2\%，所处三级行业内按由低到高排第\zhnumber{2}位）。 2026 年 6 月末货币资金 4.95 亿元、短期债务（短期借款与一年内到期非流动负债）0.14 亿元、长期债务 8.06 亿元，有息负债合计 8.21 亿元，现金短债比 34.24 倍——覆盖充分。 净有息负债 3.26 亿元，相当于其 2026 年上半年营业收入的 44\%。 流动比率 8.13、速动比率 6.20，短期偿债能力处于正常区间。 2026 年上半年经营活动现金流净额 -0.66 亿元，净利润为正但经营现金流为负，利润的现金含量偏低。 2026 年 6 月末在建工程 0.79 亿元，较上年末增加 0.12 亿元，仍有在建投入。 综合看，现金短债比 34.24 倍，短期偿付压力可控；资产负债率低于样本中位数 16.9 个百分点；有息负债中长期债务占比更高，期限结构对短债滚动的压力较小。 公司披露的股东与质押风险为：2026 年 6 月末，股东陈林艺持股 539.80 万股（2.18\%），其中质押 385.55 万股；2025 年末，陈振录持股 2，599.05 万股（10.51\%），无质押。

\pfl{重大战略转型} 公司的战略动作可归为以下几类：并购与重组（3 项）：2019 年 6 月，公司变更 6，500.00 万元募集资金用途并增资新西兰北岛小镇宠物食品有限公司 1；2019 年 2 月，公司完成对德信皮业（越南）有限公司 100\% 股权收购并改造成越南生产基地；海外与国际化（2 项）：2013 年，公司最早于 2013 年在越南设立第一家海外生产基地。 公告口径上，近三年（2023 年 9 月---2026 年 9 月）公司披露重大重组与并购类公告 0 条、对外投资与转型类公告 5 条、回购与股权激励类公告 60 条，合计公告 426 条。 第一大行业仍是“宠物食品及其他”，收入占比 97.4\%，与 2021 年年报相比没有方向性变化。 综合判断，报告期内公司有 5 条与战略相关的公告落地（重大重组与并购 0 条、对外投资与转型 5 条），资本运作与主业调整之间存在直接关联。

\begin{center}\small\setlength{\tabcolsep}{4pt}
\captionof{table}[佩蒂股份（300673）历史融资与资本运作]{佩蒂股份（300673）历史融资与资本运作（据公司公告与公开报道整理；金额为披露值，含首次公开发行、再融资、债券、银行借款与重整投资等）}
\begin{adjustbox}{max width=\textwidth}
\begin{tabular}{@{}p{1.45cm}p{2.55cm}p{4.05cm}p{4.95cm}@{}}
\toprule
时间 & 方式 & 规模 & 用途与说明 \\
\midrule
2017-07-11 & IPO（深圳证券交易所创业板） & 实际发行 2，\allowbreak{}000.00 万股 & 份有限公司，\allowbreak{}发行市盈率（按发行后总股本）22.99 倍 \\
2019-02-26 & 境外现金收购（越南生产基地） & 收购德信皮业（越南）有限公司 100\% 股权，\allowbreak{}股权取得成本 2，\allowbreak{}384.53 万元 & 扩充海外产能 \\
2019-06 & 境外投资设立子公司 & 通过子公司新西兰北岛小镇宠物食品有限公司出资 500.00 万美元设立新西兰天然纯宠物食品有限公司 & 实施“新西兰年产 4 万吨高品质宠物干粮新建项目”；该公司于 2019 年 6 月 26 日完成设立 \\
2019-06-11 & 变更部分募集资金用途并向全资子公司增资（用于境外收购） & 变更募集资金 6，\allowbreak{}500.00 万元人民币并加自有资金向新西兰北岛小镇宠物食品有限公司增资 1 & 收购 BOP Industries Limited 已发行的 77.02\% 股份（实缴）；变动前募投项目为江苏康贝宠物食品有限公司实施的“年产 2 \\
2021-12-28 & 向不特定对象发行可转换公司债券（佩蒂转债） & 发行 720 万张、\allowbreak{}每张面值 100.00 元，\allowbreak{}期限 6 年，\allowbreak{}募集资金总额 7.20 亿元 & 经深交所审核同意并经中国证监会同意佩蒂动物营养科技股份有限；债券代码 123133 \\
\bottomrule
\end{tabular}
\end{adjustbox}
\end{center}

\pfl{历史融投资情况} 从现金流量表看，2025 年公司取得借款收到的现金 0.13 亿元、偿还债务支付的现金 0.45 亿元，筹资活动现金流净额 -1.83 亿元（2023 年为取得借款 1.60 亿元、偿还债务 1.38 亿元）。 2026 年上半年筹资活动现金流净额为 -0.69 亿元（取得借款 0.10 亿元、偿还债务 0.07 亿元，未年化），已转为净流出，处于净还债状态。 2025 年分配股利、利润或偿付利息支付的现金合计 0.98 亿元。 公开披露的融资记录共 5 笔，工具类型以 IPO、可转债为主；金额与用途见下表。 其中规模最大的两笔为：2019 年 6 月，变更部分募集资金用途，变更募集资金 6，500.00 万元人民币并加自有资金向新西兰北岛小镇宠物食品有限公司增资 1；2017 年 7 月，IPO，实际发行 2，000.00 万股。 股本方面，近三年披露股本变动 17 次（可转债转股 10 次、其他 4 次、股份回购 2 次、股权激励 1 次）；总股本由 2021 年末的 2.53 亿股变为 2.49 亿股（-1.81\%）。 分红方面，近三年现金分红 2 次、合计每股派现 0.60 元，最近一次实施在 2025 年 12 月。 近三年“再融资”类公告 18 条，涉及可转债。 综合看，公司融资结构上股权与债务融资并举；2025 年筹资活动现金流净额 -1.83 亿元，融资节奏仍以维持存量债务周转为主，与前述经营与投资现金流的方向基本一致。

\pfl{未来潜在融资需求分析} 2026 年 6 月末货币资金 4.95 亿元，而短期债务 0.14 亿元（现金短债比 34.24 倍），2026 年上半年经营现金流净额 -0.66 亿元。账面货币资金可以覆盖短期债务，缺口不体现在静态偿付层面。 公司 2026 年 6 月末资产负债率 35.3\%、2026 年上半年 ROE 1.4\%（未年化），账面杠杆不高、盈利能力一般，短期内没有必须依靠外部融资才能维持的经营缺口；后续融资更可能指向技改、扩产或产业并购。 公开披露的后续资金安排线索包括：股权融资线索：2026 年 6 月，公司 2026 年半年度募集资金存放与使用情况显示：2021 年可转债募集资金累计投入 3.87 亿元。 资金需求还要叠加在建工程：期末在建工程 0.79 亿元、2026 年上半年购建固定资产等支出 0.41 亿元（未年化），这部分支出在周期底部通常会被主动放缓。 综合看，公司在静态偿付层面不存在缺口，外部融资更多是配合扩张节奏与并购整合的主动性安排，而非偿债压力驱动。




\subsubsection{开创国际（600097）}
\label{co:600097}

\pfl{公司基本情况} 公司前身可追溯至 1993 年，上海国资（光明食品集团）控制的远洋渔业上市平台。前身为 1997 年上市的海南恒泰芒果，2008 年 12 月经资产置换重组为以远洋渔业为主业的开创国际；旗下拥有 12 艘大型金枪鱼围网船和 2 艘大型拖网加工船，并控股西班牙 ALBO 等境外水产企业。2025 年归母净利润 0.74 亿元，2026 年上半年因渔场资源波动、基里巴斯部分水域关闭与国际油价上涨转为亏损。 公司于 1997-06-19 在上交所首发上市，注册资本 2.41 亿元，总股本 2.41 亿股。 截至 2026 年 9 月 24 日收盘价 9.20 元，流通市值 22.2 亿元，近一年-18.0\%，流通市值在三级行业“海洋捕捞”的 2 家公司中按由大到小排第\zhnumber{2}位，近一年涨跌幅低于全样本中位数（-10.1\%）7.9 个百分点。 按 2026 年半年报披露的行业构成，收入占比前三位为食品加工 52.0\%；远洋捕捞 33.3\%；渔货贸易 14.5\%。 与 2021 年年报相比，第一大行业口径由“海洋捕捞”（56.7\%）变为“食品加工”（52.0\%）；两期披露的分类口径有调整，占比差异中同时包含分类因素。

\pfl{历史股价} 本报告样本区间（2021 年 1 月 至 2026 年 9 月）内，股价由 9.23 元变为 9.20 元，累计 -0.3\%，区间最高 12.93 元（2022 年 3 月）、最低 7.62 元（2024 年 8 月）；当前价相当于区间最高价的 71.2\%，为区间最低价的 1.21 倍。 月度收益的年化波动率 34.1\%，低于全样本中位数 37.4\%，在三级行业“海洋捕捞”的 2 家公司中波动率由高到低排第\zhnumber{2}位。 把股价阶段与公司事件对齐看，2024 年 8 月 至 2026 年 2 月股价上涨+63.5\%，同期（2025 年 12 月）公司 2017 年完成的非公开发行为上市以来可查证的最后一次股权融资。 整体上看，区间的几处大级别拐点都能在公司的资本运作、股权变动或风险事件上找到对应，股价波动有明显的公司层面驱动，而非单纯的行业 beta。

\begin{center}
\begin{minipage}{\textwidth}\centering
\begin{tikzpicture}
\begin{axis}[
  width=12.6cm, height=3.9cm, scale only axis,
  xmin=2020.922, xmax=2026.888, ymin=7.09, ymax=13.84,
  xtick={2021,2022,2023,2024,2025,2026},
  yearticks,
  yticklabel style={font=\scriptsize},
  ylabel={元/股}, ylabel style={font=\scriptsize},
  grid=major, grid style={LightGray, line width=0.3pt},
  axis line style={InkGray, line width=0.5pt},
  every axis plot/.append style={line width=0.9pt},
]
\addplot[AgriGreen] table[x=x,y=y]{data/clean/profiles/px_600097.dat};
\addplot[SoftRed, dashed, line width=0.7pt] table[x=x,y=y]{data/clean/profiles/pxzz_600097.dat};
\addplot[only marks, mark=*, mark size=1.1pt, SoftRed] table[x=x,y=y]{data/clean/profiles/pxzz_600097.dat};
\end{axis}
\end{tikzpicture}

\captionof{figure}[开创国际（600097）股价与周期骨架]{开创国际（600097）月度收盘价（元/股）与 ZigZag 周期骨架（阈值 25\%，圆点为周期转折点）}
\end{minipage}
\end{center}

\pfl{过去周期分析} 以 25\% 的 ZigZag 阈值划分，样本区间内股价形成 3 个主升段与 4 个主跌段。 最大主升段出现在 2024 年 8 月 至 2026 年 2 月，18 个月+63.5\%；最大主跌段为 2026 年 2 月 至 2026 年 6 月，4 个月-37.6\%。 最近一次转折出现在 2026 年 6 月（下跌段自 2026 年 2 月起，4 个月-37.6\%），当前处于该段的延续之中。 公司股价较申万农林牧渔指数提前约 21 个月见底，是板块中较早定价周期反转的公司之一。 远洋与近海捕捞的产量受资源与配额约束，成本端则跟随燃油与人工，公司 2026 年 6 月末资产负债率 34.4\%、半年 ROE -1.3\%（未年化），在本轮下行中属于随行业同步调整的一类。 综合区间形态看，该股单段涨跌幅度偏大、以趋势段为主（最大主升 +63.5\%、最大主跌 -37.6\%），年化波动率 34.1\% 与全样本中位数 37.4\% 相比波动与样本中位数接近。

\pfl{盈利情况} 2026 年上半年营业收入 11.82 亿元（同比 -5.4\%），归母净利润 -0.30 亿元，由上年同期的盈利转为亏损。 以年报为趋势对照，2023---2025 年营业收入由 20.15 亿元变为 23.37 亿元（两年年均复合增速 +7.7\%），归母净利润由 1.48 亿元变为 0.74 亿元。 按半年报口径，2026 年上半年的 ROE 为 -1.3\%（未年化）、销售净利率 -2.5\%、毛利率 23.1\%，ROE 在“海洋捕捞”2 家公司中按数值由高到低排第\zhnumber{2}位。 按同一口径，三级行业“海洋捕捞”2 家公司的半年 ROE 中位数为 9.79\%，该公司低于其中位数 11.12 个百分点。 从公司披露的原因看，业绩变动主要来自：2025 年，第三季度毛利指标出现异常的核心原因是公司于当期收到上海市农业农村委员会拨付的国际履约能力提升政府补助约 5，353.86 万元；2026 年上半年，国际油价持续上涨、安排部分船只提前回国修理，同时燃油价格上涨导致船只生产成本增加。 综合看，公司仍处亏损区间、由盈转亏，半年 ROE 在三级行业内按由高到低排第\zhnumber{2}位，单位盈利能力的修复依赖行业景气与自身成本端的改善，报表弹性主要体现为价格与销量的方向性变化。

\begin{center}\small\setlength{\tabcolsep}{3.4pt}
\captionof{table}[开创国际（600097）关键财务指标]{开创国际（600097）关键财务指标（合并报表口径；两个半年报期均未年化；现金短债比 $=$ 货币资金 $\div$ 短期债务）}
\begin{adjustbox}{max width=\textwidth}
\begin{tabular}{@{}lrrrrr@{}}
\toprule
指标 & 2023 年 & 2024 年 & 2025 年 & 2025 年半年报 & 2026 年半年报 \\
\midrule
营业收入（亿元） & 20.15 & 23.16 & 23.37 & 12.50 & 11.82 \\
归母净利润（亿元） & 1.48 & 0.61 & 0.74 & 0.28 & -0.30 \\
毛利率（\%） & 32.5 & 28.1 & 29.7 & 31.5 & 23.1 \\
ROE（\%） & 6.7 & 2.7 & 3.2 & 1.2 & -1.3 \\
资产负债率（\%） & 32.6 & 32.3 & 31.5 & 30.0 & 34.4 \\
经营活动现金流净额（亿元） & 1.07 & 3.35 & 1.68 & 1.26 & 0.48 \\
货币资金（亿元） & 5.19 & 5.25 & 4.28 & 5.00 & 3.90 \\
有息负债合计（亿元） & 6.22 & 5.27 & 4.86 & 5.02 & 4.79 \\
现金短债比（倍） & 4.49 & 4.39 & 2.16 & 3.28 & 2.22 \\
\bottomrule
\end{tabular}
\end{adjustbox}
\end{center}

\begin{center}
\begin{minipage}{\textwidth}\centering
\begin{tikzpicture}
\begin{groupplot}[
  group style={group size=2 by 1, horizontal sep=0.60cm},
  width=6.85cm, height=4.60cm,
  tick label style={font=\tiny},
  title style={font=\scriptsize\heiti},
  yticklabel style={font=\tiny},
  ylabel style={font=\tiny},
  grid=major, grid style={LightGray, line width=0.3pt},
  axis line style={InkGray, line width=0.5pt},
  enlarge x limits={abs=0.8},
  legend style={font=\scriptsize, at={(0.5,-0.20)}, anchor=north,
                legend columns=2, draw=none, fill=none, column sep=6pt},
]
\nextgroupplot[
  ybar, bar width=4pt,
  ymin=-2.67, ymax=26.21,
  xtick={0,1,2,3,4,5.7},
  xticklabels={2021,2022,2023,2024,2025,2026H1},
  ylabel={亿元},
]
\addplot[fill=AgriGreen!75, draw=AgriGreen, bar shift=-2.6pt] table[x=i, y=rev]{data/clean/profiles/fin_600097.dat};
\addplot[fill=SoftRed!60, draw=SoftRed, bar shift=2.6pt] table[x=i, y=ni]{data/clean/profiles/fin_600097.dat};
\addplot[fill=AgriGreen!30, draw=AgriGreen, pattern=north east lines, bar shift=-2.6pt] table[x=x, y=rev]{data/clean/profiles/finh_600097.dat};
\addplot[fill=SoftRed!20, draw=SoftRed, pattern=north east lines, bar shift=2.6pt] table[x=x, y=ni]{data/clean/profiles/finh_600097.dat};
\legend{营业收入（年/半年）, 归母净利润（年/半年）}
\nextgroupplot[
  ymin=-4.5, ymax=42.0,
  xtick={0,1,2,3,4,5.7},
  xticklabels={2021,2022,2023,2024,2025,2026H1},
  ylabel={\%},
]
\addplot[AgriGold, mark=*, mark size=1.3pt] table[x=i, y=roe]{data/clean/profiles/fin_600097.dat};
\addplot[OilBlue, mark=square*, mark size=1.3pt] table[x=i, y=debt]{data/clean/profiles/fin_600097.dat};
\addplot[only marks, AgriGold, mark=*, mark size=2.0pt] table[x=x, y=roe]{data/clean/profiles/finh_600097.dat};
\addplot[only marks, OilBlue, mark=square*, mark size=2.0pt] table[x=x, y=debt]{data/clean/profiles/finh_600097.dat};
\legend{ROE, 资产负债率}
\end{groupplot}
\end{tikzpicture}
\begin{tikzpicture}
\begin{axis}[
  name=finbot,
  width=13.75cm, height=3.10cm, scale only axis,
  ymin=-0.76, ymax=8.49,
  xtick={0,1,2,3,4,5.7},
  xticklabels={2021,2022,2023,2024,2025,2026H1},
  tick label style={font=\tiny},
  yticklabel style={font=\tiny},
  ylabel={亿元}, ylabel style={font=\tiny},
  ybar, bar width=4pt,
  grid=major, grid style={LightGray, line width=0.3pt},
  axis line style={InkGray, line width=0.5pt},
  enlarge x limits={abs=0.8},
  legend style={font=\scriptsize, at={(0.5,-0.26)}, anchor=north,
                legend columns=2, draw=none, fill=none, column sep=10pt},
]
\addplot[fill=AgriGreen!55, draw=AgriGreen, bar shift=-3.0pt] table[x=i, y=cash]{data/clean/profiles/fin_600097.dat};
\addplot[fill=SoftRed!45, draw=SoftRed, bar shift=3.0pt] table[x=i, y=ibd]{data/clean/profiles/fin_600097.dat};
\addplot[fill=AgriGreen!25, draw=AgriGreen, pattern=north east lines, bar shift=-3.0pt] table[x=x, y=cash]{data/clean/profiles/finh_600097.dat};
\addplot[fill=SoftRed!20, draw=SoftRed, pattern=north east lines, bar shift=3.0pt] table[x=x, y=ibd]{data/clean/profiles/finh_600097.dat};
\legend{货币资金（左轴）, 有息负债（左轴）}
\end{axis}
\begin{axis}[
  at={(finbot.south east)}, anchor=south east,
  width=13.75cm, height=3.10cm, scale only axis,
  axis y line*=right, axis x line*=none, xtick=\empty,
  ymin=-0.87, ymax=4.36,
  tick label style={font=\tiny},
  yticklabel style={font=\tiny}, scaled y ticks=false,
  legend style={font=\scriptsize, at={(0.5,-0.44)}, anchor=north,
                legend columns=1, draw=none, fill=none},
]
\addplot[OilBlue, mark=*, mark size=2.2pt, line width=1.3pt] table[x=i, y=ocf]{data/clean/profiles/fin_600097.dat};
\addplot[only marks, OilBlue, mark=*, mark size=3.2pt] table[x=x, y=ocf]{data/clean/profiles/finh_600097.dat};
\legend{经营活动现金流净额（右轴，亿元，蓝点）}
\end{axis}
\end{tikzpicture}

\captionof{figure}[开创国际（600097）近年主要财务指标]{开创国际（600097）近年主要财务指标（左上：营业收入与归母净利润；右上：ROE 与资产负债率；下：货币资金、有息负债为左轴柱状，经营活动现金流净额为其右轴的蓝点折线。
2021---2025 年为年报口径，2026H1 为 2026 年半年报/半年末，与其前一年报之间留出间隔、以斜纹或空心点区分；半年报数据未年化）}
\end{minipage}
\end{center}

\pfl{财务分析} 2026 年 6 月末资产负债率 34.4\%，较 2025 年末上升 2.9 个百分点（样本中位数 52.2\%，所处三级行业内按由低到高排第\zhnumber{1}位）。 2026 年 6 月末货币资金 3.90 亿元、短期债务（短期借款与一年内到期非流动负债）1.76 亿元、长期债务 3.03 亿元，有息负债合计 4.79 亿元，现金短债比 2.22 倍——覆盖充分。 净有息负债 0.89 亿元，相当于其 2026 年上半年营业收入的 8\%。 流动比率 2.05、速动比率 1.09，短期流动性无明显缺口。 2026 年上半年经营活动现金流净额 0.48 亿元，净利润为负而经营现金流为正，差额主要来自折旧摊销与营运资本变动，说明亏损尚未完全转化为现金流出。 2026 年 6 月末在建工程 1.49 亿元，较上年末增加 0.74 亿元，仍有在建投入。 2026 年 6 月末商誉 1.04 亿元，占归母净资产的 4.7\%，是净资产中最脆弱的一块。 综合看，现金短债比 2.22 倍，短期偿付压力可控；资产负债率低于样本中位数 17.8 个百分点；有息负债中长期债务占比更高，期限结构对短债滚动的压力较小。 公司披露的监管与合规风险为：2026 年 7 月，公司就上海证券交易所对 2025 年年度报告及 2026 年一季度报告的信息披露监管问询函作出回复（致同会计师事务所出具相关回复）。

\pfl{重大战略转型} 公司的战略动作可归为以下几类：并购与重组（3 项）：2017 年，公司以非公开发行募集资金（总额调整至 6.00 亿元）用于收购西班牙 ALBO 公司股权及补充流动资金；2016 年 6 月，S.L. 100\% 股权的现金收购交割。 公告口径上，近三年（2023 年 9 月---2026 年 9 月）公司披露重大重组与并购类公告 2 条、对外投资与转型类公告 1 条、回购与股权激励类公告 0 条，合计公告 260 条。 主营结构的口径变化已经落到报表上：2021 年年报的第一大行业为“海洋捕捞”（56.7\%），2026 年半年报为“食品加工”（52.0\%）；公司在这两期的分部划分方式有过调整。 综合判断，报告期内公司有 3 条与战略相关的公告落地（重大重组与并购 2 条、对外投资与转型 1 条），资本运作与主业调整之间存在直接关联。

\begin{center}\small\setlength{\tabcolsep}{4pt}
\captionof{table}[开创国际（600097）历史融资与资本运作]{开创国际（600097）历史融资与资本运作（据公司公告与公开报道整理；金额为披露值，含首次公开发行、再融资、债券、银行借款与重整投资等）}
\begin{adjustbox}{max width=\textwidth}
\begin{tabular}{@{}p{1.45cm}p{2.55cm}p{4.05cm}p{4.95cm}@{}}
\toprule
时间 & 方式 & 规模 & 用途与说明 \\
\midrule
1997-05-21 & IPO（上海证券交易所） & 首次向社会公众发行人民币普通股 3，\allowbreak{}500 万股 & 经海南省证券管理办公室（琼证办（1997）92 号）文件与中国证券监督管理委员会（证监发审字（1997）134、\allowbreak{}235 号）文件批准 \\
2015-09 & 非公开发行股票（筹划阶段，规模后调整） & 初步方案拟募集资金约 11.00 亿元 & 收购控股股东上海远洋渔业持有的斐济金洋渔业有限公司 20\% 股权并对其增资至持有其 59.375\% 股权；收购远洋渔业持有的基里巴斯渔业公司 20\% 股权；公司股票经四次延期复牌 \\
2016-04 & 非公开发行股票（预案修订稿，拟募集资金总额不超过 84,000 万元） & 拟募集资金总额不超过 8.40 亿元（含发行费用） & 收购西班牙 ALBO 公司 100\% 股权（拟投入募集资金 4.26 亿元）、\allowbreak{}新建舟山金枪鱼食品加工基地项目（拟投入 2.00 亿元）、\allowbreak{}补充流动资金（拟投入 2.14 亿元）；ALBO 公司 100\% 股权购买价款为 60 \\
2016-05 & 重大资产购买（现金收购，募集资金置换） & 以下属全资子公司开创远洋为收购主体，\allowbreak{}S.L. 100\% 股权 & 公司拟以非公开发行募集资金 4.26 亿元投入本次收购 \\
2016-12-15 & 非公开发行股票（反馈与承诺披露） & 控股股东上海远洋渔业有限公司承诺认购不低于本次非公开发行最终确定发行股票总数量的 20\%（含 20\%） & 不适用（承诺事项）；本次非公开发行的董事会决议日为 2015 年 11 月 24 日；上海远洋及其控制的主体上海宇洋人力资源有限公司、\allowbreak{}上海蒂尔远洋渔业有限公司、\allowbreak{}上海金优远洋渔业有限公司、\allowbreak{}上海水锦洋食品有限公司、\allowbreak{}海南双海海洋渔业有限公司等一致行动人承诺自定价基准日前六个月至本次发行完成后六个月不存在主动减持计划 \\
2017-12-07 & 非公开发行股票（完成） & 发行人民币普通股 3，\allowbreak{}833.87 万股，\allowbreak{}发行价 15.65 元/\allowbreak{}股 & 收购西班牙 ALBO 公司 100\% 股权及补充流动资金（募集资金总额后经削减调整至 6.00 亿元）；经公司第七届董事会第十一次（临时）会议、\allowbreak{}第十三次（临时）会议、\allowbreak{}第十八次（临时）会议及 2016 年第一次临时股东会决议 \\
\bottomrule
\end{tabular}
\end{adjustbox}
\end{center}

\pfl{历史融投资情况} 从现金流量表看，2025 年公司取得借款收到的现金 0.24 亿元、偿还债务支付的现金 0.16 亿元，筹资活动现金流净额 -0.86 亿元（2023 年为偿还债务 0.15 亿元）。 2026 年上半年筹资活动现金流净额为 0.40 亿元（取得借款 1.23 亿元、偿还债务 0.10 亿元，未年化），仍处于净融入状态。 2025 年分配股利、利润或偿付利息支付的现金合计 0.22 亿元。 公开披露的融资记录共 6 笔，工具类型以 IPO、定向增发为主；金额与用途见下表。 其中规模最大的两笔为：2016 年 4 月，定向增发，拟募集资金总额不超过 8.40 亿元（含发行费用）；2015 年 9 月，定向增发，初步方案拟募集资金约 11.00 亿元。 股本方面，近三年披露股本变动 3 次；总股本自 2021 年末以来未发生变化（2.41 亿股）。 分红方面，近三年现金分红 4 次、合计每股派现 0.50 元，最近一次实施在 2025 年 12 月。 综合看，公司融资结构上股权与债务融资并举；2025 年筹资活动现金流净额 -0.86 亿元，融资节奏仍以维持存量债务周转为主，与前述经营与投资现金流的方向基本一致。

\pfl{未来潜在融资需求分析} 2026 年 6 月末货币资金 3.90 亿元，而短期债务 1.76 亿元（现金短债比 2.22 倍），2026 年上半年经营现金流净额 0.48 亿元。账面货币资金可以覆盖短期债务，缺口不体现在静态偿付层面。 公司 2026 年上半年归母净利润为负（-0.30 亿元），但 6 月末资产负债率 34.4\% 尚未进入第一档（亏损且负债率 $>$65\%）的区间，当前融资需求以补充流动资金、支撑低谷期经营性支出为主。 公开披露的后续资金安排线索包括：股权融资线索：2025 年 12 月，公司 2017 年完成的非公开发行为上市以来可查证的最后一次股权融资。 资金需求还要叠加在建工程：期末在建工程 1.49 亿元、2026 年上半年购建固定资产等支出 1.20 亿元（未年化），这部分支出在周期底部通常会被主动放缓。 综合看，公司在静态偿付层面不存在缺口，外部融资更多是配合扩张节奏与并购整合的主动性安排，而非偿债压力驱动。




\subsubsection{永安林业（000663）}
\label{co:000663}

\pfl{公司基本情况} 1994 年设立至今，永安林业是全国首家林业综合性股份制公司，1996 年 12 月深交所上市，主营森林资源培育与林木采伐经营、人造板与地板家居（林板一体化）；2021 年 2 月中国林业集团完成战略重组、成为间接控股股东。2023—2025 年归母净利润由盈转亏并连续两年亏损，生物质能循环利用项目停建并计提减值、引发诉讼，中林集团避免同业竞争的五年承诺已逾期未履行。 公司于 1996-12-06 在深交所主板首发上市，注册资本 3.36 亿元，总股本 3.37 亿股。 截至 2026 年 9 月 24 日收盘价 6.08 元，流通市值 18.6 亿元，近一年-5.7\%，流通市值在三级行业“林业Ⅲ”的 4 家公司中按由大到小排第\zhnumber{4}位，近一年涨跌幅高于全样本中位数（-10.1\%）4.3 个百分点。 按 2026 年半年报披露的行业构成，收入占比前三位为人造板制造业 93.1\%；其他(补充) 4.2\%；林业 2.7\%。

\pfl{历史股价} 以 2021 年 1 月为起点、2026 年 9 月为终点，股价由 4.25 元变为 6.08 元，累计 +43.1\%，区间最高 12.18 元（2023 年 6 月）、最低 4.19 元（2024 年 6 月）；当前价相当于区间最高价的 49.9\%，为区间最低价的 1.45 倍。 月度收益的年化波动率 51.7\%，高于全样本中位数 37.4\%，在三级行业“林业Ⅲ”的 4 家公司中波动率由高到低排第\zhnumber{3}位。 把股价阶段与公司事件对齐看，2024 年 6 月 至 2025 年 11 月股价上涨+108.1\%，同期（2025 年 10 月）公司全资孙公司向永安市人民法院提起对必奥新能源科技有限公司的买卖合同纠纷诉讼。 整体上看，区间的几处大级别拐点都能在公司的资本运作、股权变动或风险事件上找到对应，股价波动有明显的公司层面驱动，而非单纯的行业 beta。

\begin{center}
\begin{minipage}{\textwidth}\centering
\begin{tikzpicture}
\begin{axis}[
  width=12.6cm, height=3.9cm, scale only axis,
  xmin=2020.922, xmax=2026.888, ymin=3.90, ymax=13.03,
  xtick={2021,2022,2023,2024,2025,2026},
  yearticks,
  yticklabel style={font=\scriptsize},
  ylabel={元/股}, ylabel style={font=\scriptsize},
  grid=major, grid style={LightGray, line width=0.3pt},
  axis line style={InkGray, line width=0.5pt},
  every axis plot/.append style={line width=0.9pt},
]
\addplot[AgriGreen] table[x=x,y=y]{data/clean/profiles/px_000663.dat};
\addplot[SoftRed, dashed, line width=0.7pt] table[x=x,y=y]{data/clean/profiles/pxzz_000663.dat};
\addplot[only marks, mark=*, mark size=1.1pt, SoftRed] table[x=x,y=y]{data/clean/profiles/pxzz_000663.dat};
\end{axis}
\end{tikzpicture}

\captionof{figure}[永安林业（000663）股价与周期骨架]{永安林业（000663）月度收盘价（元/股）与 ZigZag 周期骨架（阈值 25\%，圆点为周期转折点）}
\end{minipage}
\end{center}

\pfl{过去周期分析} 以 25\% 的 ZigZag 阈值划分，样本区间内股价形成 4 个主升段与 4 个主跌段。 最大主跌段为 2023 年 6 月 至 2024 年 6 月，12 个月-65.6\%；最大主升段出现在 2021 年 1 月 至 2021 年 6 月，5 个月+183.5\%。 最近一次转折出现在 2026 年 6 月（下跌段自 2025 年 11 月起，7 个月-41.5\%），当前处于该段的延续之中。 公司股价较申万农林牧渔指数提前约 23 个月见底，是板块中较早定价周期反转的公司之一。 林业的现金流依赖林木采伐与政策补贴，周期长、变现慢，公司 2026 年 6 月末资产负债率 36.6\%、半年 ROE -2.1\%（未年化），在本轮下行中属于随行业同步调整的一类。 综合区间形态看，该股单段涨跌幅度偏大、以趋势段为主（最大主升 +183.5\%、最大主跌 -65.6\%），年化波动率 51.7\% 与全样本中位数 37.4\% 相比波动明显大于样本中位数。

\pfl{盈利情况} 2026 年上半年营业收入 1.51 亿元（同比 +14.7\%），归母净利润 -0.19 亿元，亏损同比扩大 0.02 亿元。 以年报为趋势对照，2023---2025 年营业收入由 6.93 亿元变为 3.24 亿元（两年年均复合增速 -31.6\%），归母净利润由 1.90 亿元变为 -0.95 亿元。 按半年报口径，2026 年上半年的 ROE 为 -2.1\%（未年化）、销售净利率 -13.2\%、毛利率 11.9\%，ROE 在“林业Ⅲ”4 家公司中按数值由高到低排第\zhnumber{3}位。 按同一口径，三级行业“林业Ⅲ”4 家公司的半年 ROE 中位数为 -1.96\%，该公司低于其中位数 0.15 个百分点。 从公司披露的原因看，业绩变动主要来自：2025 年度，归属于上市公司股东的净利润 -9，510.41 万元（同比下降 11.84\%）；2026 年上半年，主要系生物质项目停建后产生的诉讼费用、资产减值及政府补助同比减少等因素影响。 综合看，公司仍处亏损区间、亏损同比扩大，半年 ROE 在三级行业内按由高到低排第\zhnumber{3}位，单位盈利能力的修复依赖行业景气与自身成本端的改善，报表弹性主要体现为价格与销量的方向性变化。

\begin{center}\small\setlength{\tabcolsep}{3.4pt}
\captionof{table}[永安林业（000663）关键财务指标]{永安林业（000663）关键财务指标（合并报表口径；两个半年报期均未年化；现金短债比 $=$ 货币资金 $\div$ 短期债务）}
\begin{adjustbox}{max width=\textwidth}
\begin{tabular}{@{}lrrrrr@{}}
\toprule
指标 & 2023 年 & 2024 年 & 2025 年 & 2025 年半年报 & 2026 年半年报 \\
\midrule
营业收入（亿元） & 6.93 & 3.23 & 3.24 & 1.32 & 1.51 \\
归母净利润（亿元） & 1.90 & -0.85 & -0.95 & -0.18 & -0.19 \\
毛利率（\%） & 44.7 & 9.9 & 11.4 & 8.0 & 11.9 \\
ROE（\%） & 18.9 & -8.0 & -9.8 & -1.7 & -2.1 \\
资产负债率（\%） & 32.2 & 33.6 & 37.3 & 32.4 & 36.6 \\
经营活动现金流净额（亿元） & 3.49 & 1.17 & -0.17 & 0.09 & -0.31 \\
货币资金（亿元） & 2.76 & 2.78 & 2.13 & 2.60 & 1.68 \\
有息负债合计（亿元） & 1.92 & 0.82 & 0.55 & 0.56 & 0.55 \\
现金短债比（倍） & 12.47 & 9.32 & 6.79 & 10.43 & 5.28 \\
\bottomrule
\end{tabular}
\end{adjustbox}
\end{center}

\begin{center}
\begin{minipage}{\textwidth}\centering
\begin{tikzpicture}
\begin{groupplot}[
  group style={group size=2 by 1, horizontal sep=0.60cm},
  width=6.85cm, height=4.60cm,
  tick label style={font=\tiny},
  title style={font=\scriptsize\heiti},
  yticklabel style={font=\tiny},
  ylabel style={font=\tiny},
  grid=major, grid style={LightGray, line width=0.3pt},
  axis line style={InkGray, line width=0.5pt},
  enlarge x limits={abs=0.8},
  legend style={font=\scriptsize, at={(0.5,-0.20)}, anchor=north,
                legend columns=2, draw=none, fill=none, column sep=6pt},
]
\nextgroupplot[
  ybar, bar width=4pt,
  ymin=-1.80, ymax=8.59,
  xtick={0,1,2,3,4,5.7},
  xticklabels={2021,2022,2023,2024,2025,2026H1},
  ylabel={亿元},
]
\addplot[fill=AgriGreen!75, draw=AgriGreen, bar shift=-2.6pt] table[x=i, y=rev]{data/clean/profiles/fin_000663.dat};
\addplot[fill=SoftRed!60, draw=SoftRed, bar shift=2.6pt] table[x=i, y=ni]{data/clean/profiles/fin_000663.dat};
\addplot[fill=AgriGreen!30, draw=AgriGreen, pattern=north east lines, bar shift=-2.6pt] table[x=x, y=rev]{data/clean/profiles/finh_000663.dat};
\addplot[fill=SoftRed!20, draw=SoftRed, pattern=north east lines, bar shift=2.6pt] table[x=x, y=ni]{data/clean/profiles/finh_000663.dat};
\legend{营业收入（年/半年）, 归母净利润（年/半年）}
\nextgroupplot[
  ymin=-14.6, ymax=56.2,
  xtick={0,1,2,3,4,5.7},
  xticklabels={2021,2022,2023,2024,2025,2026H1},
  ylabel={\%},
]
\addplot[AgriGold, mark=*, mark size=1.3pt] table[x=i, y=roe]{data/clean/profiles/fin_000663.dat};
\addplot[OilBlue, mark=square*, mark size=1.3pt] table[x=i, y=debt]{data/clean/profiles/fin_000663.dat};
\addplot[only marks, AgriGold, mark=*, mark size=2.0pt] table[x=x, y=roe]{data/clean/profiles/finh_000663.dat};
\addplot[only marks, OilBlue, mark=square*, mark size=2.0pt] table[x=x, y=debt]{data/clean/profiles/finh_000663.dat};
\legend{ROE, 资产负债率}
\end{groupplot}
\end{tikzpicture}
\begin{tikzpicture}
\begin{axis}[
  name=finbot,
  width=13.75cm, height=3.10cm, scale only axis,
  ymin=-0.52, ymax=5.86,
  xtick={0,1,2,3,4,5.7},
  xticklabels={2021,2022,2023,2024,2025,2026H1},
  tick label style={font=\tiny},
  yticklabel style={font=\tiny},
  ylabel={亿元}, ylabel style={font=\tiny},
  ybar, bar width=4pt,
  grid=major, grid style={LightGray, line width=0.3pt},
  axis line style={InkGray, line width=0.5pt},
  enlarge x limits={abs=0.8},
  legend style={font=\scriptsize, at={(0.5,-0.26)}, anchor=north,
                legend columns=2, draw=none, fill=none, column sep=10pt},
]
\addplot[fill=AgriGreen!55, draw=AgriGreen, bar shift=-3.0pt] table[x=i, y=cash]{data/clean/profiles/fin_000663.dat};
\addplot[fill=SoftRed!45, draw=SoftRed, bar shift=3.0pt] table[x=i, y=ibd]{data/clean/profiles/fin_000663.dat};
\addplot[fill=AgriGreen!25, draw=AgriGreen, pattern=north east lines, bar shift=-3.0pt] table[x=x, y=cash]{data/clean/profiles/finh_000663.dat};
\addplot[fill=SoftRed!20, draw=SoftRed, pattern=north east lines, bar shift=3.0pt] table[x=x, y=ibd]{data/clean/profiles/finh_000663.dat};
\legend{货币资金（左轴）, 有息负债（左轴）}
\end{axis}
\begin{axis}[
  at={(finbot.south east)}, anchor=south east,
  width=13.75cm, height=3.10cm, scale only axis,
  axis y line*=right, axis x line*=none, xtick=\empty,
  ymin=-1.55, ymax=5.85,
  tick label style={font=\tiny},
  yticklabel style={font=\tiny}, scaled y ticks=false,
  legend style={font=\scriptsize, at={(0.5,-0.44)}, anchor=north,
                legend columns=1, draw=none, fill=none},
]
\addplot[OilBlue, mark=*, mark size=2.2pt, line width=1.3pt] table[x=i, y=ocf]{data/clean/profiles/fin_000663.dat};
\addplot[only marks, OilBlue, mark=*, mark size=3.2pt] table[x=x, y=ocf]{data/clean/profiles/finh_000663.dat};
\legend{经营活动现金流净额（右轴，亿元，蓝点）}
\end{axis}
\end{tikzpicture}

\captionof{figure}[永安林业（000663）近年主要财务指标]{永安林业（000663）近年主要财务指标（左上：营业收入与归母净利润；右上：ROE 与资产负债率；下：货币资金、有息负债为左轴柱状，经营活动现金流净额为其右轴的蓝点折线。
2021---2025 年为年报口径，2026H1 为 2026 年半年报/半年末，与其前一年报之间留出间隔、以斜纹或空心点区分；半年报数据未年化）}
\end{minipage}
\end{center}

\pfl{财务分析} 2026 年 6 月末资产负债率 36.6\%，较 2025 年末下降 0.7 个百分点（样本中位数 52.2\%，所处三级行业内按由低到高排第\zhnumber{2}位）。 2026 年 6 月末货币资金 1.68 亿元、短期债务（短期借款与一年内到期非流动负债）0.32 亿元、长期债务 0.23 亿元，有息负债合计 0.55 亿元，现金短债比 5.28 倍——覆盖充分。 净有息负债 -1.13 亿元，相当于其 2026 年上半年营业收入的 -75\%。 流动比率 1.56、速动比率 0.56，短期流动性无明显缺口。 2026 年上半年经营活动现金流净额 -0.31 亿元，经营现金流与利润同时为负，现金消耗较快。 2026 年 6 月末在建工程 2.79 亿元，较上年末增加 0.06 亿元，仍有在建投入。 综合看，现金短债比 5.28 倍，短期偿付压力可控；资产负债率低于样本中位数 15.6 个百分点；有息负债中短期债务占比更高，期限结构仍以滚动续作为主。 公司披露的诉讼与资产冻结风险为：2026 年 7 月，闽民终 658 号民事判决向最高人民法院申请再审的张强等 17 人诉讼案（涉案金额 3，172.55 万元）被最高人民法院裁定驳回再审申请；2025 年 10 月，公司全资孙公司向永安市人民法院提起对必奥新能源科技有限公司的买卖合同纠纷诉讼。

\pfl{重大战略转型} 公司的战略动作可归为以下几类：其他动作（2 项）：2021 年 2 月，林业发展集团完成股份过户登记手续；2020 年 4 月，中林控股与林业发展集团以持有的 6。 公告口径上，近三年（2023 年 9 月---2026 年 9 月）公司披露重大重组与并购类公告 10 条、对外投资与转型类公告 1 条、回购与股权激励类公告 7 条，合计公告 370 条。 第一大行业“人造板制造业”的收入占比由 2021 年年报的 57.4\% 变为 93.1\%（35.7 个百分点），收入结构出现再平衡。 综合判断，报告期内公司有 11 条与战略相关的公告落地（重大重组与并购 10 条、对外投资与转型 1 条），资本运作与主业调整之间存在直接关联。

\begin{center}\small\setlength{\tabcolsep}{4pt}
\captionof{table}[永安林业（000663）历史融资与资本运作]{永安林业（000663）历史融资与资本运作（据公司公告与公开报道整理；金额为披露值，含首次公开发行、再融资、债券、银行借款与重整投资等）}
\begin{adjustbox}{max width=\textwidth}
\begin{tabular}{@{}p{1.45cm}p{2.55cm}p{4.05cm}p{4.95cm}@{}}
\toprule
时间 & 方式 & 规模 & 用途与说明 \\
\midrule
1996-11 & IPO（深圳证券交易所） & 向社会公开募集股票 1,\allowbreak{}950 万股，\allowbreak{}共募得资金 9,\allowbreak{}150.00 万元 & 招股说明书承诺投资项目（后续发生变更，\allowbreak{}变更原因与程序已在 1997 年年报披露）；公司 1996 年 12 月 6 日上市；1997 年 5 月 12 日股东大会决议向全体股东每 10 股派送 1 股红股并以资本公积按 10:7 转增股本 \\
2015-09 & 发行股份及支付现金购买资产并募集配套资金（重大资产重组） & 以发行股份方式支付 12.50 亿元、\allowbreak{}发行价格 11.75 元/\allowbreak{}股 & 本次交易现金对价支付、\allowbreak{}森源股份信息系统升级改造和营销与服务网点建设、\allowbreak{}补充上市公司流动资金；发行股份购买资产所发行新股上市日为 2015 年 10 月 8 日，\allowbreak{}增发后公司总股本 3.09 亿股；配套融资所发行新股上市首日为 2015 年 12 月 23 日 \\
2023-11-20 & 2022 年度向特定对象发行 A 股股票（终止） & 拟募集资金总额不超过人民币 3.00 亿元，\allowbreak{}发行数量不超过 4，\allowbreak{}360.47 万股 & 涿州市生物质能循环利用项目、\allowbreak{}宁晋县生物质能循环利用项目、\allowbreak{}生物质发酵微生物研发中心项目；本次发行方案于 2022 年 9 月 30 日经第九届董事会第二十八次会议审议通过、\allowbreak{}2022 年 10 月 17 日经 2022 年第五次临时股东大会审议通过；2023 年 11 月 20 日第十届董事会第三次会议和第十届监事会第三次会议审议通过终止公司 2022 年度向特定对象发行股票事项并撤回申请文件并向深交所申请撤回相关申请文件 \\
\bottomrule
\end{tabular}
\end{adjustbox}
\end{center}

\pfl{历史融投资情况} 从现金流量表看，2025 年公司取得借款收到的现金 0.23 亿元、偿还债务支付的现金 0.49 亿元，筹资活动现金流净额 -0.29 亿元（2023 年为取得借款 1.88 亿元）。 2026 年上半年筹资活动现金流净额为 -0.01 亿元（取得借款 — 亿元、偿还债务 — 亿元，未年化），已转为净流出，处于净还债状态。 2025 年分配股利、利润或偿付利息支付的现金合计 0.02 亿元。 公开披露的融资记录共 3 笔，工具类型以 IPO、发行股份及支付现金购、定向增发为主；金额与用途见下表。 其中规模最大的两笔为：2015 年 9 月，发行股份及支付现金购，以发行股份方式支付 12.50 亿元、发行价格 11.75 元/股；1996 年 11 月，IPO，向社会公开募集股票 1,950 万股，共募得资金 9,150.00 万元。 股本方面，近三年披露股本变动 3 次；总股本自 2021 年末以来未发生变化（3.37 亿股）。 分红方面，近三年未实施现金分红，在全部样本中属于分红能力偏弱的一端。 近三年“再融资”类公告 2 条，涉及定向增发。 综合看，公司融资结构上股权与债务融资并举；2025 年筹资活动现金流净额 -0.29 亿元，融资节奏仍以维持存量债务周转为主，与前述经营与投资现金流的方向基本一致。

\pfl{未来潜在融资需求分析} 2026 年 6 月末货币资金 1.68 亿元，而短期债务 0.32 亿元（现金短债比 5.28 倍），2026 年上半年经营现金流净额 -0.31 亿元。账面货币资金可以覆盖短期债务，缺口不体现在静态偿付层面。 按第 \ref{sec:financing-need} 节的三档划分，公司属于第二档“保壳型”：近三年重大重组与并购类公告 10 条，2026 年上半年归母净利润 -0.19 亿元，融资需求更多体现为“资产置换 + 维持上市地位”，而不是扩产。 公开披露的后续资金安排线索包括：股权融资线索：2023 年 11 月，原计划由中林集团全额认购、募资不超过 3.00 亿元的股权融资方案落空。 资金需求还要叠加在建工程：期末在建工程 2.79 亿元、2026 年上半年购建固定资产等支出 0.15 亿元（未年化），这部分支出在周期底部通常会被主动放缓。 静态偿付不存在缺口，后续融资更可能服务于产能或并购安排，而不是填补流动性窟窿。






\subsubsection*{农业综合Ⅱ（2 家）}
\addcontentsline{toc}{subsubsection}{农业综合Ⅱ（2 家）}

% 农业上市公司画像：农业综合Ⅱ（共 2 家，按流通市值降序）
% 由 scripts/make_company_profiles.py 自动生成，请勿手工修改

\subsubsection{辉隆股份（002556）}
\label{co:002556}

\pfl{公司基本情况} 辉隆股份成立于 2011 年，安徽省供销社系统背景的农资流通龙头，形成化肥、农药等农资流通与精细化工“双主业”，并以 14 座现代农业综合服务中心为载体构建“耕种管收售”全周期农业服务体系。 公司于 2011-03-02 在深交所主板首发上市，注册资本 9.35 亿元，总股本 9.47 亿股。 截至 2026-09-24收盘价 5.08 元，流通市值 47.4 亿元，近一年-7.3\%，流通市值在三级行业“—”的 2 家公司中排第\zhnumber{1}位，近一年涨跌幅高于全样本中位数（-10.1\%）2.8 个百分点。 按 2026 年半年报披露的行业构成，收入占比前三位为农资产品 73.2\%；农副产品及其他 12.3\%；化工产品 8.1\%。 与 2021 年年报相比，第一大行业口径由“化肥”（62.1\%）变为“农资产品”（73.2\%）；两期披露的分类口径有调整，占比差异中同时包含分类因素。 发展过程中的关键节点包括：2020-02-25 因发行股份购买资产新增的 118,353,739 股（发行价格 5.13 元/股）上市流通前完成登记……（2020）；2023-07-28 发行 2023 年度第一期中期票据（乡村振兴）“23 辉隆农资 MTN001(乡村振兴)”……（2023）；2024-03-22 发行 2024 年度第一期中期票据（乡村振兴）“24 辉隆农资 MTN001(乡村振兴)”……（2024）。

\pfl{历史股价} 自 2021 年 1 月至 2026 年 9 月，股价由 8.11 元变为 5.08 元，累计 -37.4\%，区间最高 12.87 元（2021 年 12 月）、最低 4.17 元（2024 年 6 月）；当前价相当于区间最高价的 39.5\%，为区间最低价的 1.22 倍。 月度收益的年化波动率 30.5\%，低于全样本中位数 37.4\%，在三级行业“—”的 2 家公司中波动率排第\zhnumber{1}位。 从事件看，2024 年 6 月 至 2026 年 2 月的上涨（+61.2\%）与公司媒体报道控股股东解除 410 万股质押（与 2025 年末披露的质押 24,370,000 股口径存在差异……（2026-01-26）在时间上重合——股价的拐点多数能在公司自身的资本运作与风险事件上找到对应。

\begin{center}
\begin{minipage}{\textwidth}\centering
\begin{tikzpicture}
\begin{axis}[
  width=12.6cm, height=3.9cm, scale only axis,
  xmin=2020.922, xmax=2026.888, ymin=3.88, ymax=13.77,
  xtick={2021,2022,2023,2024,2025,2026},
  yearticks,
  yticklabel style={font=\scriptsize},
  ylabel={元/股}, ylabel style={font=\scriptsize},
  grid=major, grid style={LightGray, line width=0.3pt},
  axis line style={InkGray, line width=0.5pt},
  every axis plot/.append style={line width=0.9pt},
]
\addplot[AgriGreen] table[x=x,y=y]{data/clean/profiles/px_002556.dat};
\addplot[SoftRed, dashed, line width=0.7pt] table[x=x,y=y]{data/clean/profiles/pxzz_002556.dat};
\addplot[only marks, mark=*, mark size=1.1pt, SoftRed] table[x=x,y=y]{data/clean/profiles/pxzz_002556.dat};
\end{axis}
\end{tikzpicture}

\captionof{figure}[辉隆股份（002556）股价与周期骨架]{辉隆股份（002556）月度收盘价（元/股）与 ZigZag 周期骨架（阈值 25\%，圆点为周期转折点）}
\end{minipage}
\end{center}

\pfl{过去周期分析} 以 25\% 的 ZigZag 阈值划分，样本区间内股价形成 2 个主升段与 2 个主跌段。 最大主升段出现在 2024 年 6 月 至 2026 年 2 月，20 个月+61.2\%；最大主跌段为 2021 年 12 月 至 2024 年 6 月，30 个月-67.6\%。 最近一次转折出现在 2026 年 6 月（下跌段自 2026 年 2 月起，4 个月-33.3\%），当前处于该段的延续之中。 公司股价较申万农林牧渔指数提前约 23 个月见底，是板块中较早定价周期反转的公司之一。 农业综合类公司业务分散，估值更多取决于资产与转型预期而非单一品种价格，公司 2026 年 6 月末资产负债率 67.0\%，高于全样本中位数 52.2\%，其股价因此更多反映资金面与信用风险，而不是价格弹性本身。

\pfl{盈利情况} 2026 年上半年营业收入 85.47 亿元（同比 +3.3\%），归母净利润 1.13 亿元，同比 +1.5\%。 以年报为趋势对照，2023---2025 年营业收入由 178.33 亿元变为 151.32 亿元（两年年均复合增速 -7.9\%），归母净利润由 0.77 亿元变为 1.94 亿元。 按半年报口径，2026 年上半年的 ROE 为 3.1\%（未年化）、销售净利率 1.4\%、毛利率 6.2\%，ROE 在“—”2 家公司中按数值由高到低排第\zhnumber{1}位。 同期全样本半年 ROE 中位数 0.75\%，公司与之相差 2.33 个百分点（略好于中位数公司）。 从公开披露看，2026 年半年度计提资产减值准备 10,022.40 万元（项目主要为应收账款、其他应收款、存货、固定资产等）……；2026 年上半年经营活动产生的现金流量净额 -417,263,558.44 元……。

\begin{center}\small\setlength{\tabcolsep}{3.4pt}
\captionof{table}[辉隆股份（002556）关键财务指标]{辉隆股份（002556）关键财务指标（合并报表口径；两个半年报期均未年化；现金短债比 $=$ 货币资金 $\div$ 短期债务）}
\begin{adjustbox}{max width=\textwidth}
\begin{tabular}{@{}lrrrrr@{}}
\toprule
指标 & 2023 年 & 2024 年 & 2025 年 & 2025 年半年报 & 2026 年半年报 \\
\midrule
营业收入（亿元） & 178.33 & 156.52 & 151.32 & 82.77 & 85.47 \\
归母净利润（亿元） & 0.77 & 1.69 & 1.94 & 1.11 & 1.13 \\
毛利率（\%） & 4.8 & 5.8 & 6.3 & 5.8 & 6.2 \\
ROE（\%） & 2.0 & 4.6 & 5.2 & 3.0 & 3.1 \\
资产负债率（\%） & 66.2 & 64.6 & 64.9 & 67.9 & 67.0 \\
经营活动现金流净额（亿元） & -3.73 & -1.54 & 10.71 & -1.09 & -4.17 \\
货币资金（亿元） & 11.16 & 9.62 & 11.59 & 11.67 & 11.20 \\
有息负债合计（亿元） & 31.71 & 29.92 & 25.60 & 34.29 & 31.42 \\
现金短债比（倍） & 0.70 & 0.79 & 0.82 & 0.65 & 0.40 \\
\bottomrule
\end{tabular}
\end{adjustbox}
\end{center}

\begin{center}
\begin{minipage}{\textwidth}\centering
\begin{tikzpicture}
\begin{groupplot}[
  group style={group size=2 by 1, horizontal sep=0.60cm},
  width=6.85cm, height=4.60cm,
  tick label style={font=\tiny},
  title style={font=\scriptsize\heiti},
  yticklabel style={font=\tiny},
  ylabel style={font=\tiny},
  grid=major, grid style={LightGray, line width=0.3pt},
  axis line style={InkGray, line width=0.5pt},
  enlarge x limits={abs=0.8},
  legend style={font=\scriptsize, at={(0.5,-0.20)}, anchor=north,
                legend columns=2, draw=none, fill=none, column sep=6pt},
]
\nextgroupplot[
  ybar, bar width=4pt,
  ymin=-19.08, ymax=213.71,
  xtick={0,1,2,3,4,5.7},
  xticklabels={2021,2022,2023,2024,2025,2026H1},
  ylabel={亿元},
]
\addplot[fill=AgriGreen!75, draw=AgriGreen, bar shift=-2.6pt] table[x=i, y=rev]{data/clean/profiles/fin_002556.dat};
\addplot[fill=SoftRed!60, draw=SoftRed, bar shift=2.6pt] table[x=i, y=ni]{data/clean/profiles/fin_002556.dat};
\addplot[fill=AgriGreen!30, draw=AgriGreen, pattern=north east lines, bar shift=-2.6pt] table[x=x, y=rev]{data/clean/profiles/finh_002556.dat};
\addplot[fill=SoftRed!20, draw=SoftRed, pattern=north east lines, bar shift=2.6pt] table[x=x, y=ni]{data/clean/profiles/finh_002556.dat};
\legend{营业收入（年/半年）, 归母净利润（年/半年）}
\nextgroupplot[
  ymin=-5.4, ymax=73.7,
  xtick={0,1,2,3,4,5.7},
  xticklabels={2021,2022,2023,2024,2025,2026H1},
  ylabel={\%},
]
\addplot[AgriGold, mark=*, mark size=1.3pt] table[x=i, y=roe]{data/clean/profiles/fin_002556.dat};
\addplot[OilBlue, mark=square*, mark size=1.3pt] table[x=i, y=debt]{data/clean/profiles/fin_002556.dat};
\addplot[only marks, AgriGold, mark=*, mark size=2.0pt] table[x=x, y=roe]{data/clean/profiles/finh_002556.dat};
\addplot[only marks, OilBlue, mark=square*, mark size=2.0pt] table[x=x, y=debt]{data/clean/profiles/finh_002556.dat};
\legend{ROE, 资产负债率}
\end{groupplot}
\end{tikzpicture}
\begin{tikzpicture}
\begin{axis}[
  name=finbot,
  width=13.75cm, height=3.10cm, scale only axis,
  ymin=-3.17, ymax=35.52,
  xtick={0,1,2,3,4,5.7},
  xticklabels={2021,2022,2023,2024,2025,2026H1},
  tick label style={font=\tiny},
  yticklabel style={font=\tiny},
  ylabel={亿元}, ylabel style={font=\tiny},
  ybar, bar width=4pt,
  grid=major, grid style={LightGray, line width=0.3pt},
  axis line style={InkGray, line width=0.5pt},
  enlarge x limits={abs=0.8},
  legend style={font=\scriptsize, at={(0.5,-0.26)}, anchor=north,
                legend columns=2, draw=none, fill=none, column sep=10pt},
]
\addplot[fill=AgriGreen!55, draw=AgriGreen, bar shift=-3.0pt] table[x=i, y=cash]{data/clean/profiles/fin_002556.dat};
\addplot[fill=SoftRed!45, draw=SoftRed, bar shift=3.0pt] table[x=i, y=ibd]{data/clean/profiles/fin_002556.dat};
\addplot[fill=AgriGreen!25, draw=AgriGreen, pattern=north east lines, bar shift=-3.0pt] table[x=x, y=cash]{data/clean/profiles/finh_002556.dat};
\addplot[fill=SoftRed!20, draw=SoftRed, pattern=north east lines, bar shift=3.0pt] table[x=x, y=ibd]{data/clean/profiles/finh_002556.dat};
\legend{货币资金（左轴）, 有息负债（左轴）}
\end{axis}
\begin{axis}[
  at={(finbot.south east)}, anchor=south east,
  width=13.75cm, height=3.10cm, scale only axis,
  axis y line*=right, axis x line*=none, xtick=\empty,
  ymin=-8.49, ymax=17.41,
  tick label style={font=\tiny},
  yticklabel style={font=\tiny}, scaled y ticks=false,
  legend style={font=\scriptsize, at={(0.5,-0.44)}, anchor=north,
                legend columns=1, draw=none, fill=none},
]
\addplot[OilBlue, mark=*, mark size=2.2pt, line width=1.3pt] table[x=i, y=ocf]{data/clean/profiles/fin_002556.dat};
\addplot[only marks, OilBlue, mark=*, mark size=3.2pt] table[x=x, y=ocf]{data/clean/profiles/finh_002556.dat};
\legend{经营活动现金流净额（右轴，亿元，蓝点）}
\end{axis}
\end{tikzpicture}

\captionof{figure}[辉隆股份（002556）近年主要财务指标]{辉隆股份（002556）近年主要财务指标（左上：营业收入与归母净利润；右上：ROE 与资产负债率；下：货币资金、有息负债为左轴柱状，经营活动现金流净额为其右轴的蓝点折线。
2021---2025 年为年报口径，2026H1 为 2026 年半年报/半年末，与其前一年报之间留出间隔、以斜纹或空心点区分；半年报数据未年化）}
\end{minipage}
\end{center}

\pfl{财务分析} 2026 年 6 月末资产负债率 67.0\%，较 2025 年末 上升 2.1 个百分点（样本中位数 52.2\%，所处三级行业内按由低到高排第\zhnumber{1}位）。 2026 年 6 月末货币资金 11.20 亿元、短期债务（短期借款与一年内到期非流动负债）27.68 亿元、长期债务 3.74 亿元，有息负债合计 31.42 亿元，现金短债比 0.40 倍——覆盖偏紧。 净有息负债 20.22 亿元，相当于其 2026 年上半年营业收入的 24\%。 流动比率 1.01、速动比率 0.55，短期指标尚可。 2026 年上半年经营活动现金流净额 -4.17 亿元，净利润为正但经营现金流为负，利润的现金含量偏低。 公司披露的财务与合规风险包括：受限货币资金 384,448,852.39 元，受限类型为保证，用于开具银行承兑汇票、信用证、保函、借款……；媒体报道控股股东解除 410 万股质押（与 2025 年末披露的质押 24,370,000 股口径存在差异……。

\pfl{重大战略转型} 近三年公司的战略动作集中在以下几个方面：2026-02-04 媒体报道公司全资孙公司中成科技位于硫铁矿带；参股企业（中盟磷业）白岩磷矿储量可达 9,500 万吨；2026-04-17 披露《关于“质量回报双提升”行动方案的公告》（公告编号 2026-021），2026-08-13 披露进展：注销回购股份 11,577,228 股……；2026-06-30 聚焦“农资+精细化工”双主业提质增效，从资源端、产品端和市场端夯实主业基石，同时推进科研应用新赛道转型；2026-06-30 农服板块围绕种肥药一体化服务延伸链条，对接优质种质资源、落地绿色农资示范项目、新增秧苗培育业务。 公告口径上，近三年（2023 年 9 月---2026 年 9 月）公司披露重大重组与并购类公告 0 条、对外投资与转型类公告 1 条、回购与股权激励类公告 15 条，合计公告 341 条。 主营结构的口径变化已经落到报表上：2021 年年报的第一大行业为“化肥”（62.1\%），2026 年半年报为“农资产品”（73.2\%）；公司在这两期的分部划分方式有过调整。 综合判断，公司在报告期内有明确的外延扩张或业务结构调整动作，转型方向与融资安排之间存在直接关联。

\begin{center}\small\setlength{\tabcolsep}{4pt}
\captionof{table}[辉隆股份（002556）历史融资与资本运作]{辉隆股份（002556）历史融资与资本运作（据公司公告与公开报道整理；金额为披露值，含首次公开发行、再融资、债券、银行借款与重整投资等）}
\begin{adjustbox}{max width=\textwidth}
\begin{tabular}{@{}p{1.45cm}p{2.55cm}p{4.05cm}p{4.95cm}@{}}
\toprule
时间 & 方式 & 规模 & 用途与说明 \\
\midrule
2011-02 & IPO（深圳证券交易所中小板） & 发行 3,\allowbreak{}750 万股，\allowbreak{}发行价 38 元/\allowbreak{}股，\allowbreak{}募集资金总额 14.06 亿元；招股说明书披露…… & 配送中心建设项目和信息化系统建设项目；经证监许可[2011]197 号文核准…… \\
2018-2019 /\allowbreak{} 2020-02 & 发行股份、可转换公司债券及支付现金购买资产（关联交易） & 标的为海华科技 100\% 股权，\allowbreak{}总对价 8.28 亿元：现金 1.81 亿元、\allowbreak{}定向可转债 4…… & 收购海华科技 100\% 股权（交易对方为辉隆投资、蚌埠隆海及解凤贤、解凤苗…… \\
2020-04 & 非公开发行股份及可转换公司债券募集配套资金 & 募集配套资金 6.44 亿元，\allowbreak{}其中非公开发行股份募集 1.30 亿元、\allowbreak{}非公开发行可转换公司债…… & 向交易对方支付现金对价、海华科技项目建设及支付本次交易相关费用…… \\
2023-07-28 & 中期票据（乡村振兴） & 4.00 亿元，\allowbreak{}票面利率 3.50\%，\allowbreak{}到期日 2028-08-01（债券代码 1023818…… & 2025 年度报告披露期末余额 405,833,333.33 元…… \\
2024-03-22 & 中期票据（乡村振兴） & 4.00 亿元，\allowbreak{}票面利率 2.80\%，\allowbreak{}到期日 2027-03-26（债券代码 1024811…… & 2025 年度报告披露期末余额 408,788,888.89 元…… \\
2026-04-17 & 终止部分超募资金投资项目 & 剩余超募资金 4,\allowbreak{}343.59 万元 & 终止首发超募资金投资项目“40 万吨/年新型肥料项目”…… \\
\bottomrule
\end{tabular}
\end{adjustbox}
\end{center}

\pfl{历史融投资情况} 从现金流量表看，2025 年公司取得借款收到的现金 18.70 亿元、偿还债务支付的现金 25.51 亿元、吸收权益性投资 0.01 亿元，筹资活动现金流净额 -10.40 亿元（2023 年为取得借款 36.24 亿元、偿还债务 28.20 亿元）。 2026 年上半年筹资活动现金流净额为 0.00 亿元（取得借款 14.36 亿元、偿还债务 13.04 亿元，未年化），仍处于净融入状态。 2025 年分配股利、利润或偿付利息支付的现金合计 3.18 亿元。 公司披露的股本与债务融资脉络包括：IPO（深圳证券交易所中小板）、中期票据（乡村振兴）、发行股份、可转换公司债券及支付现金购买资产（关联交易）、终止部分超募资金投资项目、非公开发行股份及可转换公司债券募集配套资金。 其中规模较大的两笔为：2026-04-17终止部分超募资金投资项目（剩余超募资金 4,343.59 万元）；2020-04非公开发行股份及可转换公司债券募集配套资金（募集配套资金 64,406.52 万元，其中非公开发行股份募集 13,000 万元……）。 股本方面，近三年披露股本变动 9 次（其他 3 次、限售股份上市 2 次、股份回购 1 次）；总股本由 2021 年末的 9.54 亿股变为 9.47 亿股（-0.78\%）。 分红方面，近三年现金分红 5 次、累计每股派现 0.85 元（最近一次 2025-12-31）——在利润承压的阶段仍维持了分红。

\pfl{未来潜在融资需求分析} 2026 年 6 月末货币资金 11.20 亿元，而短期债务 27.68 亿元（现金短债比 0.40 倍），2026 年上半年经营现金流净额 -4.17 亿元。按“短期债务减去账面货币资金”的滚动偿付口径，静态缺口约 16.47 亿元；若再计入上半年亏损的年化额（0.00 亿元），维持一年正常经营所需的外部资金约 16.47---16.47 亿元。 公司 2026 年 6 月末资产负债率 67.0\%、2026 年上半年 ROE 3.1\%（未年化），资产负债表尚有余量，融资需求不是“补血型”的刚性需求，而是围绕并购整合、项目投入与营运资金的主动性需求。 公开披露的下一步融资线索包括：2026-05 2026-05-07 披露变更回购股份用途并注销、减少注册资本及修订《公司章程》……；2026-06-30 披露对控股公司提供担保及募集资金存放与实际使用情况专项报告；货币资金 1,030,458,766.52 元；2026-08-07 披露《关于 2023 年度第一期中期票据(乡村振兴)付息及回售兑付完成的公告》，4 亿元中期票据进入回售兑付环节。




\subsubsection{大禹节水（300021）}
\label{co:300021}

\pfl{公司基本情况} 大禹节水成立于 2009 年，国内节水灌溉与农田水利领域主要上市公司，业务覆盖节水材料装备制造、灌区 EPC/EPCO 总承包与农水设计服务，正向“设计引领+全产业链”模式转型。 公司于 2009-10-30 在深交所创业板首发上市，注册资本 10.22 亿元，总股本 8.54 亿股。 截至 2026-09-24收盘价 3.82 元，流通市值 33.4 亿元，近一年-22.7\%，流通市值在三级行业“—”的 2 家公司中排第\zhnumber{2}位，近一年涨跌幅低于全样本中位数（-10.1\%）12.6 个百分点。 按 2025 年年报披露的行业构成，收入占比前三位为智慧水利行业 100.0\%。 与 2021 年年报相比，第一大行业口径由“节水灌溉行业”（81.0\%）变为“智慧水利行业”（100.0\%）；两期披露的分类口径有调整，占比差异中同时包含分类因素。 发展过程中的关键节点包括：2010-04-14 召开 2009 年年度股东大会，审议通过以 2009 年末总股本 69,650,000 股为基数……（2010）；2020-07-28 公开发行 638 万张可转换公司债券（每张面值 100 元，发行总额 6.38 亿元）……（2020）；2022-02-24 以简易程序向特定对象发行股票 58,593,750 股（发行价 5.12 元/股）募集资金到账……（2022）。

\pfl{历史股价} 以 2021 年 1 月为起点、2026 年 9 月为终点，股价由 4.76 元变为 3.82 元，累计 -19.7\%，区间最高 6.32 元（2025 年 7 月）、最低 3.25 元（2024 年 8 月）；当前价相当于区间最高价的 60.4\%，为区间最低价的 1.18 倍。 月度收益的年化波动率 27.2\%，低于全样本中位数 37.4\%，在三级行业“—”的 2 家公司中波动率排第\zhnumber{2}位。 从事件看，2024 年 8 月 至 2025 年 7 月的上涨（+94.5\%）与公司可转债募集资金：本报告期投入 8,418.05 万元……（2025-06-30）在时间上重合；2021 年 12 月 至 2024 年 8 月的下跌（-44.9\%）与公司股东大会通过变更部分募集资金用途，将可转债募投项目尚未使用募集资金及专户存款……（2024-02-05）在时间上重合——股价的拐点多数能在公司自身的资本运作与风险事件上找到对应。

\begin{center}
\begin{minipage}{\textwidth}\centering
\begin{tikzpicture}
\begin{axis}[
  width=12.6cm, height=3.9cm, scale only axis,
  xmin=2020.922, xmax=2026.888, ymin=3.02, ymax=6.76,
  xtick={2021,2022,2023,2024,2025,2026},
  yearticks,
  yticklabel style={font=\scriptsize},
  ylabel={元/股}, ylabel style={font=\scriptsize},
  grid=major, grid style={LightGray, line width=0.3pt},
  axis line style={InkGray, line width=0.5pt},
  every axis plot/.append style={line width=0.9pt},
]
\addplot[AgriGreen] table[x=x,y=y]{data/clean/profiles/px_300021.dat};
\addplot[SoftRed, dashed, line width=0.7pt] table[x=x,y=y]{data/clean/profiles/pxzz_300021.dat};
\addplot[only marks, mark=*, mark size=1.1pt, SoftRed] table[x=x,y=y]{data/clean/profiles/pxzz_300021.dat};
\end{axis}
\end{tikzpicture}

\captionof{figure}[大禹节水（300021）股价与周期骨架]{大禹节水（300021）月度收盘价（元/股）与 ZigZag 周期骨架（阈值 25\%，圆点为周期转折点）}
\end{minipage}
\end{center}

\pfl{过去周期分析} 以 25\% 的 ZigZag 阈值划分，样本区间内股价形成 2 个主升段与 2 个主跌段。 最大主跌段为 2021 年 12 月 至 2024 年 8 月，32 个月-44.9\%；最大主升段出现在 2024 年 8 月 至 2025 年 7 月，11 个月+94.5\%。 最近一次转折出现在 2026 年 9 月（下跌段自 2025 年 7 月起，14 个月-39.6\%），当前处于该段的延续之中。 公司股价较申万农林牧渔指数提前约 21 个月见底，是板块中较早定价周期反转的公司之一。 农业综合类公司业务分散，估值更多取决于资产与转型预期而非单一品种价格，公司 2026 年 6 月末资产负债率 71.7\%，高于全样本中位数 52.2\%，其股价因此更多反映资金面与信用风险，而不是价格弹性本身。

\pfl{盈利情况} 2026 年上半年营业收入 15.86 亿元（同比 +24.3\%），归母净利润 -0.01 亿元，由上年同期的盈利转为亏损。 以年报为趋势对照，2023---2025 年营业收入由 34.53 亿元变为 37.61 亿元（两年年均复合增速 +4.4\%），归母净利润由 0.50 亿元变为 0.49 亿元。 按半年报口径，2026 年上半年的 ROE 为 -0.0\%（未年化）、销售净利率 -0.4\%、毛利率 20.6\%，ROE 在“—”2 家公司中按数值由高到低排第\zhnumber{2}位。 同期全样本半年 ROE 中位数 0.75\%，公司与之相差 0.79 个百分点（弱于中位数公司）。 从公开披露看，转让所持慧图科技 28.5\% 股权给嘉兴乾瞻衡远创业投资合伙企业（有限合伙），交易价格 19,237.50 万元……；2025 年度营业收入 376,124.83 万元，同比下降 14.10\%；归属于上市公司股东的净利润 4,890.01 万元……。

\begin{center}\small\setlength{\tabcolsep}{3.4pt}
\captionof{table}[大禹节水（300021）关键财务指标]{大禹节水（300021）关键财务指标（合并报表口径；两个半年报期均未年化；现金短债比 $=$ 货币资金 $\div$ 短期债务）}
\begin{adjustbox}{max width=\textwidth}
\begin{tabular}{@{}lrrrrr@{}}
\toprule
指标 & 2023 年 & 2024 年 & 2025 年 & 2025 年半年报 & 2026 年半年报 \\
\midrule
营业收入（亿元） & 34.53 & 43.79 & 37.61 & 12.76 & 15.86 \\
归母净利润（亿元） & 0.50 & 0.81 & 0.49 & 0.13 & -0.01 \\
毛利率（\%） & 24.3 & 17.8 & 21.4 & 23.3 & 20.6 \\
ROE（\%） & 2.5 & 3.9 & 2.2 & 0.6 & -0.0 \\
资产负债率（\%） & 71.2 & 75.1 & 70.6 & 73.6 & 71.7 \\
经营活动现金流净额（亿元） & -0.94 & 5.97 & -2.76 & -5.49 & -5.69 \\
货币资金（亿元） & 11.51 & 16.88 & 14.69 & 12.02 & 12.85 \\
有息负债合计（亿元） & 25.14 & 27.10 & 25.96 & 29.00 & 29.12 \\
现金短债比（倍） & 1.14 & 1.22 & 0.98 & 0.72 & 0.70 \\
\bottomrule
\end{tabular}
\end{adjustbox}
\end{center}

\begin{center}
\begin{minipage}{\textwidth}\centering
\begin{tikzpicture}
\begin{groupplot}[
  group style={group size=2 by 1, horizontal sep=0.60cm},
  width=6.85cm, height=4.60cm,
  tick label style={font=\tiny},
  title style={font=\scriptsize\heiti},
  yticklabel style={font=\tiny},
  ylabel style={font=\tiny},
  grid=major, grid style={LightGray, line width=0.3pt},
  axis line style={InkGray, line width=0.5pt},
  enlarge x limits={abs=0.8},
  legend style={font=\scriptsize, at={(0.5,-0.20)}, anchor=north,
                legend columns=2, draw=none, fill=none, column sep=6pt},
]
\nextgroupplot[
  ybar, bar width=4pt,
  ymin=-4.39, ymax=49.04,
  xtick={0,1,2,3,4,5.7},
  xticklabels={2021,2022,2023,2024,2025,2026H1},
  ylabel={亿元},
]
\addplot[fill=AgriGreen!75, draw=AgriGreen, bar shift=-2.6pt] table[x=i, y=rev]{data/clean/profiles/fin_300021.dat};
\addplot[fill=SoftRed!60, draw=SoftRed, bar shift=2.6pt] table[x=i, y=ni]{data/clean/profiles/fin_300021.dat};
\addplot[fill=AgriGreen!30, draw=AgriGreen, pattern=north east lines, bar shift=-2.6pt] table[x=x, y=rev]{data/clean/profiles/finh_300021.dat};
\addplot[fill=SoftRed!20, draw=SoftRed, pattern=north east lines, bar shift=2.6pt] table[x=x, y=ni]{data/clean/profiles/finh_300021.dat};
\legend{营业收入（年/半年）, 归母净利润（年/半年）}
\nextgroupplot[
  ymin=-6.0, ymax=82.6,
  xtick={0,1,2,3,4,5.7},
  xticklabels={2021,2022,2023,2024,2025,2026H1},
  ylabel={\%},
]
\addplot[AgriGold, mark=*, mark size=1.3pt] table[x=i, y=roe]{data/clean/profiles/fin_300021.dat};
\addplot[OilBlue, mark=square*, mark size=1.3pt] table[x=i, y=debt]{data/clean/profiles/fin_300021.dat};
\addplot[only marks, AgriGold, mark=*, mark size=2.0pt] table[x=x, y=roe]{data/clean/profiles/finh_300021.dat};
\addplot[only marks, OilBlue, mark=square*, mark size=2.0pt] table[x=x, y=debt]{data/clean/profiles/finh_300021.dat};
\legend{ROE, 资产负债率}
\end{groupplot}
\end{tikzpicture}
\begin{tikzpicture}
\begin{axis}[
  name=finbot,
  width=13.75cm, height=3.10cm, scale only axis,
  ymin=-2.91, ymax=32.62,
  xtick={0,1,2,3,4,5.7},
  xticklabels={2021,2022,2023,2024,2025,2026H1},
  tick label style={font=\tiny},
  yticklabel style={font=\tiny},
  ylabel={亿元}, ylabel style={font=\tiny},
  ybar, bar width=4pt,
  grid=major, grid style={LightGray, line width=0.3pt},
  axis line style={InkGray, line width=0.5pt},
  enlarge x limits={abs=0.8},
  legend style={font=\scriptsize, at={(0.5,-0.26)}, anchor=north,
                legend columns=2, draw=none, fill=none, column sep=10pt},
]
\addplot[fill=AgriGreen!55, draw=AgriGreen, bar shift=-3.0pt] table[x=i, y=cash]{data/clean/profiles/fin_300021.dat};
\addplot[fill=SoftRed!45, draw=SoftRed, bar shift=3.0pt] table[x=i, y=ibd]{data/clean/profiles/fin_300021.dat};
\addplot[fill=AgriGreen!25, draw=AgriGreen, pattern=north east lines, bar shift=-3.0pt] table[x=x, y=cash]{data/clean/profiles/finh_300021.dat};
\addplot[fill=SoftRed!20, draw=SoftRed, pattern=north east lines, bar shift=3.0pt] table[x=x, y=ibd]{data/clean/profiles/finh_300021.dat};
\legend{货币资金（左轴）, 有息负债（左轴）}
\end{axis}
\begin{axis}[
  at={(finbot.south east)}, anchor=south east,
  width=13.75cm, height=3.10cm, scale only axis,
  axis y line*=right, axis x line*=none, xtick=\empty,
  ymin=-8.72, ymax=9.46,
  tick label style={font=\tiny},
  yticklabel style={font=\tiny}, scaled y ticks=false,
  legend style={font=\scriptsize, at={(0.5,-0.44)}, anchor=north,
                legend columns=1, draw=none, fill=none},
]
\addplot[OilBlue, mark=*, mark size=2.2pt, line width=1.3pt] table[x=i, y=ocf]{data/clean/profiles/fin_300021.dat};
\addplot[only marks, OilBlue, mark=*, mark size=3.2pt] table[x=x, y=ocf]{data/clean/profiles/finh_300021.dat};
\legend{经营活动现金流净额（右轴，亿元，蓝点）}
\end{axis}
\end{tikzpicture}

\captionof{figure}[大禹节水（300021）近年主要财务指标]{大禹节水（300021）近年主要财务指标（左上：营业收入与归母净利润；右上：ROE 与资产负债率；下：货币资金、有息负债为左轴柱状，经营活动现金流净额为其右轴的蓝点折线。
2021---2025 年为年报口径，2026H1 为 2026 年半年报/半年末，与其前一年报之间留出间隔、以斜纹或空心点区分；半年报数据未年化）}
\end{minipage}
\end{center}

\pfl{财务分析} 2026 年 6 月末资产负债率 71.7\%，较 2025 年末 上升 1.1 个百分点（样本中位数 52.2\%，所处三级行业内按由低到高排第\zhnumber{2}位）。 2026 年 6 月末货币资金 12.85 亿元、短期债务（短期借款与一年内到期非流动负债）18.48 亿元、长期债务 10.64 亿元，有息负债合计 29.12 亿元，现金短债比 0.70 倍——覆盖偏紧。 净有息负债 16.27 亿元，相当于其 2026 年上半年营业收入的 103\%。 流动比率 1.25、速动比率 1.20，短期指标尚可。 2026 年上半年经营活动现金流净额 -5.69 亿元，经营现金流与利润同时为负，日常运营对债务与股东投入的依赖度较高。 2026 年 6 月末在建工程 0.61 亿元，较上年末增加 0.54 亿元，仍有在建投入。 2026 年 6 月末商誉 1.34 亿元，占归母净资产的 4.6\%，是净资产中最脆弱的一块。

\pfl{重大战略转型} 公司在报告期内的战略动作可归为 4 项：2026-03 转让所持慧图科技 28.5\% 股权给嘉兴乾瞻衡远创业投资合伙企业（有限合伙），交易价格 19,237.50 万元……；2026-06-30 依托“设计引领+全产业链”协同模式，设计业务成为第一增长引擎（农水设计服务收入同比 +306.84\%）……；2026-06-30 在生态治理与再生水新赛道取得突破，取得新疆阿克苏地区塔里木河干流盖孜库木片区生态综合治理项目……；2026-06-30 全面落地全域一体化营销机制（“大营销”协同机制），打通各业务板块资源壁垒，统筹全国区域项目联合开发，持续储备大中型灌区、流域治理……。 公告口径上，近三年（2023 年 9 月---2026 年 9 月）公司披露重大重组与并购类公告 9 条、对外投资与转型类公告 35 条、回购与股权激励类公告 42 条，合计公告 633 条。 主营结构的口径变化已经落到报表上：2021 年年报的第一大行业为“节水灌溉行业”（81.0\%），2025 年年报为“智慧水利行业”（100.0\%）；公司在这两期的分部划分方式有过调整。 综合判断，公司在报告期内有明确的外延扩张或业务结构调整动作，转型方向与融资安排之间存在直接关联。

\begin{center}\small\setlength{\tabcolsep}{4pt}
\captionof{table}[大禹节水（300021）历史融资与资本运作]{大禹节水（300021）历史融资与资本运作（据公司公告与公开报道整理；金额为披露值，含首次公开发行、再融资、债券、银行借款与重整投资等）}
\begin{adjustbox}{max width=\textwidth}
\begin{tabular}{@{}p{1.45cm}p{2.55cm}p{4.05cm}p{4.95cm}@{}}
\toprule
时间 & 方式 & 规模 & 用途与说明 \\
\midrule
2009-10 & IPO（深圳证券交易所创业板） & 发行 1,\allowbreak{}800 万股，\allowbreak{}发行价 14 元/\allowbreak{}股，\allowbreak{}预计募集资金总额 2.52 亿元、\allowbreak{}净额 2.2…… & 在酒泉总部及全资子公司新疆公司、武威公司建设滴灌管（带）及配套节水器材生产线技改扩建项目…… \\
2020-07-28 & 公开发行可转换公司债券“大禹转债” & 638 万张、\allowbreak{}每张面值 100 元，\allowbreak{}发行总额 6.38 亿元，\allowbreak{}期限 6 年（2020-07-…… & 高端节水灌溉产品智能工厂建设项目（拟配备以色列先进压力补偿式滴灌生产线及核心智能制造设备…… \\
2022-02-24 & 以简易程序向特定对象发行股票 & 58,\allowbreak{}593,\allowbreak{}750 股，\allowbreak{}发行价 5 元/\allowbreak{}股，\allowbreak{}募集资金总额 3.00 亿元，\allowbreak{}扣除发行费用 9…… & 未查得（募投项目及使用情况见年度报告募集资金专项说明）…… \\
2025-08-21 & 可转债提前赎回 & 赎回价格 100 元/\allowbreak{}张（含税）；公司披露本年度实际赎回 10,\allowbreak{}724.00 张大禹转债 & 终止“大禹转债”存续并在深交所摘牌；本年度因可转债转股增加股本 17,234.74 万股…… \\
2026-07 & 2025 年度利润分配 & 以 1,\allowbreak{}022,\allowbreak{}367,\allowbreak{}469 股为基数，\allowbreak{}每 10 股派发现金红利 0 元（含税），\allowbreak{}合计 2…… & 股东回报；2026 年 7 月完成派发 \\
\bottomrule
\end{tabular}
\end{adjustbox}
\end{center}

\pfl{历史融投资情况} 从现金流量表看，2025 年公司取得借款收到的现金 21.76 亿元、偿还债务支付的现金 16.85 亿元、吸收权益性投资 0.08 亿元，筹资活动现金流净额 2.39 亿元（2023 年为取得借款 14.48 亿元、偿还债务 9.61 亿元）。 2026 年上半年筹资活动现金流净额为 2.82 亿元（取得借款 10.47 亿元、偿还债务 7.10 亿元，未年化），仍处于净融入状态。 2025 年分配股利、利润或偿付利息支付的现金合计 1.46 亿元。 公司披露的股本与债务融资脉络包括：2025 年度利润分配、IPO（深圳证券交易所创业板）、以简易程序向特定对象发行股票、公开发行可转换公司债券“大禹转债”、可转债提前赎回。 其中规模较大的两笔为：2020-07-28公开发行可转换公司债券“大禹转债”（638 万张、每张面值 100 元，发行总额 6.38 亿元……）；2009-10IPO（深圳证券交易所创业板）（发行 1,800 万股，发行价 14.00 元/股，预计募集资金总额 25,200 万元……）。 股本方面，近三年披露股本变动 17 次（可转债转股 11 次、股份回购 4 次、其他 1 次）；总股本由 2021 年末的 8.01 亿股变为 8.54 亿股（+6.59\%）。 分红方面，近三年现金分红 4 次、累计每股派现 0.21 元（最近一次 2025-12-31）——分红规模与利润的下滑幅度相比仍然有限。 近三年“再融资”类公告 18 条，涉及定向增发、可转债。

\pfl{未来潜在融资需求分析} 2026 年 6 月末货币资金 12.85 亿元，而短期债务 18.48 亿元（现金短债比 0.70 倍），2026 年上半年经营现金流净额 -5.69 亿元。按“短期债务减去账面货币资金”的滚动偿付口径，静态缺口约 5.62 亿元；若再计入上半年亏损的年化额（0.02 亿元），维持一年正常经营所需的外部资金约 5.62---5.65 亿元。 按第 \ref{sec:financing-need} 节的三档划分，公司属于第一档“补血型”：2026 年上半年归母净利润 -0.01 亿元、6 月末资产负债率 71.7\%。 可行的融资路径按现实性排序为：定向增发（需股价配合，且控股股东需放弃优先认购或同步增资）、出售非核心资产与参股股权、债务重组或展期（与银行和债券持有人协商）。 公开披露的下一步融资线索包括：2025-06-30 可转债募集资金：本报告期投入 8,418.05 万元，截至 2025-06-30 累计使用 33,062.88 万元……；2025-08-26 公司提出高标准农田建设由整改前单体项目拆分招标转为整县整装推进，将显著提升单个订单体量，对具有规模优势的上市公司参与此类项目更有利……；2025-12-31 简易程序发行募集资金：本报告期投入 4,375.47 万元，截至 2025-12-31 累计使用 20,635.73 万元……。 资金需求还要叠加在建工程：期末在建工程 0.61 亿元、2026 年上半年购建固定资产等支出 0.45 亿元（未年化），这部分支出在周期底部通常会被主动放缓。





